% Recovered from the previously completed PDF; no new mathematical content.
% Prose is editable; formula boxes retain original symbol positions.
% Build from repository root: xelatex -no-shell-escape -halt-on-error -output-directory=... print/fudan/NAME.tex
\documentclass[11pt,a4paper]{article}
\usepackage[margin=23mm,headheight=14pt,headsep=8.5mm,footskip=10mm]{geometry}
\usepackage{amsmath,amssymb}
\DeclareMathAlphabet{\mathsf}{OT1}{cmss}{m}{n}
\usepackage[no-math]{fontspec}
\usepackage{xeCJK}
\setmainfont{Latin Modern Roman}
\setCJKmainfont{FandolSong-Regular.otf}[Path=assets/fonts/fandol/,BoldFont=FandolSong-Bold.otf,ItalicFont=FandolSong-Regular.otf,Script=CustomDefault]
\xeCJKsetup{PunctStyle=quanjiao,CheckSingle=true}
\usepackage{setspace}\setstretch{1.12}
\setlength{\parindent}{2em}\setlength{\parskip}{3pt}
\emergencystretch=6em
\hfuzz=9pt
\usepackage{titlesec}
\titleformat{\section}{\centering\normalfont\Large\bfseries}{}{0pt}{}
\titlespacing*{\section}{0pt}{1.8\baselineskip}{1.0\baselineskip}
\usepackage{fancyhdr}\pagestyle{fancy}\fancyhf{}
\fancyhead[L]{\small 数学手记·高等代数}\fancyhead[R]{\small 22S16–22S18、22A01–22A04}
\fancyfoot[C]{\thepage}\renewcommand{\headrulewidth}{0.4pt}
\usepackage[hidelinks,unicode]{hyperref}
\hypersetup{pdftitle={22S16–22S18、22A01–22A04：题目、解答与小结},pdfauthor={Vincent Arqain}}
\newcommand{\RB}[4]{\raisebox{0pt}[#2bp][#3bp]{\hbox to #1bp{#4\hss}}}
\newcommand{\RG}[5]{\hbox to 0pt{\kern#1bp\raisebox{#2bp}[0pt][0pt]{\fontsize{#3bp}{#3bp}\csname font#4\endcsname\char#5}\hss}}
\newcommand{\RT}[4]{\hbox to 0pt{\kern#1bp\raisebox{#2bp}[0pt][0pt]{\fontsize{#3bp}{#3bp}\selectfont#4}\hss}}
\newcommand{\RS}[4]{\hbox to 0pt{\kern#1bp\raisebox{#2bp}[0pt][0pt]{\fontsize{#3bp}{#3bp}\selectfont\ensuremath{#4}}\hss}}
\newcommand{\RR}[4]{\hbox to 0pt{\kern#1bp\raisebox{#2bp}[0pt][0pt]{\rule{#3bp}{#4bp}}\hss}}
\newcommand{\fontmi}{\usefont{OML}{cmm}{m}{it}}
\newcommand{\fontsy}{\usefont{OMS}{cmsy}{m}{n}}
\newcommand{\fontex}{\usefont{OMX}{cmex}{m}{n}}
\newcommand{\fontrm}{\usefont{OT1}{cmr}{m}{n}}
\newcommand{\fontbf}{\usefont{OT1}{cmr}{bx}{n}}
\newcommand{\fontam}{\usefont{U}{msa}{m}{n}}
\newcommand{\fontbb}{\usefont{U}{msb}{m}{n}}
\widowpenalty=10000\clubpenalty=10000\raggedbottom
\begin{document}
\begin{center}{\LARGE\bfseries 22S16–22S18、22A01–22A04\par}\vspace{8pt}复旦大学高等代数习题集·题目与解答\end{center}
% Original page 1.
\section*{2022S16}
将问题2021A06中的一一对应推广到\allowbreak{}\hskip.1em \RB{6.729}{5.322}{0.620}{%
\RG{0.000}{0.000}{10.9091}{mi}{110}%
}\hskip.1em\allowbreak{} 阶酉阵与\allowbreak{}\hskip.1em \RB{6.729}{5.322}{0.620}{%
\RG{0.000}{0.000}{10.9091}{mi}{110}%
}\hskip.1em\allowbreak{} 阶斜Hermite阵，并用二者的酉相似标准型说明这一对应的本质。解：以\allowbreak{}\hskip.1em \RB{15.939}{8.166}{0.500}{%
\RG{0.000}{0.001}{10.9091}{mi}{77}\RG{11.773}{3.960}{7.9701}{sy}{3}%
}\hskip.1em\allowbreak{} 表示共轭转置。令

\[
\RB{397.655}{9.290}{3.230}{%
\RG{102.524}{4.501}{7.9701}{sy}{3}\RG{276.296}{4.501}{7.9701}{sy}{3}\RG{0.000}{-0.003}{10.9091}{sy}{83}\RG{10.455}{-0.003}{10.9091}{rm}{61}\RG{21.970}{-0.003}{10.9091}{sy}{102}\RG{27.424}{-0.003}{10.9091}{mi}{75}\RG{40.500}{-0.003}{10.9091}{sy}{50}\RG{50.803}{-0.003}{10.9091}{mi}{77}\RG{67.024}{-0.003}{10.9091}{rm}{40}\RG{71.266}{-0.003}{10.9091}{bb}{67}\RG{79.145}{-0.003}{10.9091}{rm}{41}\RG{86.399}{-0.003}{10.9091}{rm}{58}\RG{92.478}{-0.003}{10.9091}{mi}{75}\RG{110.286}{-0.003}{10.9091}{rm}{61}\RG{121.801}{-0.003}{10.9091}{sy}{0}\RG{130.286}{-0.003}{10.9091}{mi}{75}\RG{140.332}{-0.003}{10.9091}{sy}{103}\RG{145.786}{-0.003}{10.9091}{mi}{59}\RG{172.453}{-0.003}{10.9091}{sy}{85}\RG{187.044}{-0.003}{10.9091}{rm}{61}\RG{198.559}{-0.003}{10.9091}{sy}{102}\RG{204.013}{-0.003}{10.9091}{mi}{85}\RG{215.681}{-0.003}{10.9091}{sy}{50}\RG{225.984}{-0.003}{10.9091}{mi}{77}\RG{242.204}{-0.003}{10.9091}{rm}{40}\RG{246.447}{-0.003}{10.9091}{bb}{67}\RG{254.325}{-0.003}{10.9091}{rm}{41}\RG{261.580}{-0.003}{10.9091}{rm}{58}\RG{267.659}{-0.003}{10.9091}{mi}{85}\RG{281.029}{-0.003}{10.9091}{mi}{85}\RG{292.697}{-0.003}{10.9091}{rm}{61}\RG{304.213}{-0.003}{10.9091}{mi}{73}\RG{309.853}{-0.003}{10.9091}{mi}{59}\RT{318.346}{-0.003}{10.9091}{d}\RT{324.411}{-0.003}{10.9091}{e}\RT{329.255}{-0.003}{10.9091}{t}\RG{333.498}{-0.003}{10.9091}{rm}{40}\RG{337.741}{-0.003}{10.9091}{mi}{73}\RG{345.816}{-0.003}{10.9091}{rm}{43}\RG{356.726}{-0.003}{10.9091}{mi}{85}\RG{365.363}{-0.003}{10.9091}{rm}{41}\RG{372.636}{-0.003}{10.9091}{sy}{54}\RG{372.636}{-0.003}{10.9091}{rm}{61}\RG{384.134}{-0.003}{10.9091}{rm}{48}\RG{389.606}{-0.003}{10.9091}{sy}{103}\RG{395.060}{-0.003}{10.9091}{mi}{58}\RG{61.387}{-1.640}{7.9701}{mi}{110}\RG{179.280}{-1.640}{7.9701}{rm}{48}\RG{236.568}{-1.640}{7.9701}{mi}{110}%
}
\]
所求的Cayley变换及其逆变换为

\[
\RB{280.254}{10.303}{3.228}{%
\RG{115.986}{4.503}{7.9701}{sy}{0}\RG{122.572}{4.503}{7.9701}{rm}{49}\RG{266.340}{4.503}{7.9701}{sy}{0}\RG{272.927}{4.503}{7.9701}{rm}{49}\RG{0.000}{-0.001}{10.9091}{rm}{8}\RG{7.876}{-0.001}{10.9091}{rm}{40}\RG{12.121}{-0.001}{10.9091}{mi}{75}\RG{22.167}{-0.001}{10.9091}{rm}{41}\RG{29.422}{-0.001}{10.9091}{rm}{61}\RG{40.931}{-0.001}{10.9091}{rm}{40}\RG{45.197}{-0.001}{10.9091}{mi}{73}\RG{53.273}{-0.001}{10.9091}{sy}{0}\RG{64.183}{-0.001}{10.9091}{mi}{75}\RG{74.228}{-0.001}{10.9091}{rm}{41}\RG{78.460}{-0.001}{10.9091}{rm}{40}\RG{82.713}{-0.001}{10.9091}{mi}{73}\RG{90.789}{-0.001}{10.9091}{rm}{43}\RG{101.698}{-0.001}{10.9091}{mi}{75}\RG{111.743}{-0.001}{10.9091}{rm}{41}\RG{127.305}{-0.001}{10.9091}{mi}{59}\RG{153.971}{-0.001}{10.9091}{rm}{9}\RG{162.448}{-0.001}{10.9091}{rm}{40}\RG{166.699}{-0.001}{10.9091}{mi}{85}\RG{175.337}{-0.001}{10.9091}{rm}{41}\RG{182.592}{-0.001}{10.9091}{rm}{61}\RG{194.101}{-0.001}{10.9091}{rm}{40}\RG{198.368}{-0.001}{10.9091}{mi}{73}\RG{206.443}{-0.001}{10.9091}{sy}{0}\RG{217.353}{-0.001}{10.9091}{mi}{85}\RG{225.990}{-0.001}{10.9091}{rm}{41}\RG{230.223}{-0.001}{10.9091}{rm}{40}\RG{234.475}{-0.001}{10.9091}{mi}{73}\RG{242.551}{-0.001}{10.9091}{rm}{43}\RG{253.460}{-0.001}{10.9091}{mi}{85}\RG{262.098}{-0.001}{10.9091}{rm}{41}\RG{277.659}{-0.001}{10.9091}{mi}{58}%
}
\]
这里必须从酉阵中排除具有特征值\allowbreak{}\hskip.1em \RB{13.556}{7.771}{0.500}{%
\RG{0.000}{0.005}{10.9091}{sy}{0}\RG{8.485}{0.005}{10.9091}{rm}{49}%
}\hskip.1em\allowbreak{} 的矩阵。

先验证\allowbreak{}\hskip.1em \RB{7.755}{7.952}{0.500}{%
\RG{0.000}{0.001}{10.9091}{rm}{8}%
}\hskip.1em\allowbreak{} 的良定义性。若\allowbreak{}\hskip.1em \RB{64.430}{8.683}{3.226}{%
\RG{0.000}{0.001}{10.9091}{rm}{40}\RG{4.242}{0.001}{10.9091}{mi}{73}\RG{12.488}{0.001}{10.9091}{rm}{43}\RG{23.567}{0.001}{10.9091}{mi}{75}\RG{33.612}{0.001}{10.9091}{rm}{41}\RG{37.854}{0.001}{10.9091}{mi}{118}\RG{46.988}{0.001}{10.9091}{rm}{61}\RG{58.912}{0.001}{10.9091}{rm}{48}%
}\hskip.1em\allowbreak{} ，则\allowbreak{}\hskip.1em \RB{70.906}{9.760}{3.226}{%
\RG{0.000}{0.001}{10.9091}{mi}{118}\RG{5.680}{3.960}{7.9701}{sy}{3}\RG{10.412}{0.001}{10.9091}{mi}{75}\RG{20.449}{0.001}{10.9091}{mi}{118}\RG{29.591}{0.001}{10.9091}{rm}{61}\RG{41.531}{0.001}{10.9091}{sy}{0}\RG{50.008}{0.001}{10.9091}{sy}{107}\RG{55.470}{0.001}{10.9091}{mi}{118}\RG{61.150}{0.001}{10.9091}{sy}{107}\RG{66.605}{3.960}{7.9701}{rm}{50}%
}\hskip.1em\allowbreak{} ；而\allowbreak{}\hskip.1em \RB{122.332}{8.166}{0.619}{%
\RG{0.000}{0.001}{10.9091}{mi}{118}\RG{5.679}{3.153}{7.9701}{sy}{3}\RG{10.411}{0.001}{10.9091}{mi}{75}\RG{20.447}{0.001}{10.9091}{mi}{118}\RG{29.590}{0.001}{10.9091}{rm}{61}\RG{41.530}{0.001}{10.9091}{mi}{118}\RG{47.209}{3.960}{7.9701}{sy}{3}\RG{51.942}{0.001}{10.9091}{mi}{75}\RG{61.987}{3.960}{7.9701}{sy}{3}\RG{66.720}{0.001}{10.9091}{mi}{118}\RG{75.854}{0.001}{10.9091}{rm}{61}\RG{87.794}{0.001}{10.9091}{sy}{0}\RG{96.278}{0.001}{10.9091}{mi}{118}\RG{101.958}{3.960}{7.9701}{sy}{3}\RG{106.690}{0.001}{10.9091}{mi}{75}\RG{116.726}{0.001}{10.9091}{mi}{118}\RR{0.000}{8.983}{26.136}{0.436}%
}\hskip.1em\allowbreak{} ，故\allowbreak{}\hskip.1em \RB{26.054}{8.161}{0.624}{%
\RG{0.000}{-0.004}{10.9091}{mi}{118}\RG{5.679}{3.955}{7.9701}{sy}{3}\RG{10.412}{-0.004}{10.9091}{mi}{75}\RG{20.449}{-0.004}{10.9091}{mi}{118}%
}\hskip.1em\allowbreak{} 为纯虚数，迫使\allowbreak{}\hskip.1em \RB{25.726}{7.762}{0.744}{%
\RG{0.000}{-0.004}{10.9091}{mi}{118}\RG{8.709}{-0.004}{10.9091}{rm}{61}\RG{20.207}{-0.004}{10.9091}{rm}{48}%
}\hskip.1em\allowbreak{} 。因此\allowbreak{}\hskip.1em \RB{28.519}{7.947}{1.409}{%
\RG{0.000}{-0.004}{10.9091}{mi}{73}\RG{7.743}{-0.004}{10.9091}{rm}{43}\RG{18.321}{-0.004}{10.9091}{mi}{75}%
}\hskip.1em\allowbreak{} 可逆，同理\allowbreak{}\hskip.1em \RB{28.519}{7.947}{0.504}{%
\RG{0.000}{-0.004}{10.9091}{mi}{73}\RG{7.744}{-0.004}{10.9091}{sy}{0}\RG{18.321}{-0.004}{10.9091}{mi}{75}%
}\hskip.1em\allowbreak{} 可逆。令\allowbreak{}\hskip.1em \RB{49.003}{8.678}{3.231}{%
\RG{0.000}{-0.004}{10.9091}{mi}{85}\RG{11.668}{-0.004}{10.9091}{rm}{61}\RG{23.166}{-0.004}{10.9091}{rm}{8}\RG{31.043}{-0.004}{10.9091}{rm}{40}\RG{35.305}{-0.004}{10.9091}{mi}{75}\RG{45.350}{-0.004}{10.9091}{rm}{41}%
}\hskip.1em\allowbreak{} ，利用\allowbreak{}\hskip.1em \RB{28.518}{7.947}{0.504}{%
\RG{0.000}{-0.004}{10.9091}{mi}{73}\RG{7.743}{-0.004}{10.9091}{sy}{6}\RG{18.320}{-0.004}{10.9091}{mi}{75}%
}\hskip.1em\allowbreak{}及其逆彼此交换，得

\[
\RB{308.856}{10.301}{3.230}{%
\RG{8.638}{4.502}{7.9701}{sy}{3}\RG{65.432}{4.502}{7.9701}{sy}{0}\RG{72.018}{4.501}{7.9701}{rm}{49}\RG{149.570}{4.502}{7.9701}{sy}{3}\RG{294.943}{4.502}{7.9701}{sy}{0}\RG{301.529}{4.501}{7.9701}{rm}{49}\RG{0.000}{-0.003}{10.9091}{mi}{85}\RG{16.401}{-0.003}{10.9091}{rm}{61}\RG{27.899}{-0.003}{10.9091}{rm}{40}\RG{32.158}{-0.003}{10.9091}{mi}{73}\RG{40.234}{-0.003}{10.9091}{sy}{0}\RG{51.143}{-0.003}{10.9091}{mi}{75}\RG{61.189}{-0.003}{10.9091}{rm}{41}\RG{76.751}{-0.003}{10.9091}{rm}{40}\RG{80.993}{-0.003}{10.9091}{mi}{73}\RG{89.069}{-0.003}{10.9091}{rm}{43}\RG{99.978}{-0.003}{10.9091}{mi}{75}\RG{110.023}{-0.003}{10.9091}{rm}{41}\RG{114.266}{-0.003}{10.9091}{mi}{59}\RG{140.917}{-0.003}{10.9091}{mi}{85}\RG{154.302}{-0.003}{10.9091}{mi}{85}\RG{165.971}{-0.003}{10.9091}{rm}{61}\RG{177.486}{-0.003}{10.9091}{mi}{73}\RG{183.127}{-0.003}{10.9091}{mi}{59}\RG{209.789}{-0.003}{10.9091}{mi}{73}\RG{217.881}{-0.003}{10.9091}{rm}{43}\RG{228.790}{-0.003}{10.9091}{mi}{85}\RG{240.458}{-0.003}{10.9091}{rm}{61}\RG{251.956}{-0.003}{10.9091}{rm}{50}\RG{257.411}{-0.003}{10.9091}{rm}{40}\RG{261.670}{-0.003}{10.9091}{mi}{73}\RG{269.746}{-0.003}{10.9091}{rm}{43}\RG{280.655}{-0.003}{10.9091}{mi}{75}\RG{290.700}{-0.003}{10.9091}{rm}{41}\RG{306.262}{-0.003}{10.9091}{mi}{58}%
}
\]
故\allowbreak{}\hskip.1em \RB{33.196}{7.954}{2.300}{%
\RG{0.000}{0.003}{10.9091}{mi}{85}\RG{11.668}{0.003}{10.9091}{sy}{50}\RG{21.955}{0.003}{10.9091}{sy}{85}\RG{28.798}{-1.633}{7.9701}{rm}{48}%
}\hskip.1em\allowbreak{} 。

反之，若\allowbreak{}\hskip.1em \RB{33.197}{7.950}{2.304}{%
\RG{0.000}{-0.001}{10.9091}{mi}{85}\RG{11.668}{-0.001}{10.9091}{sy}{50}\RG{21.955}{-0.001}{10.9091}{sy}{85}\RG{28.799}{-1.637}{7.9701}{rm}{48}%
}\hskip.1em\allowbreak{} ，令\allowbreak{}\hskip.1em \RB{49.609}{8.681}{3.228}{%
\RG{0.000}{-0.001}{10.9091}{mi}{75}\RG{13.076}{-0.001}{10.9091}{rm}{61}\RG{24.574}{-0.001}{10.9091}{rm}{9}\RG{33.051}{-0.001}{10.9091}{rm}{40}\RG{37.318}{-0.001}{10.9091}{mi}{85}\RG{45.956}{-0.001}{10.9091}{rm}{41}%
}\hskip.1em\allowbreak{} 。由\allowbreak{}\hskip.1em \RB{47.203}{9.758}{0.741}{%
\RG{0.000}{-0.001}{10.9091}{mi}{85}\RG{8.638}{3.958}{7.9701}{sy}{3}\RG{16.401}{-0.001}{10.9091}{rm}{61}\RG{27.916}{-0.001}{10.9091}{mi}{85}\RG{36.554}{3.958}{7.9701}{sy}{0}\RG{43.140}{3.958}{7.9701}{rm}{49}%
}\hskip.1em\allowbreak{} ，

\[
\RB{255.454}{10.300}{3.231}{%
\RG{10.046}{4.500}{7.9701}{sy}{3}\RG{61.190}{4.500}{7.9701}{sy}{3}\RG{70.164}{4.500}{7.9701}{sy}{0}\RG{76.751}{4.500}{7.9701}{rm}{49}\RG{113.348}{4.500}{7.9701}{sy}{3}\RG{172.976}{4.500}{7.9701}{sy}{0}\RG{179.562}{4.500}{7.9701}{rm}{49}\RG{0.000}{-0.004}{10.9091}{mi}{75}\RG{17.808}{-0.004}{10.9091}{rm}{61}\RG{29.306}{-0.004}{10.9091}{rm}{40}\RG{33.566}{-0.004}{10.9091}{mi}{73}\RG{41.642}{-0.004}{10.9091}{rm}{43}\RG{52.551}{-0.004}{10.9091}{mi}{85}\RG{65.922}{-0.004}{10.9091}{rm}{41}\RG{81.483}{-0.004}{10.9091}{rm}{40}\RG{85.726}{-0.004}{10.9091}{mi}{73}\RG{93.801}{-0.004}{10.9091}{sy}{0}\RG{104.710}{-0.004}{10.9091}{mi}{85}\RG{118.081}{-0.004}{10.9091}{rm}{41}\RG{125.336}{-0.004}{10.9091}{rm}{61}\RG{136.845}{-0.004}{10.9091}{rm}{40}\RG{141.111}{-0.004}{10.9091}{mi}{73}\RG{149.187}{-0.004}{10.9091}{rm}{43}\RG{160.096}{-0.004}{10.9091}{mi}{85}\RG{168.734}{-0.004}{10.9091}{rm}{41}\RG{184.296}{-0.004}{10.9091}{rm}{40}\RG{188.538}{-0.004}{10.9091}{mi}{85}\RG{199.600}{-0.004}{10.9091}{sy}{0}\RG{210.509}{-0.004}{10.9091}{mi}{73}\RG{216.161}{-0.004}{10.9091}{rm}{41}\RG{223.415}{-0.004}{10.9091}{rm}{61}\RG{234.949}{-0.004}{10.9091}{sy}{0}\RG{243.434}{-0.004}{10.9091}{mi}{75}\RG{252.860}{-0.004}{10.9091}{mi}{58}%
}
\]
因此\allowbreak{}\hskip.1em \RB{30.867}{8.193}{0.934}{%
\RG{0.000}{0.002}{10.9091}{mi}{75}\RG{13.076}{0.002}{10.9091}{sy}{50}\RG{23.363}{0.002}{10.9091}{sy}{83}%
}\hskip.1em\allowbreak{} 。又因

\[
\RB{267.885}{10.302}{3.229}{%
\RG{112.955}{4.502}{7.9701}{sy}{0}\RG{119.542}{4.502}{7.9701}{rm}{49}\RG{253.851}{4.502}{7.9701}{sy}{0}\RG{260.437}{4.502}{7.9701}{rm}{49}\RG{0.000}{-0.002}{10.9091}{mi}{73}\RG{8.076}{-0.002}{10.9091}{sy}{0}\RG{18.985}{-0.002}{10.9091}{rm}{8}\RG{26.862}{-0.002}{10.9091}{rm}{40}\RG{31.106}{-0.002}{10.9091}{mi}{75}\RG{41.151}{-0.002}{10.9091}{rm}{41}\RG{48.406}{-0.002}{10.9091}{rm}{61}\RG{59.915}{-0.002}{10.9091}{rm}{50}\RG{65.395}{-0.002}{10.9091}{mi}{75}\RG{75.440}{-0.002}{10.9091}{rm}{40}\RG{79.683}{-0.002}{10.9091}{mi}{73}\RG{87.758}{-0.002}{10.9091}{rm}{43}\RG{98.667}{-0.002}{10.9091}{mi}{75}\RG{108.713}{-0.002}{10.9091}{rm}{41}\RG{124.274}{-0.002}{10.9091}{mi}{59}\RG{150.925}{-0.002}{10.9091}{mi}{73}\RG{159.017}{-0.002}{10.9091}{rm}{43}\RG{169.916}{-0.002}{10.9091}{rm}{8}\RG{177.792}{-0.002}{10.9091}{rm}{40}\RG{182.048}{-0.002}{10.9091}{mi}{75}\RG{192.093}{-0.002}{10.9091}{rm}{41}\RG{199.347}{-0.002}{10.9091}{rm}{61}\RG{210.856}{-0.002}{10.9091}{rm}{50}\RG{216.311}{-0.002}{10.9091}{rm}{40}\RG{220.578}{-0.002}{10.9091}{mi}{73}\RG{228.654}{-0.002}{10.9091}{rm}{43}\RG{239.563}{-0.002}{10.9091}{mi}{75}\RG{249.608}{-0.002}{10.9091}{rm}{41}\RG{265.170}{-0.002}{10.9091}{mi}{59}%
}
\]
所以\allowbreak{}\hskip.1em \RB{68.123}{8.686}{3.223}{%
\RG{0.000}{0.004}{10.9091}{rm}{9}\RG{8.476}{0.004}{10.9091}{rm}{40}\RG{12.709}{0.004}{10.9091}{rm}{8}\RG{20.586}{0.004}{10.9091}{rm}{40}\RG{24.848}{0.004}{10.9091}{mi}{75}\RG{34.894}{0.004}{10.9091}{rm}{41}\RG{39.127}{0.004}{10.9091}{rm}{41}\RG{46.382}{0.004}{10.9091}{rm}{61}\RG{57.924}{0.004}{10.9091}{mi}{75}%
}\hskip.1em\allowbreak{} ；将\allowbreak{}\hskip.1em \RB{10.198}{7.955}{0.500}{%
\RG{0.000}{0.004}{10.9091}{mi}{75}%
}\hskip.1em\allowbreak{} 换成\allowbreak{}\hskip.1em \RB{8.791}{7.955}{0.736}{%
\RG{0.000}{0.004}{10.9091}{mi}{85}%
}\hskip.1em\allowbreak{} ，同一计算也给出\allowbreak{}\hskip.1em \RB{65.308}{8.686}{3.223}{%
\RG{0.000}{0.004}{10.9091}{rm}{8}\RG{7.876}{0.004}{10.9091}{rm}{40}\RG{12.109}{0.004}{10.9091}{rm}{9}\RG{20.586}{0.004}{10.9091}{rm}{40}\RG{24.848}{0.004}{10.9091}{mi}{85}\RG{33.486}{0.004}{10.9091}{rm}{41}\RG{37.719}{0.004}{10.9091}{rm}{41}\RG{44.974}{0.004}{10.9091}{rm}{61}\RG{56.517}{0.004}{10.9091}{mi}{85}%
}\hskip.1em\allowbreak{} 。二者确为互逆的双射。

最后说明标准型意义。正规矩阵的谱定理给出：若\allowbreak{}\hskip.1em \RB{48.007}{8.165}{0.500}{%
\RG{0.000}{-0.000}{10.9091}{mi}{75}\RG{10.045}{3.959}{7.9701}{sy}{3}\RG{17.808}{-0.000}{10.9091}{rm}{61}\RG{29.323}{-0.000}{10.9091}{sy}{0}\RG{37.808}{-0.000}{10.9091}{mi}{75}%
}\hskip.1em\allowbreak{} ，则存在酉阵\allowbreak{}\hskip.1em \RB{8.573}{8.191}{2.616}{%
\RG{0.000}{-0.000}{10.9091}{mi}{81}%
}\hskip.1em\allowbreak{} 使

\[
\RB{185.685}{9.289}{3.766}{%
\RG{122.157}{4.500}{7.9701}{sy}{3}\RG{0.000}{-0.004}{10.9091}{mi}{75}\RG{13.076}{-0.004}{10.9091}{rm}{61}\RG{24.591}{-0.004}{10.9091}{mi}{81}\RT{35.034}{-0.004}{10.9091}{d}\RT{41.100}{-0.004}{10.9091}{i}\RT{44.133}{-0.004}{10.9091}{a}\RT{49.587}{-0.004}{10.9091}{g}\RG{55.041}{-0.004}{10.9091}{rm}{40}\RG{59.284}{-0.004}{10.9091}{mi}{105}\RG{63.036}{-0.004}{10.9091}{mi}{116}\RG{71.714}{-0.004}{10.9091}{mi}{59}\RG{76.546}{-0.004}{10.9091}{mi}{58}\RG{81.390}{-0.004}{10.9091}{mi}{58}\RG{86.234}{-0.004}{10.9091}{mi}{58}\RG{91.066}{-0.004}{10.9091}{mi}{59}\RG{95.910}{-0.004}{10.9091}{mi}{105}\RG{99.663}{-0.004}{10.9091}{mi}{116}\RG{109.291}{-0.004}{10.9091}{rm}{41}\RG{113.533}{-0.004}{10.9091}{mi}{81}\RG{126.890}{-0.004}{10.9091}{mi}{59}\RG{153.541}{-0.004}{10.9091}{mi}{116}\RG{164.908}{-0.004}{10.9091}{sy}{50}\RG{175.211}{-0.004}{10.9091}{bb}{82}\RG{183.090}{-0.004}{10.9091}{mi}{58}\RG{66.982}{-1.640}{7.9701}{rm}{49}\RG{103.654}{-1.640}{7.9701}{mi}{110}\RG{157.496}{-1.640}{7.9701}{mi}{106}%
}
\]
这里斜Hermite阵是正规阵，其特征值为纯虚数。代入可得

\[
\RB{197.599}{16.321}{10.850}{%
\RG{71.405}{15.385}{10.9092}{ex}{18}\RG{171.799}{15.385}{10.9092}{ex}{19}\RG{80.631}{7.383}{10.9091}{rm}{49}\RG{88.510}{7.383}{10.9091}{sy}{0}\RG{99.419}{7.383}{10.9091}{mi}{105}\RG{103.172}{7.383}{10.9091}{mi}{116}\RG{138.482}{7.383}{10.9091}{rm}{49}\RG{146.361}{7.383}{10.9091}{sy}{0}\RG{157.270}{7.383}{10.9091}{mi}{105}\RG{161.023}{7.383}{10.9091}{mi}{116}\RG{107.116}{5.747}{7.9701}{rm}{49}\RG{164.968}{5.747}{7.9701}{mi}{110}\RG{190.273}{4.507}{7.9701}{sy}{3}\RG{0.000}{0.003}{10.9091}{rm}{8}\RG{7.876}{0.003}{10.9091}{rm}{40}\RG{12.121}{0.003}{10.9091}{mi}{75}\RG{22.167}{0.003}{10.9091}{rm}{41}\RG{29.422}{0.003}{10.9091}{rm}{61}\RG{40.954}{0.003}{10.9091}{mi}{81}\RT{51.397}{0.003}{10.9091}{d}\RT{57.463}{0.003}{10.9091}{i}\RT{60.496}{0.003}{10.9091}{a}\RT{65.950}{0.003}{10.9091}{g}\RG{113.044}{0.003}{10.9091}{mi}{59}\RG{117.877}{0.003}{10.9091}{mi}{58}\RG{122.721}{0.003}{10.9091}{mi}{58}\RG{127.564}{0.003}{10.9091}{mi}{58}\RG{132.397}{0.003}{10.9091}{mi}{59}\RG{181.649}{0.003}{10.9091}{mi}{81}\RG{195.005}{0.003}{10.9091}{mi}{58}\RG{80.631}{-7.480}{10.9091}{rm}{49}\RG{88.507}{-7.480}{10.9091}{rm}{43}\RG{99.419}{-7.480}{10.9091}{mi}{105}\RG{103.172}{-7.480}{10.9091}{mi}{116}\RG{138.482}{-7.480}{10.9091}{rm}{49}\RG{146.359}{-7.480}{10.9091}{rm}{43}\RG{157.270}{-7.480}{10.9091}{mi}{105}\RG{161.023}{-7.480}{10.9091}{mi}{116}\RG{107.116}{-9.116}{7.9701}{rm}{49}\RG{164.968}{-9.116}{7.9701}{mi}{110}\RR{80.631}{2.730}{31.218}{0.436}\RR{138.482}{2.730}{32.122}{0.436}%
}
\]
若\allowbreak{}\hskip.1em \RB{56.266}{8.187}{0.940}{%
\RG{0.000}{-0.004}{10.9091}{sy}{0}\RG{8.485}{-0.004}{10.9091}{mi}{25}\RG{18.118}{-0.004}{10.9091}{mi}{60}\RG{29.627}{-0.004}{10.9091}{mi}{18}\RG{38.071}{-0.004}{10.9091}{mi}{60}\RG{49.580}{-0.004}{10.9091}{mi}{25}%
}\hskip.1em\allowbreak{} ，则

\[
\RB{159.563}{15.570}{8.891}{%
\RG{0.000}{7.379}{10.9091}{rm}{49}\RG{7.879}{7.379}{10.9091}{sy}{0}\RG{18.788}{7.379}{10.9091}{mi}{105}\RG{22.541}{7.379}{10.9091}{mi}{116}\RG{150.333}{7.379}{10.9091}{mi}{18}\RG{47.306}{4.503}{7.9701}{mi}{105}\RG{50.184}{4.503}{7.9701}{mi}{18}\RG{30.711}{-0.001}{10.9091}{rm}{61}\RG{42.226}{-0.001}{10.9091}{mi}{101}\RG{68.813}{-0.001}{10.9091}{sy}{40}\RG{77.911}{-0.001}{10.9091}{sy}{41}\RG{102.753}{-0.001}{10.9091}{mi}{116}\RG{109.722}{-0.001}{10.9091}{rm}{61}\RG{121.237}{-0.001}{10.9091}{sy}{0}\RT{131.541}{-0.001}{10.9091}{t}\RT{135.785}{-0.001}{10.9091}{a}\RT{141.239}{-0.001}{10.9091}{n}\RG{156.968}{-0.001}{10.9091}{mi}{58}\RG{0.000}{-7.485}{10.9091}{rm}{49}\RG{7.876}{-7.485}{10.9091}{rm}{43}\RG{18.788}{-7.485}{10.9091}{mi}{105}\RG{22.541}{-7.485}{10.9091}{mi}{116}\RG{150.318}{-7.485}{10.9091}{rm}{50}\RR{0.000}{2.726}{26.485}{0.436}\RR{150.318}{2.726}{5.454}{0.436}%
}
\]
所以这一对应保留同一组标准正交特征向量，只将每个特征值\allowbreak{}\hskip.1em \RB{7.853}{7.710}{0.621}{%
\RG{0.000}{-0.001}{10.9091}{mi}{105}\RG{3.753}{-0.001}{10.9091}{mi}{116}%
}\hskip.1em\allowbreak{} 按上述分式变为单位圆上除\allowbreak{}\hskip.1em \RB{13.556}{7.765}{0.501}{%
\RG{0.000}{-0.001}{10.9091}{sy}{0}\RG{8.485}{-0.001}{10.9091}{rm}{49}%
}\hskip.1em\allowbreak{}外的点。谱定理在此也可由归纳得到：正规矩阵的特征向量同时是伴随矩阵的特征向量，其正交补仍不变，逐次取特征向量即可。\allowbreak{}\hskip.1em \RB{8.366}{7.994}{0.501}{%
\RG{0.000}{-0.001}{10.9091}{am}{3}%
}\hskip.1em\allowbreak{}

方法小结：矩阵层面的双射来自同一个分式变换；在酉相似标准型下，它就是纯虚轴与去掉\allowbreak{}\hskip.1em \RB{13.556}{7.763}{0.503}{%
\RG{0.000}{-0.003}{10.9091}{sy}{0}\RG{8.485}{-0.003}{10.9091}{rm}{49}%
}\hskip.1em\allowbreak{} 的单位圆之间的逐点对应。

题目来源：谢启鸿《复旦大学高等代数习题集（2023版）》，第117页。

% Original page 2.
\section*{2022S17}
设\allowbreak{}\hskip.1em \RB{21.762}{8.310}{2.606}{%
\RG{0.000}{-0.001}{10.9091}{mi}{65}\RG{8.182}{-0.001}{10.9091}{mi}{59}\RG{13.015}{-0.001}{10.9091}{mi}{66}%
}\hskip.1em\allowbreak{} 为\allowbreak{}\hskip.1em \RB{6.729}{5.321}{0.621}{%
\RG{0.000}{-0.001}{10.9091}{mi}{110}%
}\hskip.1em\allowbreak{} 阶实方阵，\allowbreak{}\hskip.1em \RB{8.366}{8.310}{0.501}{%
\RG{0.000}{-0.001}{10.9091}{mi}{65}%
}\hskip.1em\allowbreak{} 的\allowbreak{}\hskip.1em \RB{6.729}{5.321}{0.621}{%
\RG{0.000}{-0.001}{10.9091}{mi}{110}%
}\hskip.1em\allowbreak{} 个特征值均为正实数，且

\[
\RB{98.868}{11.423}{1.410}{%
\RT{47.341}{4.499}{7.9701}{T}\RT{89.990}{4.499}{7.9701}{T}\RG{0.000}{-0.005}{10.9091}{mi}{65}\RG{8.182}{-0.005}{10.9091}{mi}{66}\RG{19.428}{-0.005}{10.9091}{rm}{43}\RG{30.337}{-0.005}{10.9091}{mi}{66}\RG{39.152}{-0.005}{10.9091}{mi}{65}\RG{56.656}{-0.005}{10.9091}{rm}{61}\RG{68.155}{-0.005}{10.9091}{rm}{50}\RG{73.626}{-0.005}{10.9091}{mi}{65}\RG{81.807}{-0.005}{10.9091}{mi}{65}\RG{96.274}{-0.005}{10.9091}{mi}{58}%
}
\]
证明：

(1)\allowbreak{}\hskip.1em \RB{8.747}{7.949}{0.502}{%
\RG{0.000}{-0.002}{10.9091}{mi}{66}%
}\hskip.1em\allowbreak{} 为对称阵；(2)\allowbreak{}\hskip.1em \RB{8.366}{8.312}{0.500}{%
\RG{0.000}{0.001}{10.9091}{mi}{65}%
}\hskip.1em\allowbreak{} 为对称阵，当且仅当\allowbreak{}\hskip.1em \RB{31.475}{8.312}{0.500}{%
\RG{0.000}{0.001}{10.9091}{mi}{65}\RG{11.212}{0.001}{10.9091}{rm}{61}\RG{22.727}{0.001}{10.9091}{mi}{66}%
}\hskip.1em\allowbreak{} ，也当且仅当tr\allowbreak{}\hskip.1em \RB{33.416}{9.759}{3.226}{%
\RG{0.000}{0.001}{10.9091}{rm}{40}\RG{4.242}{0.001}{10.9091}{mi}{66}\RG{13.064}{3.959}{7.9701}{rm}{50}\RG{17.796}{0.001}{10.9091}{rm}{41}\RG{25.051}{0.001}{10.9091}{rm}{61}%
}\hskip.1em\allowbreak{} tr\allowbreak{}\hskip.1em \RB{20.790}{8.683}{3.226}{%
\RG{0.000}{0.001}{10.9091}{rm}{40}\RG{4.243}{0.001}{10.9091}{mi}{65}\RG{12.425}{0.001}{10.9091}{mi}{65}%
}\hskip.1em\allowbreak{} T\allowbreak{}\hskip.1em \RB{3.653}{8.683}{3.226}{%
\RG{0.000}{0.001}{10.9091}{rm}{41}%
}\hskip.1em\allowbreak{} ；(3) det\allowbreak{}\hskip.1em \RB{19.923}{7.954}{1.992}{%
\RG{0.000}{0.003}{10.9091}{mi}{66}\RG{11.852}{0.003}{10.9091}{sy}{21}%
}\hskip.1em\allowbreak{} det\allowbreak{}\hskip.1em \RB{8.366}{8.314}{0.500}{%
\RG{0.000}{0.003}{10.9091}{mi}{65}%
}\hskip.1em\allowbreak{} ，且等号成立当且仅当\allowbreak{}\hskip.1em \RB{31.475}{8.314}{0.500}{%
\RG{0.000}{0.003}{10.9091}{mi}{65}\RG{11.212}{0.003}{10.9091}{rm}{61}\RG{22.727}{0.003}{10.9091}{mi}{66}%
}\hskip.1em\allowbreak{} 。

证明：先建立一个唯一性结论：

\[
\RB{153.714}{11.424}{1.409}{%
\RT{49.485}{4.500}{7.9701}{T}\RG{0.000}{-0.004}{10.9091}{mi}{65}\RG{8.182}{-0.004}{10.9091}{mi}{88}\RG{20.500}{-0.004}{10.9091}{rm}{43}\RG{31.409}{-0.004}{10.9091}{mi}{88}\RG{41.293}{-0.004}{10.9091}{mi}{65}\RG{58.799}{-0.004}{10.9091}{rm}{61}\RG{70.298}{-0.004}{10.9091}{rm}{48}\RG{89.694}{-0.004}{10.9091}{rm}{61}\RG{96.376}{-0.004}{10.9091}{sy}{41}\RG{121.225}{-0.004}{10.9091}{mi}{88}\RG{134.149}{-0.004}{10.9091}{rm}{61}\RG{145.647}{-0.004}{10.9091}{rm}{48}\RG{151.119}{-0.004}{10.9091}{mi}{58}%
}
\]
\noindent \RB{465.385}{10.885}{18.400}{%
\RG{271.842}{8.184}{10.9092}{ex}{81}\RG{282.145}{5.029}{7.9701}{mi}{110}\RT{421.230}{3.961}{7.9701}{T}\RT{0.000}{0.002}{10.9091}{设}\RG{14.975}{0.002}{10.9091}{mi}{65}\RT{27.222}{0.002}{10.9091}{的}\RT{38.415}{0.002}{10.9091}{特}\RT{49.619}{0.002}{10.9091}{征}\RT{60.812}{0.002}{10.9091}{值}\RT{72.015}{0.002}{10.9091}{为}\RG{86.991}{0.002}{10.9091}{mi}{21}\RG{98.087}{0.002}{10.9091}{mi}{59}\RG{102.920}{0.002}{10.9091}{mi}{58}\RG{107.763}{0.002}{10.9091}{mi}{58}\RG{112.607}{0.002}{10.9091}{mi}{58}\RG{117.440}{0.002}{10.9091}{mi}{59}\RG{122.283}{0.002}{10.9091}{mi}{21}\RG{138.082}{0.002}{10.9091}{mi}{62}\RG{150.319}{0.002}{10.9091}{rm}{48}\RT{155.774}{0.002}{10.9091}{，}\RT{166.967}{0.002}{10.9091}{特}\RT{178.171}{0.002}{10.9091}{征}\RT{189.364}{0.002}{10.9091}{多}\RT{200.567}{0.002}{10.9091}{项}\RT{211.760}{0.002}{10.9091}{式}\RT{222.964}{0.002}{10.9091}{为}\RG{237.940}{0.002}{10.9091}{mi}{112}\RG{243.428}{0.002}{10.9091}{rm}{40}\RG{247.671}{0.002}{10.9091}{mi}{116}\RG{251.610}{0.002}{10.9091}{rm}{41}\RG{259.585}{0.002}{10.9091}{rm}{61}\RG{297.348}{0.002}{10.9091}{rm}{40}\RG{301.591}{0.002}{10.9091}{mi}{116}\RG{308.244}{0.002}{10.9091}{sy}{0}\RG{319.442}{0.002}{10.9091}{mi}{21}\RG{330.188}{0.002}{10.9091}{rm}{41}\RT{334.430}{0.002}{10.9091}{。}\RT{345.623}{0.002}{10.9091}{由}\RG{360.603}{0.002}{10.9091}{mi}{65}\RG{368.785}{0.002}{10.9091}{mi}{88}\RG{382.431}{0.002}{10.9091}{rm}{61}\RG{394.668}{0.002}{10.9091}{sy}{0}\RG{403.153}{0.002}{10.9091}{mi}{88}\RG{413.037}{0.002}{10.9091}{mi}{65}\RT{431.580}{0.002}{10.9091}{归}\RT{442.772}{0.002}{10.9091}{纳}\RT{453.976}{0.002}{10.9091}{得}\RG{93.355}{-1.634}{7.9701}{rm}{49}\RG{128.693}{-1.634}{7.9701}{mi}{110}\RG{325.805}{-1.634}{7.9701}{mi}{106}\RG{282.145}{-3.224}{7.9701}{mi}{106}\RG{286.029}{-3.224}{7.9701}{rm}{61}\RG{292.613}{-3.224}{7.9701}{rm}{49}\RG{8.182}{-11.214}{7.9701}{mi}{107}\RG{60.166}{-11.214}{7.9701}{mi}{107}\RT{87.604}{-11.214}{7.9701}{T}\RG{98.131}{-11.214}{7.9701}{mi}{107}\RG{0.000}{-15.173}{10.9091}{mi}{65}\RG{13.302}{-15.173}{10.9091}{mi}{88}\RG{26.226}{-15.173}{10.9091}{rm}{61}\RG{37.725}{-15.173}{10.9091}{rm}{40}\RG{41.983}{-15.173}{10.9091}{sy}{0}\RG{50.468}{-15.173}{10.9091}{rm}{49}\RG{55.923}{-15.173}{10.9091}{rm}{41}\RG{65.286}{-15.173}{10.9091}{mi}{88}\RG{75.180}{-15.173}{10.9091}{rm}{40}\RG{79.422}{-15.173}{10.9091}{mi}{65}\RG{93.889}{-15.173}{10.9091}{rm}{41}\RT{103.251}{-15.173}{10.9091}{，}\RT{114.160}{-15.173}{10.9091}{再}\RT{125.069}{-15.173}{10.9091}{由}\RT{139.611}{-15.173}{10.9091}{C}\RT{147.488}{-15.173}{10.9091}{a}\RT{152.637}{-15.173}{10.9091}{y}\RT{158.397}{-15.173}{10.9091}{l}\RT{161.430}{-15.173}{10.9091}{e}\RT{166.273}{-15.173}{10.9091}{y}\RT{172.033}{-15.173}{10.9091}{–}\RT{177.488}{-15.173}{10.9091}{H}\RT{185.670}{-15.173}{10.9091}{a}\RT{191.124}{-15.173}{10.9091}{m}\RT{200.212}{-15.173}{10.9091}{i}\RT{203.245}{-15.173}{10.9091}{l}\RT{206.277}{-15.173}{10.9091}{t}\RT{210.521}{-15.173}{10.9091}{o}\RT{215.976}{-15.173}{10.9091}{n}\RT{225.674}{-15.173}{10.9091}{定}\RT{236.583}{-15.173}{10.9091}{理}\RG{251.125}{-15.173}{10.9091}{mi}{112}\RG{256.613}{-15.173}{10.9091}{rm}{40}\RG{260.856}{-15.173}{10.9091}{mi}{65}\RG{269.037}{-15.173}{10.9091}{rm}{41}\RG{276.292}{-15.173}{10.9091}{rm}{61}\RG{287.801}{-15.173}{10.9091}{rm}{48}\RT{293.280}{-15.173}{10.9091}{，}\RT{304.190}{-15.173}{10.9091}{有}%
}

\[
\RB{291.413}{17.647}{15.218}{%
\RG{226.214}{13.632}{7.9701}{mi}{110}\RG{221.814}{10.359}{10.9092}{ex}{89}\RT{102.888}{4.500}{7.9701}{T}\RT{166.479}{4.500}{7.9701}{T}\RG{213.977}{4.500}{7.9701}{mi}{110}\RT{248.560}{4.500}{7.9701}{T}\RG{0.000}{-0.005}{10.9091}{rm}{48}\RG{8.476}{-0.005}{10.9091}{rm}{61}\RG{20.000}{-0.005}{10.9091}{mi}{112}\RG{25.489}{-0.005}{10.9091}{rm}{40}\RG{29.731}{-0.005}{10.9091}{mi}{65}\RG{37.913}{-0.005}{10.9091}{rm}{41}\RG{42.156}{-0.005}{10.9091}{mi}{88}\RG{55.080}{-0.005}{10.9091}{rm}{61}\RG{66.596}{-0.005}{10.9091}{mi}{88}\RG{76.479}{-0.005}{10.9091}{mi}{112}\RG{81.979}{-0.005}{10.9091}{rm}{40}\RG{86.221}{-0.005}{10.9091}{sy}{0}\RG{94.706}{-0.005}{10.9091}{mi}{65}\RG{109.172}{-0.005}{10.9091}{rm}{41}\RG{113.415}{-0.005}{10.9091}{mi}{59}\RG{140.066}{-0.005}{10.9091}{mi}{112}\RG{145.570}{-0.005}{10.9091}{rm}{40}\RG{149.812}{-0.005}{10.9091}{sy}{0}\RG{158.297}{-0.005}{10.9091}{mi}{65}\RG{172.763}{-0.005}{10.9091}{rm}{41}\RG{180.018}{-0.005}{10.9091}{rm}{61}\RG{191.527}{-0.005}{10.9091}{rm}{40}\RG{195.795}{-0.005}{10.9091}{sy}{0}\RG{204.280}{-0.005}{10.9091}{rm}{49}\RG{209.734}{-0.005}{10.9091}{rm}{41}\RG{236.136}{-0.005}{10.9091}{rm}{40}\RG{240.378}{-0.005}{10.9091}{mi}{65}\RG{257.269}{-0.005}{10.9091}{rm}{43}\RG{268.178}{-0.005}{10.9091}{mi}{21}\RG{278.924}{-0.005}{10.9091}{mi}{73}\RG{284.575}{-0.005}{10.9091}{rm}{41}\RG{288.818}{-0.005}{10.9091}{mi}{58}\RG{274.541}{-1.641}{7.9701}{mi}{106}\RG{221.431}{-13.092}{7.9701}{mi}{106}\RG{225.315}{-13.092}{7.9701}{rm}{61}\RG{231.898}{-13.092}{7.9701}{rm}{49}%
}
\]
每个\allowbreak{}\hskip.1em \RB{8.366}{8.309}{0.502}{%
\RG{0.000}{-0.002}{10.9091}{mi}{65}%
}\hskip.1em\allowbreak{} T\allowbreak{}\hskip.1em \RB{27.767}{8.069}{3.765}{%
\RG{0.000}{-0.002}{10.9091}{rm}{43}\RG{11.099}{-0.002}{10.9091}{mi}{21}\RG{17.463}{-1.639}{7.9701}{mi}{106}\RG{21.845}{-0.002}{10.9091}{mi}{73}%
}\hskip.1em\allowbreak{} 的特征值均为正数，因而可逆；于是\allowbreak{}\hskip.1em \RB{30.896}{7.949}{0.742}{%
\RG{0.000}{-0.002}{10.9091}{mi}{88}\RG{13.399}{-0.002}{10.9091}{rm}{61}\RG{25.377}{-0.002}{10.9091}{rm}{48}%
}\hskip.1em\allowbreak{} 。这里的Cayley–Hamilton定理可由\allowbreak{}\hskip.1em \RB{39.002}{8.685}{3.224}{%
\RG{0.000}{0.003}{10.9091}{rm}{40}\RG{4.243}{0.003}{10.9091}{mi}{116}\RG{8.181}{0.003}{10.9091}{mi}{73}\RG{16.258}{0.003}{10.9091}{sy}{0}\RG{27.167}{0.003}{10.9091}{mi}{65}\RG{35.349}{0.003}{10.9091}{rm}{41}%
}\hskip.1em\allowbreak{} adj\allowbreak{}\hskip.1em \RB{77.972}{8.685}{3.224}{%
\RG{0.000}{0.003}{10.9091}{rm}{40}\RG{4.243}{0.003}{10.9091}{mi}{116}\RG{8.181}{0.003}{10.9091}{mi}{73}\RG{16.258}{0.003}{10.9091}{sy}{0}\RG{27.167}{0.003}{10.9091}{mi}{65}\RG{35.349}{0.003}{10.9091}{rm}{41}\RG{42.604}{0.003}{10.9091}{rm}{61}\RG{54.137}{0.003}{10.9091}{mi}{112}\RG{59.625}{0.003}{10.9091}{rm}{40}\RG{63.869}{0.003}{10.9091}{mi}{116}\RG{67.808}{0.003}{10.9091}{rm}{41}\RG{72.051}{0.003}{10.9091}{mi}{73}%
}\hskip.1em\allowbreak{} 比较系数并递推消去得到。

(1)将已知等式转置后相减，得

\[
\RB{150.070}{11.423}{3.232}{%
\RT{43.401}{4.499}{7.9701}{T}\RT{102.482}{4.499}{7.9701}{T}\RT{121.191}{4.499}{7.9701}{T}\RG{0.000}{-0.005}{10.9091}{mi}{65}\RG{8.182}{-0.005}{10.9091}{rm}{40}\RG{12.424}{-0.005}{10.9091}{mi}{66}\RG{23.670}{-0.005}{10.9091}{sy}{0}\RG{34.579}{-0.005}{10.9091}{mi}{66}\RG{49.686}{-0.005}{10.9091}{rm}{41}\RG{56.341}{-0.005}{10.9091}{rm}{43}\RG{67.239}{-0.005}{10.9091}{rm}{40}\RG{71.505}{-0.005}{10.9091}{mi}{66}\RG{82.751}{-0.005}{10.9091}{sy}{0}\RG{93.660}{-0.005}{10.9091}{mi}{66}\RG{108.766}{-0.005}{10.9091}{rm}{41}\RG{113.009}{-0.005}{10.9091}{mi}{65}\RG{130.505}{-0.005}{10.9091}{rm}{61}\RG{142.003}{-0.005}{10.9091}{rm}{48}\RG{147.475}{-0.005}{10.9091}{mi}{58}%
}
\]
由唯一性结论，\allowbreak{}\hskip.1em \RB{32.115}{7.951}{0.500}{%
\RG{0.000}{0.000}{10.9091}{mi}{66}\RG{11.852}{0.000}{10.9091}{rm}{61}\RG{23.367}{0.000}{10.9091}{mi}{66}%
}\hskip.1em\allowbreak{} T。

(2)若\allowbreak{}\hskip.1em \RB{31.094}{8.307}{0.504}{%
\RG{0.000}{-0.004}{10.9091}{mi}{65}\RG{11.212}{-0.004}{10.9091}{rm}{61}\RG{22.728}{-0.004}{10.9091}{mi}{65}%
}\hskip.1em\allowbreak{} T，则\allowbreak{}\hskip.1em \RB{31.733}{8.307}{0.504}{%
\RG{0.000}{-0.004}{10.9091}{mi}{66}\RG{11.852}{-0.004}{10.9091}{rm}{61}\RG{23.367}{-0.004}{10.9091}{mi}{65}%
}\hskip.1em\allowbreak{} 是原方程的一个解；任意另一解与它的差满足上述齐次方程，故原来的\allowbreak{}\hskip.1em \RB{8.747}{7.952}{0.500}{%
\RG{0.000}{0.001}{10.9091}{mi}{66}%
}\hskip.1em\allowbreak{} 必等于\allowbreak{}\hskip.1em \RB{8.366}{8.312}{0.500}{%
\RG{0.000}{0.001}{10.9091}{mi}{65}%
}\hskip.1em\allowbreak{} 。反之，若\allowbreak{}\hskip.1em \RB{31.733}{8.312}{0.500}{%
\RG{0.000}{0.001}{10.9091}{mi}{66}\RG{11.852}{0.001}{10.9091}{rm}{61}\RG{23.367}{0.001}{10.9091}{mi}{65}%
}\hskip.1em\allowbreak{} ，则由(1)得\allowbreak{}\hskip.1em \RB{8.366}{8.312}{0.500}{%
\RG{0.000}{0.001}{10.9091}{mi}{65}%
}\hskip.1em\allowbreak{} 为对称阵。

再由(1)及迹的循环性质，tr\allowbreak{}\hskip.1em \RB{21.423}{8.680}{3.229}{%
\RG{0.000}{-0.002}{10.9091}{rm}{40}\RG{4.243}{-0.002}{10.9091}{mi}{66}\RG{13.058}{-0.002}{10.9091}{mi}{65}%
}\hskip.1em\allowbreak{} T\allowbreak{}\hskip.1em \RB{15.620}{8.680}{3.229}{%
\RG{0.000}{-0.002}{10.9091}{rm}{41}\RG{7.255}{-0.002}{10.9091}{rm}{61}%
}\hskip.1em\allowbreak{} tr\allowbreak{}\hskip.1em \RB{24.899}{8.680}{3.229}{%
\RG{0.000}{-0.002}{10.9091}{rm}{40}\RG{4.242}{-0.002}{10.9091}{mi}{65}\RG{12.424}{-0.002}{10.9091}{mi}{66}\RG{21.246}{-0.002}{10.9091}{rm}{41}%
}\hskip.1em\allowbreak{} 。对原方程取迹，得

\[
\RB{90.802}{11.432}{3.224}{%
\RT{77.682}{4.508}{7.9701}{T}\RT{0.000}{0.003}{10.9091}{t}\RT{4.244}{0.003}{10.9091}{r}\RG{8.520}{0.003}{10.9091}{rm}{40}\RG{12.763}{0.003}{10.9091}{mi}{65}\RG{20.945}{0.003}{10.9091}{mi}{66}\RG{29.767}{0.003}{10.9091}{rm}{41}\RG{37.022}{0.003}{10.9091}{rm}{61}\RT{48.554}{0.003}{10.9091}{t}\RT{52.798}{0.003}{10.9091}{r}\RG{57.074}{0.003}{10.9091}{rm}{40}\RG{61.318}{0.003}{10.9091}{mi}{65}\RG{69.499}{0.003}{10.9091}{mi}{65}\RG{83.966}{0.003}{10.9091}{rm}{41}\RG{88.208}{0.003}{10.9091}{mi}{58}%
}
\]
记\allowbreak{}\hskip.1em \RB{26.984}{9.758}{3.228}{%
\RG{0.000}{-0.001}{10.9091}{sy}{107}\RG{5.455}{-0.001}{10.9091}{mi}{77}\RG{17.227}{-0.001}{10.9091}{sy}{107}\RG{22.682}{3.958}{7.9701}{rm}{50}%
}\hskip.1em\allowbreak{} F\allowbreak{}\hskip.1em \RB{8.366}{4.503}{0.500}{%
\RG{0.000}{-0.001}{10.9091}{rm}{61}%
}\hskip.1em\allowbreak{} tr\allowbreak{}\hskip.1em \RB{16.131}{8.681}{3.228}{%
\RG{0.000}{-0.001}{10.9091}{rm}{40}\RG{4.242}{-0.001}{10.9091}{mi}{77}%
}\hskip.1em\allowbreak{} T\allowbreak{}\hskip.1em \RB{69.759}{9.758}{5.352}{%
\RG{0.000}{-0.001}{10.9091}{mi}{77}\RG{11.773}{-0.001}{10.9091}{rm}{41}\RG{19.028}{-0.001}{10.9091}{rm}{61}\RG{30.561}{8.181}{10.9092}{ex}{80}\RG{42.076}{-3.226}{7.9701}{mi}{105}\RG{44.954}{-3.226}{7.9701}{mi}{59}\RG{47.305}{-3.226}{7.9701}{mi}{106}\RG{53.512}{-0.001}{10.9091}{mi}{109}\RG{63.091}{3.958}{7.9701}{rm}{50}\RG{63.091}{-3.056}{7.9701}{mi}{105}\RG{65.968}{-3.056}{7.9701}{mi}{106}%
}\hskip.1em\allowbreak{} 。展开平方可得

\[
\RB{245.739}{11.424}{21.395}{%
\RG{41.246}{4.500}{7.9701}{rm}{50}\RG{83.398}{4.500}{7.9701}{rm}{50}\RT{173.992}{4.500}{7.9701}{T}\RT{227.620}{4.500}{7.9701}{T}\RG{0.000}{-0.004}{10.9091}{sy}{107}\RG{5.454}{-0.004}{10.9091}{mi}{66}\RG{16.700}{-0.004}{10.9091}{sy}{0}\RG{27.609}{-0.004}{10.9091}{mi}{65}\RG{35.791}{-0.004}{10.9091}{sy}{107}\RG{50.297}{-0.004}{10.9091}{rm}{61}\RT{61.814}{-0.004}{10.9091}{t}\RT{66.057}{-0.004}{10.9091}{r}\RG{70.334}{-0.004}{10.9091}{rm}{40}\RG{74.576}{-0.004}{10.9091}{mi}{66}\RG{88.130}{-0.004}{10.9091}{rm}{41}\RG{94.797}{-0.004}{10.9091}{sy}{0}\RT{105.706}{-0.004}{10.9091}{t}\RT{109.950}{-0.004}{10.9091}{r}\RG{114.226}{-0.004}{10.9091}{rm}{40}\RG{118.468}{-0.004}{10.9091}{mi}{66}\RG{127.283}{-0.004}{10.9091}{mi}{65}\RG{135.472}{-0.004}{10.9091}{rm}{41}\RG{142.139}{-0.004}{10.9091}{sy}{0}\RT{153.048}{-0.004}{10.9091}{t}\RT{157.292}{-0.004}{10.9091}{r}\RG{161.568}{-0.004}{10.9091}{rm}{40}\RG{165.810}{-0.004}{10.9091}{mi}{65}\RG{180.278}{-0.004}{10.9091}{mi}{66}\RG{189.100}{-0.004}{10.9091}{rm}{41}\RG{195.754}{-0.004}{10.9091}{rm}{43}\RT{206.675}{-0.004}{10.9091}{t}\RT{210.918}{-0.004}{10.9091}{r}\RG{215.195}{-0.004}{10.9091}{rm}{40}\RG{219.438}{-0.004}{10.9091}{mi}{65}\RG{233.904}{-0.004}{10.9091}{mi}{65}\RG{242.086}{-0.004}{10.9091}{rm}{41}\RT{41.246}{-2.701}{7.9701}{F}\RG{83.398}{-13.664}{7.9701}{rm}{50}\RT{134.832}{-13.664}{7.9701}{T}\RG{50.297}{-18.168}{10.9091}{rm}{61}\RT{61.814}{-18.168}{10.9091}{t}\RT{66.057}{-18.168}{10.9091}{r}\RG{70.334}{-18.168}{10.9091}{rm}{40}\RG{74.576}{-18.168}{10.9091}{mi}{66}\RG{88.130}{-18.168}{10.9091}{rm}{41}\RG{94.797}{-18.168}{10.9091}{sy}{0}\RT{105.706}{-18.168}{10.9091}{t}\RT{109.950}{-18.168}{10.9091}{r}\RG{114.226}{-18.168}{10.9091}{rm}{40}\RG{118.468}{-18.168}{10.9091}{mi}{65}\RG{126.650}{-18.168}{10.9091}{mi}{65}\RG{141.117}{-18.168}{10.9091}{rm}{41}\RG{145.359}{-18.168}{10.9091}{mi}{58}%
}
\]
因此tr\allowbreak{}\hskip.1em \RB{33.417}{9.757}{3.229}{%
\RG{0.000}{-0.002}{10.9091}{rm}{40}\RG{4.242}{-0.002}{10.9091}{mi}{66}\RG{13.064}{3.957}{7.9701}{rm}{50}\RG{17.797}{-0.002}{10.9091}{rm}{41}\RG{25.052}{-0.002}{10.9091}{rm}{61}%
}\hskip.1em\allowbreak{} tr\allowbreak{}\hskip.1em \RB{20.789}{8.680}{3.229}{%
\RG{0.000}{-0.002}{10.9091}{rm}{40}\RG{4.242}{-0.002}{10.9091}{mi}{65}\RG{12.424}{-0.002}{10.9091}{mi}{65}%
}\hskip.1em\allowbreak{} T\allowbreak{}\hskip.1em \RB{3.653}{8.680}{3.229}{%
\RG{0.000}{-0.002}{10.9091}{rm}{41}%
}\hskip.1em\allowbreak{} 当且仅当\allowbreak{}\hskip.1em \RB{31.476}{8.309}{0.502}{%
\RG{0.000}{-0.002}{10.9091}{mi}{65}\RG{11.212}{-0.002}{10.9091}{rm}{61}\RG{22.728}{-0.002}{10.9091}{mi}{66}%
}\hskip.1em\allowbreak{} ，所需三个条件等价。

(3)因\allowbreak{}\hskip.1em \RB{8.366}{8.315}{0.500}{%
\RG{0.000}{0.004}{10.9091}{mi}{65}%
}\hskip.1em\allowbreak{} 可逆，在原方程左乘\allowbreak{}\hskip.1em \RB{18.831}{9.763}{0.500}{%
\RG{0.000}{0.004}{10.9091}{mi}{65}\RG{8.182}{3.963}{7.9701}{sy}{0}\RG{14.768}{3.963}{7.9701}{rm}{49}%
}\hskip.1em\allowbreak{} 、右乘\allowbreak{}\hskip.1em \RB{14.556}{8.315}{0.500}{%
\RG{0.000}{0.004}{10.9091}{mi}{65}\RG{8.182}{3.963}{7.9701}{sy}{0}%
}\hskip.1em\allowbreak{} T，得到

\[
\RB{99.765}{11.428}{1.405}{%
\RG{8.182}{4.504}{7.9701}{sy}{0}\RG{14.768}{4.504}{7.9701}{rm}{49}\RG{58.659}{4.504}{7.9701}{sy}{0}\RT{65.247}{4.504}{7.9701}{T}\RG{0.000}{0.000}{10.9091}{mi}{65}\RG{19.500}{0.000}{10.9091}{mi}{66}\RG{30.747}{0.000}{10.9091}{rm}{43}\RG{41.656}{0.000}{10.9091}{mi}{66}\RG{50.471}{0.000}{10.9091}{mi}{65}\RG{74.562}{0.000}{10.9091}{rm}{61}\RG{86.060}{0.000}{10.9091}{rm}{50}\RG{91.531}{0.000}{10.9091}{mi}{73}\RG{97.171}{0.000}{10.9091}{mi}{58}%
}
\]
令\allowbreak{}\hskip.1em \RB{72.169}{9.754}{0.505}{%
\RG{0.000}{-0.005}{10.9091}{mi}{75}\RG{13.076}{-0.005}{10.9091}{rm}{61}\RG{24.591}{-0.005}{10.9091}{mi}{65}\RG{32.773}{3.954}{7.9701}{sy}{0}\RG{39.359}{3.954}{7.9701}{rm}{49}\RG{44.091}{-0.005}{10.9091}{mi}{66}\RG{55.338}{-0.005}{10.9091}{sy}{0}\RG{66.248}{-0.005}{10.9091}{mi}{73}%
}\hskip.1em\allowbreak{} 。由\allowbreak{}\hskip.1em \RB{32.115}{7.946}{0.505}{%
\RG{0.000}{-0.005}{10.9091}{mi}{66}\RG{11.852}{-0.005}{10.9091}{rm}{61}\RG{23.367}{-0.005}{10.9091}{mi}{66}%
}\hskip.1em\allowbreak{} T可知\allowbreak{}\hskip.1em \RB{10.198}{7.946}{0.505}{%
\RG{0.000}{-0.005}{10.9091}{mi}{75}%
}\hskip.1em\allowbreak{} T\allowbreak{}\hskip.1em \RB{30.198}{7.946}{0.505}{%
\RG{0.000}{-0.005}{10.9091}{rm}{61}\RG{11.515}{-0.005}{10.9091}{sy}{0}\RG{20.000}{-0.005}{10.9091}{mi}{75}%
}\hskip.1em\allowbreak{} ，且

\[
\RB{96.796}{16.677}{10.096}{%
\RT{0.000}{7.384}{10.9091}{d}\RT{6.066}{7.384}{10.9091}{e}\RT{10.909}{7.384}{10.9091}{t}\RG{16.971}{7.384}{10.9091}{mi}{66}\RG{30.018}{0.004}{10.9091}{rm}{61}\RT{41.533}{0.004}{10.9091}{d}\RT{47.599}{0.004}{10.9091}{e}\RT{52.442}{0.004}{10.9091}{t}\RG{56.686}{0.004}{10.9091}{rm}{40}\RG{60.929}{0.004}{10.9091}{mi}{73}\RG{69.004}{0.004}{10.9091}{rm}{43}\RG{79.913}{0.004}{10.9091}{mi}{75}\RG{89.959}{0.004}{10.9091}{rm}{41}\RG{94.201}{0.004}{10.9091}{mi}{58}\RT{0.320}{-7.480}{10.9091}{d}\RT{6.385}{-7.480}{10.9091}{e}\RT{11.229}{-7.480}{10.9091}{t}\RG{17.291}{-7.480}{10.9091}{mi}{65}\RR{-0.000}{2.731}{25.792}{0.436}%
}
\]
实反对称阵的非零特征值成对为\allowbreak{}\hskip.1em \RB{36.903}{7.713}{3.760}{%
\RG{0.000}{0.002}{10.9091}{mi}{105}\RG{3.753}{0.002}{10.9091}{mi}{116}\RG{7.698}{-1.634}{7.9701}{mi}{106}\RG{12.080}{0.002}{10.9091}{mi}{59}\RG{16.928}{0.002}{10.9091}{sy}{0}\RG{25.413}{0.002}{10.9091}{mi}{105}\RG{29.166}{0.002}{10.9091}{mi}{116}\RG{33.111}{-1.634}{7.9701}{mi}{106}%
}\hskip.1em\allowbreak{} ，\allowbreak{}\hskip.1em \RB{28.461}{7.768}{3.760}{%
\RG{0.000}{0.002}{10.9091}{mi}{116}\RG{3.940}{-1.634}{7.9701}{mi}{106}\RG{11.390}{0.002}{10.9091}{mi}{62}\RG{22.943}{0.002}{10.9091}{rm}{48}%
}\hskip.1em\allowbreak{} ；这是上一题斜Hermite阵谱定理加上实矩阵的共轭根性质。因此

\[
\RB{140.458}{17.646}{15.218}{%
\RG{70.821}{13.631}{7.9701}{mi}{109}\RG{67.596}{10.358}{10.9092}{ex}{89}\RG{108.888}{4.499}{7.9701}{rm}{50}\RT{0.000}{-0.005}{10.9091}{d}\RT{6.066}{-0.005}{10.9091}{e}\RT{10.909}{-0.005}{10.9091}{t}\RG{15.152}{-0.005}{10.9091}{rm}{40}\RG{19.395}{-0.005}{10.9091}{mi}{73}\RG{27.470}{-0.005}{10.9091}{rm}{43}\RG{38.379}{-0.005}{10.9091}{mi}{75}\RG{48.425}{-0.005}{10.9091}{rm}{41}\RG{55.680}{-0.005}{10.9091}{rm}{61}\RG{81.918}{-0.005}{10.9091}{rm}{40}\RG{86.150}{-0.005}{10.9091}{rm}{49}\RG{94.027}{-0.005}{10.9091}{rm}{43}\RG{104.949}{-0.005}{10.9091}{mi}{116}\RG{113.620}{-0.005}{10.9091}{rm}{41}\RG{120.893}{-0.005}{10.9091}{sy}{21}\RG{132.408}{-0.005}{10.9091}{rm}{49}\RG{137.863}{-0.005}{10.9091}{mi}{58}\RG{108.888}{-2.702}{7.9701}{mi}{106}\RG{67.214}{-13.092}{7.9701}{mi}{106}\RG{71.098}{-13.092}{7.9701}{rm}{61}\RG{77.681}{-13.092}{7.9701}{rm}{49}%
}
\]
又det\allowbreak{}\hskip.1em \RB{70.393}{8.679}{5.354}{%
\RG{0.000}{-0.003}{10.9091}{mi}{65}\RG{11.300}{-0.003}{10.9091}{rm}{61}\RG{22.904}{8.179}{10.9092}{ex}{81}\RG{33.207}{-3.228}{7.9701}{mi}{106}\RG{39.407}{-0.003}{10.9091}{mi}{21}\RG{45.771}{-1.639}{7.9701}{mi}{106}\RG{53.271}{-0.003}{10.9091}{mi}{62}\RG{64.875}{-0.003}{10.9091}{rm}{48}%
}\hskip.1em\allowbreak{} ，故det\allowbreak{}\hskip.1em \RB{20.011}{7.948}{1.998}{%
\RG{0.000}{-0.003}{10.9091}{mi}{66}\RG{11.940}{-0.003}{10.9091}{sy}{21}%
}\hskip.1em\allowbreak{} det\allowbreak{}\hskip.1em \RB{8.366}{8.308}{0.503}{%
\RG{0.000}{-0.003}{10.9091}{mi}{65}%
}\hskip.1em\allowbreak{} 。等号成立当且仅当\allowbreak{}\hskip.1em \RB{10.198}{7.948}{0.503}{%
\RG{0.000}{-0.003}{10.9091}{mi}{75}%
}\hskip.1em\allowbreak{} 的全部特征值为零；\allowbreak{}\hskip.1em \RB{10.198}{7.948}{0.503}{%
\RG{0.000}{-0.003}{10.9091}{mi}{75}%
}\hskip.1em\allowbreak{} 可酉对角化，故这等价于\allowbreak{}\hskip.1em \RB{30.093}{7.953}{0.738}{%
\RG{0.000}{0.002}{10.9091}{mi}{75}\RG{13.076}{0.002}{10.9091}{rm}{61}\RG{24.574}{0.002}{10.9091}{rm}{48}%
}\hskip.1em\allowbreak{} ，亦即\allowbreak{}\hskip.1em \RB{31.733}{8.313}{0.500}{%
\RG{0.000}{0.002}{10.9091}{mi}{66}\RG{11.852}{0.002}{10.9091}{rm}{61}\RG{23.367}{0.002}{10.9091}{mi}{65}%
}\hskip.1em\allowbreak{} 。\allowbreak{}\hskip.1em \RB{8.366}{7.997}{0.500}{%
\RG{0.000}{0.002}{10.9091}{am}{3}%
}\hskip.1em\allowbreak{}

方法小结：先以特征多项式证明齐次矩阵方程只有零解；迹的等号化为Frobenius范数为零，行列式的等号化为反对称阵为零。题目来源：谢启鸿《复旦大学高等代数习题集（2023版）》，第117页。

% Original page 3.
\section*{2022S18}
设\allowbreak{}\hskip.1em \RB{8.366}{8.311}{0.500}{%
\RG{0.000}{0.000}{10.9091}{mi}{65}%
}\hskip.1em\allowbreak{} 为\allowbreak{}\hskip.1em \RB{6.729}{5.322}{0.620}{%
\RG{0.000}{0.000}{10.9091}{mi}{110}%
}\hskip.1em\allowbreak{} 阶实方阵，证明

\[
\RB{132.026}{11.913}{4.313}{%
\RG{0.000}{8.831}{10.9092}{ex}{0}\RG{23.520}{8.831}{10.9092}{ex}{1}\RG{28.520}{6.113}{7.9701}{rm}{50}\RT{110.602}{4.499}{7.9701}{T}\RT{5.000}{-0.005}{10.9091}{t}\RT{9.244}{-0.005}{10.9091}{r}\RG{15.338}{-0.005}{10.9091}{mi}{65}\RG{36.282}{-0.005}{10.9091}{sy}{20}\RT{47.797}{-0.005}{10.9091}{r}\RT{52.074}{-0.005}{10.9091}{a}\RT{57.528}{-0.005}{10.9091}{n}\RT{63.594}{-0.005}{10.9091}{k}\RG{69.355}{-0.005}{10.9091}{rm}{40}\RG{73.597}{-0.005}{10.9091}{mi}{65}\RG{81.779}{-0.005}{10.9091}{rm}{41}\RT{89.658}{-0.005}{10.9091}{t}\RT{93.901}{-0.005}{10.9091}{r}\RG{98.178}{-0.005}{10.9091}{rm}{40}\RG{102.420}{-0.005}{10.9091}{mi}{65}\RG{116.887}{-0.005}{10.9091}{mi}{65}\RG{125.068}{-0.005}{10.9091}{rm}{41}\RG{129.311}{-0.005}{10.9091}{mi}{59}%
}
\]
并求等号成立的充要条件。

证明：记\allowbreak{}\hskip.1em \RB{16.621}{5.322}{0.620}{%
\RG{0.000}{-0.000}{10.9091}{mi}{114}\RG{8.255}{-0.000}{10.9091}{rm}{61}%
}\hskip.1em\allowbreak{} rank\allowbreak{}\hskip.1em \RB{16.077}{8.682}{3.227}{%
\RG{0.000}{-0.000}{10.9091}{rm}{40}\RG{4.242}{-0.000}{10.9091}{mi}{65}\RG{12.424}{-0.000}{10.9091}{rm}{41}%
}\hskip.1em\allowbreak{} 。当\allowbreak{}\hskip.1em \RB{25.272}{7.766}{0.740}{%
\RG{0.000}{-0.000}{10.9091}{mi}{114}\RG{8.255}{-0.000}{10.9091}{rm}{61}\RG{19.753}{-0.000}{10.9091}{rm}{48}%
}\hskip.1em\allowbreak{} 时，\allowbreak{}\hskip.1em \RB{28.229}{8.311}{0.740}{%
\RG{0.000}{-0.000}{10.9091}{mi}{65}\RG{11.212}{-0.000}{10.9091}{rm}{61}\RG{22.710}{-0.000}{10.9091}{rm}{48}%
}\hskip.1em\allowbreak{} ，等号成立。以下设\allowbreak{}\hskip.1em \RB{25.288}{7.766}{0.936}{%
\RG{0.000}{-0.000}{10.9091}{mi}{114}\RG{8.247}{-0.000}{10.9091}{mi}{62}\RG{19.770}{-0.000}{10.9091}{rm}{48}%
}\hskip.1em\allowbreak{} 。

取Im\allowbreak{}\hskip.1em \RB{8.366}{8.307}{0.504}{%
\RG{0.000}{-0.004}{10.9091}{mi}{65}%
}\hskip.1em\allowbreak{} 的一组标准正交基\allowbreak{}\hskip.1em \RB{45.647}{5.318}{2.609}{%
\RG{0.000}{-0.004}{10.9091}{mi}{117}\RG{6.245}{-1.640}{7.9701}{rm}{49}\RG{10.977}{-0.004}{10.9091}{mi}{59}\RG{15.810}{-0.004}{10.9091}{mi}{58}\RG{20.654}{-0.004}{10.9091}{mi}{58}\RG{25.486}{-0.004}{10.9091}{mi}{58}\RG{30.330}{-0.004}{10.9091}{mi}{59}\RG{35.174}{-0.004}{10.9091}{mi}{117}\RG{41.464}{-1.640}{7.9701}{mi}{114}%
}\hskip.1em\allowbreak{} ，令

\[
\RB{66.392}{17.652}{15.212}{%
\RG{28.915}{13.637}{7.9701}{mi}{114}\RG{23.065}{10.365}{10.9092}{ex}{88}\RT{57.513}{4.505}{7.9701}{T}\RG{0.000}{0.001}{10.9091}{mi}{80}\RG{11.550}{0.001}{10.9091}{rm}{61}\RG{40.641}{0.001}{10.9091}{mi}{117}\RG{51.268}{0.001}{10.9091}{mi}{117}\RG{63.798}{0.001}{10.9091}{mi}{58}\RG{46.885}{-1.635}{7.9701}{mi}{106}\RG{57.513}{-2.696}{7.9701}{mi}{106}\RG{23.591}{-13.086}{7.9701}{mi}{106}\RG{27.475}{-13.086}{7.9701}{rm}{61}\RG{34.059}{-13.086}{7.9701}{rm}{49}%
}
\]
则\allowbreak{}\hskip.1em \RB{8.726}{7.948}{0.503}{%
\RG{0.000}{-0.003}{10.9091}{mi}{80}%
}\hskip.1em\allowbreak{} 为到Im\allowbreak{}\hskip.1em \RB{8.366}{8.308}{0.503}{%
\RG{0.000}{-0.003}{10.9091}{mi}{65}%
}\hskip.1em\allowbreak{} 的正交投影，满足

\[
\RB{231.811}{11.431}{2.613}{%
\RT{8.519}{4.507}{7.9701}{T}\RG{71.842}{4.507}{7.9701}{rm}{50}\RG{0.000}{0.003}{10.9091}{mi}{80}\RG{17.833}{0.003}{10.9091}{rm}{61}\RG{29.348}{0.003}{10.9091}{mi}{80}\RG{36.646}{0.003}{10.9091}{mi}{59}\RG{63.309}{0.003}{10.9091}{mi}{80}\RG{79.604}{0.003}{10.9091}{rm}{61}\RG{91.120}{0.003}{10.9091}{mi}{80}\RG{98.418}{0.003}{10.9091}{mi}{59}\RG{125.080}{0.003}{10.9091}{mi}{80}\RG{133.600}{0.003}{10.9091}{mi}{65}\RG{144.824}{0.003}{10.9091}{rm}{61}\RG{156.339}{0.003}{10.9091}{mi}{65}\RG{164.521}{0.003}{10.9091}{mi}{59}\RT{191.189}{0.003}{10.9091}{t}\RT{195.432}{0.003}{10.9091}{r}\RG{201.527}{0.003}{10.9091}{mi}{80}\RG{213.076}{0.003}{10.9091}{rm}{61}\RG{224.591}{0.003}{10.9091}{mi}{114}\RG{229.216}{0.003}{10.9091}{mi}{58}%
}
\]
在实矩阵空间上使用Frobenius内积

\[
\RB{154.718}{11.434}{15.208}{%
\RG{108.433}{10.369}{10.9092}{ex}{88}\RT{74.603}{4.510}{7.9701}{T}\RG{0.000}{0.005}{10.9091}{sy}{104}\RG{4.242}{0.005}{10.9091}{mi}{88}\RG{13.515}{0.005}{10.9091}{mi}{59}\RG{18.359}{0.005}{10.9091}{mi}{89}\RG{27.136}{0.005}{10.9091}{sy}{105}\RG{40.430}{0.005}{10.9091}{rm}{61}\RT{51.946}{0.005}{10.9091}{t}\RT{56.190}{0.005}{10.9091}{r}\RG{60.466}{0.005}{10.9091}{rm}{40}\RG{64.709}{0.005}{10.9091}{mi}{88}\RG{80.887}{0.005}{10.9091}{mi}{89}\RG{89.645}{0.005}{10.9091}{rm}{41}\RG{96.899}{0.005}{10.9091}{rm}{61}\RG{126.009}{0.005}{10.9091}{mi}{120}\RG{139.509}{0.005}{10.9091}{mi}{121}\RG{152.123}{0.005}{10.9091}{mi}{58}\RT{31.379}{-1.657}{7.9701}{F}\RG{132.243}{-1.631}{7.9701}{mi}{105}\RG{135.120}{-1.631}{7.9701}{mi}{106}\RG{144.857}{-1.631}{7.9701}{mi}{105}\RG{147.734}{-1.631}{7.9701}{mi}{106}\RG{111.752}{-13.082}{7.9701}{mi}{105}\RG{114.629}{-13.082}{7.9701}{mi}{59}\RG{116.980}{-13.082}{7.9701}{mi}{106}%
}
\]
因为\allowbreak{}\hskip.1em \RB{39.612}{8.312}{0.500}{%
\RG{0.000}{0.001}{10.9091}{mi}{80}\RG{8.509}{0.001}{10.9091}{mi}{65}\RG{19.731}{0.001}{10.9091}{rm}{61}\RG{31.246}{0.001}{10.9091}{mi}{65}%
}\hskip.1em\allowbreak{} ，所以

\[
\RB{125.040}{11.425}{3.712}{%
\RT{54.346}{4.501}{7.9701}{T}\RT{0.000}{-0.004}{10.9091}{t}\RT{4.244}{-0.004}{10.9091}{r}\RG{10.338}{-0.004}{10.9091}{mi}{65}\RG{21.550}{-0.004}{10.9091}{rm}{61}\RT{33.065}{-0.004}{10.9091}{t}\RT{37.309}{-0.004}{10.9091}{r}\RG{41.585}{-0.004}{10.9091}{rm}{40}\RG{45.827}{-0.004}{10.9091}{mi}{80}\RG{60.631}{-0.004}{10.9091}{mi}{65}\RG{68.814}{-0.004}{10.9091}{rm}{41}\RG{76.068}{-0.004}{10.9091}{rm}{61}\RG{87.601}{-0.004}{10.9091}{sy}{104}\RG{91.844}{-0.004}{10.9091}{mi}{80}\RG{99.142}{-0.004}{10.9091}{mi}{59}\RG{103.986}{-0.004}{10.9091}{mi}{65}\RG{112.181}{-0.004}{10.9091}{sy}{105}\RG{122.445}{-0.004}{10.9091}{mi}{58}\RT{116.423}{-1.666}{7.9701}{F}%
}
\]
对\allowbreak{}\hskip.1em \RB{20.507}{8.307}{2.609}{%
\RG{0.000}{-0.004}{10.9091}{mi}{80}\RG{7.298}{-0.004}{10.9091}{mi}{59}\RG{12.142}{-0.004}{10.9091}{mi}{65}%
}\hskip.1em\allowbreak{} 的全部元素应用Cauchy–Schwarz不等式，得到

\[
\RB{186.999}{30.080}{3.229}{%
\RG{0.000}{26.998}{10.9092}{ex}{0}\RG{23.520}{26.998}{10.9092}{ex}{1}\RG{28.520}{24.280}{7.9701}{rm}{50}\RT{5.000}{18.162}{10.9091}{t}\RT{9.244}{18.162}{10.9091}{r}\RG{15.339}{18.162}{10.9091}{mi}{65}\RG{36.283}{18.162}{10.9091}{sy}{20}\RG{47.782}{18.162}{10.9091}{sy}{104}\RG{52.041}{18.162}{10.9091}{mi}{80}\RG{59.340}{18.162}{10.9091}{mi}{59}\RG{64.183}{18.162}{10.9091}{mi}{80}\RG{72.716}{18.162}{10.9091}{sy}{105}\RG{84.798}{18.162}{10.9091}{sy}{104}\RG{89.040}{18.162}{10.9091}{mi}{65}\RG{97.222}{18.162}{10.9091}{mi}{59}\RG{102.054}{18.162}{10.9091}{mi}{65}\RG{110.252}{18.162}{10.9091}{sy}{105}\RT{76.958}{16.500}{7.9701}{F}\RT{114.495}{16.500}{7.9701}{F}\RG{69.081}{4.502}{7.9701}{rm}{50}\RT{102.636}{4.502}{7.9701}{T}\RT{165.695}{4.502}{7.9701}{T}\RG{36.283}{-0.002}{10.9091}{rm}{61}\RT{47.798}{-0.002}{10.9091}{t}\RT{52.042}{-0.002}{10.9091}{r}\RG{56.318}{-0.002}{10.9091}{rm}{40}\RG{60.562}{-0.002}{10.9091}{mi}{80}\RG{73.813}{-0.002}{10.9091}{rm}{41}\RT{81.692}{-0.002}{10.9091}{t}\RT{85.935}{-0.002}{10.9091}{r}\RG{90.212}{-0.002}{10.9091}{rm}{40}\RG{94.454}{-0.002}{10.9091}{mi}{65}\RG{108.920}{-0.002}{10.9091}{mi}{65}\RG{117.102}{-0.002}{10.9091}{rm}{41}\RG{124.357}{-0.002}{10.9091}{rm}{61}\RG{135.890}{-0.002}{10.9091}{mi}{114}\RT{144.751}{-0.002}{10.9091}{t}\RT{148.995}{-0.002}{10.9091}{r}\RG{153.271}{-0.002}{10.9091}{rm}{40}\RG{157.513}{-0.002}{10.9091}{mi}{65}\RG{171.980}{-0.002}{10.9091}{mi}{65}\RG{180.162}{-0.002}{10.9091}{rm}{41}\RG{184.405}{-0.002}{10.9091}{mi}{58}%
}
\]
由于\allowbreak{}\hskip.1em \RB{28.566}{8.306}{2.861}{%
\RG{0.000}{-0.005}{10.9091}{mi}{80}\RG{11.549}{-0.005}{10.9091}{sy}{54}\RG{11.549}{-0.005}{10.9091}{rm}{61}\RG{23.047}{-0.005}{10.9091}{rm}{48}%
}\hskip.1em\allowbreak{} ，Cauchy–Schwarz不等式取等当且仅当存在实数\allowbreak{}\hskip.1em \RB{5.191}{5.317}{0.625}{%
\RG{0.000}{-0.005}{10.9091}{mi}{99}%
}\hskip.1em\allowbreak{} ，使\allowbreak{}\hskip.1em \RB{36.166}{8.306}{0.625}{%
\RG{0.000}{-0.005}{10.9091}{mi}{65}\RG{11.212}{-0.005}{10.9091}{rm}{61}\RG{22.727}{-0.005}{10.9091}{mi}{99}\RG{27.440}{-0.005}{10.9091}{mi}{80}%
}\hskip.1em\allowbreak{} 。这里\allowbreak{}\hskip.1em \RB{28.229}{8.306}{2.861}{%
\RG{0.000}{-0.005}{10.9091}{mi}{65}\RG{11.212}{-0.005}{10.9091}{sy}{54}\RG{11.212}{-0.005}{10.9091}{rm}{61}\RG{22.710}{-0.005}{10.9091}{rm}{48}%
}\hskip.1em\allowbreak{} ，故\allowbreak{}\hskip.1em \RB{24.768}{8.311}{2.856}{%
\RG{0.000}{0.000}{10.9091}{mi}{99}\RG{7.751}{0.000}{10.9091}{sy}{54}\RG{7.751}{0.000}{10.9091}{rm}{61}\RG{19.249}{0.000}{10.9091}{rm}{48}%
}\hskip.1em\allowbreak{} 。反之，若\allowbreak{}\hskip.1em \RB{36.165}{8.311}{0.620}{%
\RG{0.000}{0.000}{10.9091}{mi}{65}\RG{11.212}{0.000}{10.9091}{rm}{61}\RG{22.727}{0.000}{10.9091}{mi}{99}\RG{27.440}{0.000}{10.9091}{mi}{80}%
}\hskip.1em\allowbreak{} ，其中\allowbreak{}\hskip.1em \RB{8.726}{7.951}{0.500}{%
\RG{0.000}{0.000}{10.9091}{mi}{80}%
}\hskip.1em\allowbreak{} 为秩\allowbreak{}\hskip.1em \RB{5.256}{5.322}{0.620}{%
\RG{0.000}{0.000}{10.9091}{mi}{114}%
}\hskip.1em\allowbreak{} 的正交投影、\allowbreak{}\hskip.1em \RB{24.768}{8.311}{2.856}{%
\RG{0.000}{0.000}{10.9091}{mi}{99}\RG{7.751}{0.000}{10.9091}{sy}{54}\RG{7.751}{0.000}{10.9091}{rm}{61}\RG{19.249}{0.000}{10.9091}{rm}{48}%
}\hskip.1em\allowbreak{} ，则

\[
\RB{140.064}{11.424}{3.231}{%
\RT{90.017}{4.500}{7.9701}{T}\RG{127.992}{4.500}{7.9701}{rm}{50}\RT{0.000}{-0.004}{10.9091}{t}\RT{4.244}{-0.004}{10.9091}{r}\RG{10.338}{-0.004}{10.9091}{mi}{65}\RG{21.550}{-0.004}{10.9091}{rm}{61}\RG{33.065}{-0.004}{10.9091}{mi}{99}\RG{37.778}{-0.004}{10.9091}{mi}{114}\RG{42.404}{-0.004}{10.9091}{mi}{59}\RT{69.072}{-0.004}{10.9091}{t}\RT{73.315}{-0.004}{10.9091}{r}\RG{77.592}{-0.004}{10.9091}{rm}{40}\RG{81.835}{-0.004}{10.9091}{mi}{65}\RG{96.301}{-0.004}{10.9091}{mi}{65}\RG{104.483}{-0.004}{10.9091}{rm}{41}\RG{111.737}{-0.004}{10.9091}{rm}{61}\RG{123.271}{-0.004}{10.9091}{mi}{99}\RG{132.724}{-0.004}{10.9091}{mi}{114}\RG{137.350}{-0.004}{10.9091}{mi}{59}%
}
\]
原式确实取等。

综上，等号条件为：\allowbreak{}\hskip.1em \RB{28.229}{8.309}{0.742}{%
\RG{0.000}{-0.002}{10.9091}{mi}{65}\RG{11.212}{-0.002}{10.9091}{rm}{61}\RG{22.710}{-0.002}{10.9091}{rm}{48}%
}\hskip.1em\allowbreak{} ，或者\allowbreak{}\hskip.1em \RB{8.366}{8.309}{0.502}{%
\RG{0.000}{-0.002}{10.9091}{mi}{65}%
}\hskip.1em\allowbreak{} 是某个非零正交投影矩阵的非零实数倍。等价地，当\allowbreak{}\hskip.1em \RB{25.288}{7.764}{0.938}{%
\RG{0.000}{-0.002}{10.9091}{mi}{114}\RG{8.247}{-0.002}{10.9091}{mi}{62}\RG{19.770}{-0.002}{10.9091}{rm}{48}%
}\hskip.1em\allowbreak{}时，存在正交阵\allowbreak{}\hskip.1em \RB{8.573}{8.194}{2.613}{%
\RG{0.000}{0.003}{10.9091}{mi}{81}%
}\hskip.1em\allowbreak{} 及实数\allowbreak{}\hskip.1em \RB{24.769}{8.314}{2.853}{%
\RG{0.000}{0.003}{10.9091}{mi}{99}\RG{7.751}{0.003}{10.9091}{sy}{54}\RG{7.752}{0.003}{10.9091}{rm}{61}\RG{19.250}{0.003}{10.9091}{rm}{48}%
}\hskip.1em\allowbreak{} ，使

\[
\RB{99.252}{19.588}{14.129}{%
\RG{33.170}{18.651}{10.9092}{ex}{32}\RG{71.295}{18.651}{10.9092}{ex}{33}\RG{41.806}{7.276}{10.9091}{mi}{99}\RG{46.519}{7.276}{10.9091}{mi}{73}\RG{65.841}{7.276}{10.9091}{rm}{48}\RG{51.322}{5.639}{7.9701}{mi}{114}\RT{90.374}{4.500}{7.9701}{T}\RG{0.000}{-0.004}{10.9091}{mi}{65}\RG{11.212}{-0.004}{10.9091}{rm}{61}\RG{22.727}{-0.004}{10.9091}{mi}{81}\RG{81.750}{-0.004}{10.9091}{mi}{81}\RG{96.658}{-0.004}{10.9091}{mi}{58}\RG{46.114}{-7.899}{10.9091}{rm}{48}\RG{65.838}{-7.899}{10.9091}{rm}{48}%
}
\]
也就是说，\allowbreak{}\hskip.1em \RB{8.366}{8.313}{0.500}{%
\RG{0.000}{0.002}{10.9091}{mi}{65}%
}\hskip.1em\allowbreak{} 必须对称，并且其全部非零特征值都相等。\allowbreak{}\hskip.1em \RB{8.366}{7.997}{0.500}{%
\RG{0.000}{0.002}{10.9091}{am}{3}%
}\hskip.1em\allowbreak{}

方法小结：把恒等矩阵换成\allowbreak{}\hskip.1em \RB{8.366}{8.311}{0.500}{%
\RG{0.000}{-0.000}{10.9091}{mi}{65}%
}\hskip.1em\allowbreak{} 的像空间上的正交投影，便将通常估计中的维数\allowbreak{}\hskip.1em \RB{6.729}{5.322}{0.620}{%
\RG{0.000}{-0.000}{10.9091}{mi}{110}%
}\hskip.1em\allowbreak{} 降为秩\allowbreak{}\hskip.1em \RB{5.256}{5.322}{0.620}{%
\RG{0.000}{-0.000}{10.9091}{mi}{114}%
}\hskip.1em\allowbreak{} ；取等条件直接来自内积不等式。

题目来源：谢启鸿《复旦大学高等代数习题集（2023版）》，第118页。

% Original page 4.
\section*{2022A01}
求\allowbreak{}\hskip.1em \RB{6.729}{5.322}{0.620}{%
\RG{0.000}{0.000}{10.9091}{mi}{110}%
}\hskip.1em\allowbreak{} 阶行列式

\[
\RB{296.543}{16.316}{13.295}{%
\RG{44.366}{15.380}{10.9092}{ex}{18}\RG{92.923}{15.380}{10.9092}{ex}{19}\RG{69.933}{7.378}{10.9091}{rm}{49}\RG{0.000}{-0.002}{10.9091}{mi}{68}\RG{17.698}{-0.002}{10.9091}{rm}{61}\RT{29.213}{-0.002}{10.9091}{d}\RT{35.279}{-0.002}{10.9091}{e}\RT{40.122}{-0.002}{10.9091}{t}\RG{134.934}{-0.002}{10.9091}{mi}{59}\RG{161.601}{-0.002}{10.9091}{rm}{49}\RG{169.480}{-0.002}{10.9091}{sy}{0}\RG{180.389}{-0.002}{10.9091}{mi}{120}\RG{190.006}{-0.002}{10.9091}{mi}{121}\RG{202.767}{-0.002}{10.9091}{sy}{54}\RG{202.767}{-0.002}{10.9091}{rm}{61}\RG{214.265}{-0.002}{10.9091}{rm}{48}\RG{230.629}{-0.002}{10.9091}{rm}{40}\RG{234.861}{-0.002}{10.9091}{rm}{49}\RG{243.373}{-0.002}{10.9091}{sy}{20}\RG{254.888}{-0.002}{10.9091}{mi}{105}\RG{258.641}{-0.002}{10.9091}{mi}{59}\RG{263.474}{-0.002}{10.9091}{mi}{106}\RG{271.642}{-0.002}{10.9091}{sy}{20}\RG{283.157}{-0.002}{10.9091}{mi}{110}\RG{289.705}{-0.002}{10.9091}{rm}{41}\RG{293.948}{-0.002}{10.9091}{mi}{58}\RG{9.032}{-1.638}{7.9701}{mi}{110}\RG{186.625}{-1.638}{7.9701}{mi}{105}\RG{195.355}{-1.638}{7.9701}{mi}{106}\RG{53.592}{-7.485}{10.9091}{rm}{49}\RG{61.472}{-7.485}{10.9091}{sy}{0}\RG{72.381}{-7.485}{10.9091}{mi}{120}\RG{81.997}{-7.485}{10.9091}{mi}{121}\RG{78.616}{-9.121}{7.9701}{mi}{105}\RG{87.345}{-9.121}{7.9701}{mi}{106}\RG{100.953}{-11.169}{7.9701}{rm}{49}\RG{105.187}{-11.169}{7.9701}{sy}{20}\RG{111.774}{-11.169}{7.9701}{mi}{105}\RG{114.651}{-11.169}{7.9701}{mi}{59}\RG{117.002}{-11.169}{7.9701}{mi}{106}\RG{120.893}{-11.169}{7.9701}{sy}{20}\RG{127.480}{-11.169}{7.9701}{mi}{110}\RR{53.592}{2.725}{38.134}{0.436}%
}
\]
解：结果为

\[
\RB{163.231}{8.683}{21.864}{%
\RG{30.409}{8.183}{10.9092}{ex}{81}\RG{77.111}{0.001}{10.9091}{rm}{40}\RG{81.353}{0.001}{10.9091}{mi}{120}\RG{94.394}{0.001}{10.9091}{sy}{0}\RG{105.303}{0.001}{10.9091}{mi}{120}\RG{114.919}{0.001}{10.9091}{rm}{41}\RG{119.152}{0.001}{10.9091}{rm}{40}\RG{123.404}{0.001}{10.9091}{mi}{121}\RG{135.559}{0.001}{10.9091}{sy}{0}\RG{146.468}{0.001}{10.9091}{mi}{121}\RG{155.198}{0.001}{10.9091}{rm}{41}\RG{87.588}{-1.636}{7.9701}{mi}{106}\RG{111.538}{-1.636}{7.9701}{mi}{105}\RG{128.753}{-1.636}{7.9701}{mi}{106}\RG{151.817}{-1.636}{7.9701}{mi}{105}\RG{40.712}{-3.225}{7.9701}{rm}{49}\RG{44.947}{-3.225}{7.9701}{sy}{20}\RG{51.533}{-3.225}{7.9701}{mi}{105}\RG{54.411}{-3.225}{7.9701}{mi}{60}\RG{60.994}{-3.225}{7.9701}{mi}{106}\RG{64.888}{-3.225}{7.9701}{sy}{20}\RG{71.474}{-3.225}{7.9701}{mi}{110}\RG{45.700}{-8.330}{10.9092}{ex}{81}\RG{72.024}{-8.330}{10.9092}{ex}{81}\RG{0.000}{-9.029}{10.9091}{mi}{68}\RG{17.699}{-9.029}{10.9091}{rm}{61}\RG{160.636}{-9.029}{10.9091}{mi}{58}\RG{9.032}{-10.665}{7.9701}{mi}{110}\RG{56.003}{-11.485}{7.9701}{mi}{110}\RG{82.327}{-11.485}{7.9701}{mi}{110}\RG{97.530}{-16.512}{10.9091}{rm}{40}\RG{101.763}{-16.512}{10.9091}{rm}{49}\RG{109.651}{-16.512}{10.9091}{sy}{0}\RG{120.560}{-16.512}{10.9091}{mi}{120}\RG{130.177}{-16.512}{10.9091}{mi}{121}\RG{139.907}{-16.512}{10.9091}{rm}{41}\RG{126.795}{-18.149}{7.9701}{mi}{105}\RG{135.525}{-18.149}{7.9701}{mi}{106}\RG{56.003}{-19.738}{7.9701}{mi}{105}\RG{58.886}{-19.738}{7.9701}{rm}{61}\RG{65.470}{-19.738}{7.9701}{rm}{49}\RG{82.327}{-19.738}{7.9701}{mi}{106}\RG{86.211}{-19.738}{7.9701}{rm}{61}\RG{92.794}{-19.738}{7.9701}{rm}{49}\RR{30.409}{-6.302}{129.031}{0.436}%
}
\]
下面给出不要求\allowbreak{}\hskip.1em \RB{23.605}{5.326}{3.758}{%
\RG{0.000}{0.004}{10.9091}{mi}{120}\RG{6.235}{-1.632}{7.9701}{mi}{105}\RG{9.616}{0.004}{10.9091}{mi}{59}\RG{14.449}{0.004}{10.9091}{mi}{121}\RG{19.813}{-1.632}{7.9701}{mi}{106}%
}\hskip.1em\allowbreak{} 非零或互异的推导。

记

\[
\RB{326.085}{10.859}{15.218}{%
\RG{23.169}{10.359}{10.9092}{ex}{89}\RG{148.752}{10.359}{10.9092}{ex}{89}\RG{269.272}{10.359}{10.9092}{ex}{89}\RG{0.000}{-0.005}{10.9091}{mi}{81}\RG{11.654}{-0.005}{10.9091}{rm}{61}\RG{37.109}{-0.005}{10.9091}{rm}{40}\RG{41.342}{-0.005}{10.9091}{rm}{49}\RG{49.230}{-0.005}{10.9091}{sy}{0}\RG{60.139}{-0.005}{10.9091}{mi}{120}\RG{69.756}{-0.005}{10.9091}{mi}{121}\RG{79.487}{-0.005}{10.9091}{rm}{41}\RG{83.729}{-0.005}{10.9091}{mi}{59}\RG{110.396}{-0.005}{10.9091}{rm}{1}\RG{119.483}{-0.005}{10.9091}{rm}{40}\RG{123.729}{-0.005}{10.9091}{mi}{120}\RG{129.964}{-0.005}{10.9091}{rm}{41}\RG{137.219}{-0.005}{10.9091}{rm}{61}\RG{162.692}{-0.005}{10.9091}{rm}{40}\RG{166.934}{-0.005}{10.9091}{mi}{120}\RG{179.975}{-0.005}{10.9091}{sy}{0}\RG{190.885}{-0.005}{10.9091}{mi}{120}\RG{200.501}{-0.005}{10.9091}{rm}{41}\RG{204.744}{-0.005}{10.9091}{mi}{59}\RG{231.410}{-0.005}{10.9091}{rm}{1}\RG{240.497}{-0.005}{10.9091}{rm}{40}\RG{244.744}{-0.005}{10.9091}{mi}{121}\RG{250.484}{-0.005}{10.9091}{rm}{41}\RG{257.739}{-0.005}{10.9091}{rm}{61}\RG{283.211}{-0.005}{10.9091}{rm}{40}\RG{287.453}{-0.005}{10.9091}{mi}{121}\RG{299.608}{-0.005}{10.9091}{sy}{0}\RG{310.518}{-0.005}{10.9091}{mi}{121}\RG{319.248}{-0.005}{10.9091}{rm}{41}\RG{323.491}{-0.005}{10.9091}{mi}{58}\RG{66.375}{-1.641}{7.9701}{mi}{105}\RG{75.105}{-1.641}{7.9701}{mi}{106}\RG{173.169}{-1.641}{7.9701}{mi}{106}\RG{197.120}{-1.641}{7.9701}{mi}{105}\RG{292.802}{-1.641}{7.9701}{mi}{106}\RG{315.867}{-1.641}{7.9701}{mi}{105}\RG{25.579}{-13.092}{7.9701}{mi}{105}\RG{28.456}{-13.092}{7.9701}{mi}{59}\RG{30.808}{-13.092}{7.9701}{mi}{106}\RG{149.045}{-13.092}{7.9701}{mi}{105}\RG{151.922}{-13.092}{7.9701}{mi}{60}\RG{158.506}{-13.092}{7.9701}{mi}{106}\RG{269.564}{-13.092}{7.9701}{mi}{105}\RG{272.441}{-13.092}{7.9701}{mi}{60}\RG{279.025}{-13.092}{7.9701}{mi}{106}%
}
\]
按行列式展开，\allowbreak{}\hskip.1em \RB{46.110}{8.191}{2.616}{%
\RG{0.000}{0.000}{10.9091}{mi}{70}\RG{11.561}{0.000}{10.9091}{rm}{61}\RG{23.076}{0.000}{10.9091}{mi}{81}\RG{31.694}{0.000}{10.9091}{mi}{68}\RG{40.732}{-1.636}{7.9701}{mi}{110}%
}\hskip.1em\allowbreak{} 是多项式；每一项在每行、每列中都恰好消去一个分母，故deg\allowbreak{}\hskip.1em \RB{53.709}{7.918}{4.520}{%
\RG{0.000}{-2.745}{7.9701}{mi}{120}\RG{4.766}{-3.960}{5.9776}{mi}{105}\RG{10.244}{0.000}{10.9091}{mi}{70}\RG{21.805}{0.000}{10.9091}{sy}{20}\RG{33.320}{0.000}{10.9091}{mi}{110}\RG{40.010}{0.000}{10.9091}{sy}{0}\RG{48.638}{0.000}{10.9091}{rm}{49}%
}\hskip.1em\allowbreak{} 、deg\allowbreak{}\hskip.1em \RB{59.102}{7.913}{5.684}{%
\RG{0.000}{-2.750}{7.9701}{mi}{121}\RG{4.185}{-3.965}{5.9776}{mi}{106}\RG{10.404}{-0.005}{10.9091}{mi}{70}\RG{22.204}{-0.005}{10.9091}{sy}{20}\RG{33.959}{-0.005}{10.9091}{mi}{110}\RG{43.027}{-0.005}{10.9091}{sy}{0}\RG{54.031}{-0.005}{10.9091}{rm}{49}%
}\hskip.1em\allowbreak{} 。交换任意两个\allowbreak{}\hskip.1em \RB{9.262}{5.317}{2.221}{%
\RG{0.000}{-0.005}{10.9091}{mi}{120}\RG{6.235}{-1.641}{7.9701}{mi}{105}%
}\hskip.1em\allowbreak{} 相当于交换两行，\allowbreak{}\hskip.1em \RB{14.409}{7.946}{2.221}{%
\RG{0.000}{-0.005}{10.9091}{mi}{68}\RG{9.031}{-1.641}{7.9701}{mi}{110}%
}\hskip.1em\allowbreak{} 变号而\allowbreak{}\hskip.1em \RB{8.573}{8.186}{2.621}{%
\RG{0.000}{-0.005}{10.9091}{mi}{81}%
}\hskip.1em\allowbreak{} 不变；所以\allowbreak{}\hskip.1em \RB{8.669}{7.913}{0.505}{%
\RG{0.000}{-0.005}{10.9091}{mi}{70}%
}\hskip.1em\allowbreak{} 是关于\allowbreak{}\hskip.1em \RB{6.249}{5.317}{0.625}{%
\RG{0.000}{-0.005}{10.9091}{mi}{120}%
}\hskip.1em\allowbreak{} 的交错多项式。同理，它关于\allowbreak{}\hskip.1em \RB{5.846}{5.322}{2.736}{%
\RG{0.000}{0.000}{10.9091}{mi}{121}%
}\hskip.1em\allowbreak{} 也交错。

因此，\allowbreak{}\hskip.1em \RB{34.189}{5.318}{3.766}{%
\RG{0.000}{-0.004}{10.9091}{mi}{120}\RG{6.235}{-1.640}{7.9701}{mi}{105}\RG{12.647}{-0.004}{10.9091}{rm}{61}\RG{24.162}{-0.004}{10.9091}{mi}{120}\RG{30.397}{-1.640}{7.9701}{mi}{106}%
}\hskip.1em\allowbreak{} 或\allowbreak{}\hskip.1em \RB{32.416}{5.318}{3.766}{%
\RG{0.000}{-0.004}{10.9091}{mi}{121}\RG{5.349}{-1.640}{7.9701}{mi}{105}\RG{11.760}{-0.004}{10.9091}{rm}{61}\RG{23.276}{-0.004}{10.9091}{mi}{121}\RG{28.624}{-1.640}{7.9701}{mi}{106}%
}\hskip.1em\allowbreak{} 时\allowbreak{}\hskip.1em \RB{28.577}{7.914}{0.744}{%
\RG{0.000}{-0.004}{10.9091}{mi}{70}\RG{11.560}{-0.004}{10.9091}{rm}{61}\RG{23.058}{-0.004}{10.9091}{rm}{48}%
}\hskip.1em\allowbreak{} 。由多项式因式定理，各一次因子\allowbreak{}\hskip.1em \RB{32.933}{5.318}{3.766}{%
\RG{0.000}{-0.004}{10.9091}{mi}{120}\RG{6.234}{-1.640}{7.9701}{mi}{106}\RG{12.902}{-0.004}{10.9091}{sy}{0}\RG{23.672}{-0.004}{10.9091}{mi}{120}\RG{29.906}{-1.640}{7.9701}{mi}{105}%
}\hskip.1em\allowbreak{} 、\allowbreak{}\hskip.1em \RB{31.161}{5.318}{3.766}{%
\RG{0.000}{-0.004}{10.9091}{mi}{121}\RG{5.348}{-1.640}{7.9701}{mi}{106}\RG{12.016}{-0.004}{10.9091}{sy}{0}\RG{22.785}{-0.004}{10.9091}{mi}{121}\RG{28.134}{-1.640}{7.9701}{mi}{105}%
}\hskip.1em\allowbreak{} 均整除\allowbreak{}\hskip.1em \RB{8.669}{7.919}{0.500}{%
\RG{0.000}{0.001}{10.9091}{mi}{70}%
}\hskip.1em\allowbreak{} ；这些因子两两不相伴，故其乘积整除\allowbreak{}\hskip.1em \RB{8.669}{7.919}{0.500}{%
\RG{0.000}{0.001}{10.9091}{mi}{70}%
}\hskip.1em\allowbreak{} 。又

\[
\RB{209.649}{9.290}{5.682}{%
\RT{0.000}{-0.003}{10.9091}{d}\RT{6.066}{-0.003}{10.9091}{e}\RT{10.909}{-0.003}{10.9091}{g}\RG{26.609}{-0.003}{10.9091}{rm}{1}\RG{35.697}{-0.003}{10.9091}{rm}{40}\RG{39.942}{-0.003}{10.9091}{mi}{120}\RG{46.177}{-0.003}{10.9091}{rm}{41}\RG{53.432}{-0.003}{10.9091}{rm}{61}\RG{64.966}{-0.003}{10.9091}{mi}{110}\RG{73.938}{-0.003}{10.9091}{sy}{0}\RG{84.847}{-0.003}{10.9091}{rm}{49}\RG{90.302}{-0.003}{10.9091}{mi}{59}\RT{116.969}{-0.003}{10.9091}{d}\RT{123.034}{-0.003}{10.9091}{e}\RT{127.878}{-0.003}{10.9091}{g}\RG{143.737}{-0.003}{10.9091}{rm}{1}\RG{152.825}{-0.003}{10.9091}{rm}{40}\RG{157.070}{-0.003}{10.9091}{mi}{121}\RG{162.810}{-0.003}{10.9091}{rm}{41}\RG{170.065}{-0.003}{10.9091}{rm}{61}\RG{181.598}{-0.003}{10.9091}{mi}{110}\RG{190.571}{-0.003}{10.9091}{sy}{0}\RG{201.480}{-0.003}{10.9091}{rm}{49}\RG{206.935}{-0.003}{10.9091}{mi}{59}\RG{16.364}{-2.748}{7.9701}{mi}{120}\RG{133.332}{-2.748}{7.9701}{mi}{121}\RG{21.131}{-3.963}{5.9776}{mi}{105}\RG{137.518}{-3.963}{5.9776}{mi}{106}%
}
\]
因此商不依赖于任何\allowbreak{}\hskip.1em \RB{9.262}{5.325}{2.213}{%
\RG{0.000}{0.003}{10.9091}{mi}{120}\RG{6.235}{-1.633}{7.9701}{mi}{105}%
}\hskip.1em\allowbreak{} 或\allowbreak{}\hskip.1em \RB{9.141}{5.325}{3.759}{%
\RG{0.000}{0.003}{10.9091}{mi}{121}\RG{5.349}{-1.633}{7.9701}{mi}{106}%
}\hskip.1em\allowbreak{} 。于是存在常数\allowbreak{}\hskip.1em \RB{10.099}{5.325}{2.213}{%
\RG{0.000}{0.003}{10.9091}{mi}{99}\RG{4.721}{-1.633}{7.9701}{mi}{110}%
}\hskip.1em\allowbreak{} ，使

\[
\RB{83.154}{8.681}{3.228}{%
\RG{0.000}{-0.001}{10.9091}{mi}{70}\RG{11.561}{-0.001}{10.9091}{rm}{61}\RG{23.076}{-0.001}{10.9091}{mi}{99}\RG{33.433}{-0.001}{10.9091}{rm}{1}\RG{42.521}{-0.001}{10.9091}{rm}{40}\RG{46.766}{-0.001}{10.9091}{mi}{120}\RG{53.001}{-0.001}{10.9091}{rm}{41}\RG{57.234}{-0.001}{10.9091}{rm}{1}\RG{66.322}{-0.001}{10.9091}{rm}{40}\RG{70.578}{-0.001}{10.9091}{mi}{121}\RG{76.318}{-0.001}{10.9091}{rm}{41}\RG{80.560}{-0.001}{10.9091}{mi}{58}\RG{27.797}{-1.637}{7.9701}{mi}{110}%
}
\]
为确定常数，在形式幂级数意义下按各行展开

\[
\RB{96.260}{15.142}{15.212}{%
\RG{53.875}{10.359}{10.9092}{ex}{88}\RG{16.340}{7.376}{10.9091}{rm}{49}\RG{77.687}{4.500}{7.9701}{mi}{107}\RG{88.546}{4.500}{7.9701}{mi}{107}\RG{42.360}{-0.004}{10.9091}{rm}{61}\RG{71.452}{-0.004}{10.9091}{mi}{120}\RG{82.807}{-0.004}{10.9091}{mi}{121}\RG{93.666}{-0.004}{10.9091}{mi}{58}\RG{77.687}{-2.701}{7.9701}{mi}{105}\RG{88.155}{-2.701}{7.9701}{mi}{106}\RG{0.000}{-7.487}{10.9091}{rm}{49}\RG{7.878}{-7.487}{10.9091}{sy}{0}\RG{18.787}{-7.487}{10.9091}{mi}{120}\RG{28.404}{-7.487}{10.9091}{mi}{121}\RG{25.022}{-9.124}{7.9701}{mi}{105}\RG{33.752}{-9.124}{7.9701}{mi}{106}\RG{54.033}{-13.357}{7.9701}{mi}{107}\RG{58.655}{-13.357}{7.9701}{sy}{21}\RG{65.242}{-13.357}{7.9701}{rm}{48}\RR{0.000}{2.723}{38.134}{0.436}%
}
\]
由行列式的多重线性，\allowbreak{}\hskip.1em \RB{14.410}{7.954}{2.213}{%
\RG{0.000}{0.003}{10.9091}{mi}{68}\RG{9.032}{-1.633}{7.9701}{mi}{110}%
}\hskip.1em\allowbreak{} 中单项式\allowbreak{}\hskip.1em \RB{47.034}{10.899}{3.621}{%
\RG{0.000}{0.003}{10.9091}{mi}{120}\RG{6.235}{4.868}{7.9701}{mi}{107}\RG{6.235}{-3.121}{7.9701}{rm}{49}\RG{10.644}{3.761}{5.9776}{rm}{49}\RG{17.111}{0.003}{10.9091}{sy}{1}\RG{21.944}{0.003}{10.9091}{sy}{1}\RG{26.788}{0.003}{10.9091}{sy}{1}\RG{31.657}{0.003}{10.9091}{mi}{120}\RG{37.892}{3.962}{7.9701}{mi}{107}\RG{37.892}{-2.693}{7.9701}{mi}{110}\RG{42.301}{2.966}{5.9776}{mi}{110}%
}\hskip.1em\allowbreak{} 的系数为

\[
\RB{60.480}{11.003}{5.383}{%
\RG{25.135}{4.972}{7.9701}{mi}{107}\RG{37.447}{4.504}{7.9701}{mi}{110}\RG{29.545}{3.757}{5.9776}{mi}{105}\RT{0.000}{-0.001}{10.9091}{d}\RT{6.066}{-0.001}{10.9091}{e}\RT{10.909}{-0.001}{10.9091}{t}\RG{15.153}{0.000}{10.9091}{rm}{40}\RG{19.395}{-0.001}{10.9091}{mi}{121}\RG{33.205}{-0.001}{10.9091}{rm}{41}\RG{57.886}{-0.001}{10.9091}{mi}{58}\RG{37.447}{-2.697}{7.9701}{mi}{105}\RG{40.325}{-2.697}{7.9701}{mi}{59}\RG{42.676}{-2.697}{7.9701}{mi}{106}\RG{46.566}{-2.697}{7.9701}{rm}{61}\RG{53.150}{-2.697}{7.9701}{rm}{49}\RG{24.744}{-3.257}{7.9701}{mi}{106}%
}
\]
若两个\allowbreak{}\hskip.1em \RB{8.707}{8.067}{2.221}{%
\RG{0.000}{-0.004}{10.9091}{mi}{107}\RG{5.680}{-1.641}{7.9701}{mi}{105}%
}\hskip.1em\allowbreak{} 相等，该系数为零；所以\allowbreak{}\hskip.1em \RB{14.410}{7.947}{2.221}{%
\RG{0.000}{-0.004}{10.9091}{mi}{68}\RG{9.032}{-1.641}{7.9701}{mi}{110}%
}\hskip.1em\allowbreak{} 关于\allowbreak{}\hskip.1em \RB{6.249}{5.318}{0.624}{%
\RG{0.000}{-0.004}{10.9091}{mi}{120}%
}\hskip.1em\allowbreak{} 的最低可能总次数为\allowbreak{}\hskip.1em \RB{118.082}{8.678}{3.231}{%
\RG{0.000}{-0.004}{10.9091}{mi}{109}\RG{12.609}{-0.004}{10.9091}{rm}{61}\RG{24.107}{-0.004}{10.9091}{rm}{48}\RG{31.624}{-0.004}{10.9091}{rm}{43}\RG{42.162}{-0.004}{10.9091}{rm}{49}\RG{49.679}{-0.004}{10.9091}{rm}{43}\RG{60.245}{-0.004}{10.9091}{sy}{1}\RG{65.078}{-0.004}{10.9091}{sy}{1}\RG{69.922}{-0.004}{10.9091}{sy}{1}\RG{75.033}{-0.004}{10.9091}{rm}{43}\RG{85.560}{-0.004}{10.9091}{rm}{40}\RG{89.821}{-0.004}{10.9091}{mi}{110}\RG{98.430}{-0.004}{10.9091}{sy}{0}\RG{108.975}{-0.004}{10.9091}{rm}{49}\RG{114.430}{-0.004}{10.9091}{rm}{41}%
}\hskip.1em\allowbreak{} 。又\allowbreak{}\hskip.1em \RB{8.573}{8.192}{2.615}{%
\RG{0.000}{0.001}{10.9091}{mi}{81}%
}\hskip.1em\allowbreak{} 关于\allowbreak{}\hskip.1em \RB{6.249}{5.323}{0.619}{%
\RG{0.000}{0.001}{10.9091}{mi}{120}%
}\hskip.1em\allowbreak{} 的常数项为\allowbreak{}\hskip.1em \RB{5.071}{7.767}{0.500}{%
\RG{0.000}{0.001}{10.9091}{rm}{49}%
}\hskip.1em\allowbreak{} ，故\allowbreak{}\hskip.1em \RB{46.146}{8.192}{2.615}{%
\RG{0.000}{0.001}{10.9091}{mi}{70}\RG{11.579}{0.001}{10.9091}{rm}{61}\RG{23.112}{0.001}{10.9091}{mi}{81}\RG{31.730}{0.001}{10.9091}{mi}{68}\RG{40.768}{-1.636}{7.9701}{mi}{110}%
}\hskip.1em\allowbreak{} 中总次数\allowbreak{}\hskip.1em \RB{9.751}{5.323}{0.619}{%
\RG{0.000}{0.001}{10.9091}{mi}{109}%
}\hskip.1em\allowbreak{} 的部分与\allowbreak{}\hskip.1em \RB{14.410}{7.952}{2.216}{%
\RG{0.000}{0.001}{10.9091}{mi}{68}\RG{9.032}{-1.636}{7.9701}{mi}{110}%
}\hskip.1em\allowbreak{} 相同。特别地，\allowbreak{}\hskip.1em \RB{60.321}{9.759}{3.422}{%
\RG{0.000}{0.001}{10.9091}{mi}{120}\RG{6.235}{3.959}{7.9701}{rm}{48}\RG{6.235}{-2.922}{7.9701}{rm}{49}\RG{10.967}{0.001}{10.9091}{mi}{120}\RG{17.202}{3.959}{7.9701}{rm}{49}\RG{17.202}{-2.922}{7.9701}{rm}{50}\RG{23.753}{0.001}{10.9091}{sy}{1}\RG{28.586}{0.001}{10.9091}{sy}{1}\RG{33.430}{0.001}{10.9091}{sy}{1}\RG{38.298}{0.001}{10.9091}{mi}{120}\RG{44.533}{3.959}{7.9701}{mi}{110}\RG{44.533}{-2.696}{7.9701}{mi}{110}\RG{49.671}{3.959}{7.9701}{sy}{0}\RG{56.258}{3.959}{7.9701}{rm}{49}%
}\hskip.1em\allowbreak{} 的系数为

\[
\RB{104.474}{10.388}{5.382}{%
\RG{25.135}{4.588}{7.9701}{mi}{105}\RG{28.018}{4.588}{7.9701}{sy}{0}\RG{34.605}{4.588}{7.9701}{rm}{49}\RG{43.580}{4.505}{7.9701}{mi}{110}\RT{0.000}{0.000}{10.9091}{d}\RT{6.066}{0.000}{10.9091}{e}\RT{10.909}{0.000}{10.9091}{t}\RG{15.153}{0.000}{10.9091}{rm}{40}\RG{19.395}{0.000}{10.9091}{mi}{121}\RG{39.337}{0.000}{10.9091}{rm}{41}\RG{67.049}{0.000}{10.9091}{rm}{61}\RG{78.547}{0.000}{10.9091}{rm}{1}\RG{87.634}{0.000}{10.9091}{rm}{40}\RG{91.898}{0.000}{10.9091}{mi}{121}\RG{97.638}{0.000}{10.9091}{rm}{41}\RG{101.880}{0.000}{10.9091}{mi}{58}\RG{43.580}{-2.696}{7.9701}{mi}{105}\RG{46.458}{-2.696}{7.9701}{mi}{59}\RG{48.809}{-2.696}{7.9701}{mi}{106}\RG{52.699}{-2.696}{7.9701}{rm}{61}\RG{59.283}{-2.696}{7.9701}{rm}{49}\RG{24.744}{-3.256}{7.9701}{mi}{106}%
}
\]
最后一步是Vandermonde行列式公式：其左端交错、次数为\allowbreak{}\hskip.1em \RB{9.751}{5.322}{0.620}{%
\RG{0.000}{0.000}{10.9091}{mi}{109}%
}\hskip.1em\allowbreak{} ，而\allowbreak{}\hskip.1em \RB{58.835}{9.759}{3.423}{%
\RG{0.000}{0.000}{10.9091}{mi}{121}\RG{5.348}{-2.923}{7.9701}{rm}{49}\RG{5.739}{3.959}{7.9701}{rm}{48}\RG{10.472}{-0.000}{10.9091}{mi}{121}\RG{15.820}{-2.923}{7.9701}{rm}{50}\RG{16.212}{3.959}{7.9701}{rm}{49}\RG{22.762}{-0.000}{10.9091}{sy}{1}\RG{27.595}{-0.000}{10.9091}{sy}{1}\RG{32.439}{-0.000}{10.9091}{sy}{1}\RG{37.308}{-0.000}{10.9091}{mi}{121}\RG{42.656}{-2.697}{7.9701}{mi}{110}\RG{43.047}{3.959}{7.9701}{mi}{110}\RG{48.186}{3.959}{7.9701}{sy}{0}\RG{54.772}{3.959}{7.9701}{rm}{49}%
}\hskip.1em\allowbreak{} 的系数为\allowbreak{}\hskip.1em \RB{5.071}{7.765}{0.500}{%
\RG{0.000}{-0.000}{10.9091}{rm}{49}%
}\hskip.1em\allowbreak{} ，故等于\allowbreak{}\hskip.1em \RB{22.726}{8.687}{3.222}{%
\RG{0.000}{0.005}{10.9091}{rm}{1}\RG{9.087}{0.005}{10.9091}{rm}{40}\RG{13.333}{0.005}{10.9091}{mi}{121}\RG{19.073}{0.005}{10.9091}{rm}{41}%
}\hskip.1em\allowbreak{} 。同样，该单项式关于\allowbreak{}\hskip.1em \RB{6.249}{5.327}{0.615}{%
\RG{0.000}{0.005}{10.9091}{mi}{120}%
}\hskip.1em\allowbreak{} 在\allowbreak{}\hskip.1em \RB{23.221}{8.687}{3.222}{%
\RG{0.000}{0.005}{10.9091}{rm}{1}\RG{9.087}{0.005}{10.9091}{rm}{40}\RG{13.333}{0.005}{10.9091}{mi}{120}\RG{19.568}{0.005}{10.9091}{rm}{41}%
}\hskip.1em\allowbreak{} 中的系数也是\allowbreak{}\hskip.1em \RB{5.071}{7.770}{0.500}{%
\RG{0.000}{0.005}{10.9091}{rm}{49}%
}\hskip.1em\allowbreak{} 。比较\allowbreak{}\hskip.1em \RB{80.227}{8.687}{3.222}{%
\RG{0.000}{0.005}{10.9091}{mi}{70}\RG{11.689}{0.005}{10.9091}{rm}{61}\RG{23.332}{0.005}{10.9091}{mi}{99}\RG{28.053}{-1.632}{7.9701}{mi}{110}\RG{33.690}{0.005}{10.9091}{rm}{1}\RG{42.778}{0.005}{10.9091}{rm}{40}\RG{47.023}{0.005}{10.9091}{mi}{120}\RG{53.258}{0.005}{10.9091}{rm}{41}\RG{57.491}{0.005}{10.9091}{rm}{1}\RG{66.579}{0.005}{10.9091}{rm}{40}\RG{70.835}{0.005}{10.9091}{mi}{121}\RG{76.575}{0.005}{10.9091}{rm}{41}%
}\hskip.1em\allowbreak{} 两侧系数得\allowbreak{}\hskip.1em \RB{29.956}{7.765}{2.217}{%
\RG{0.000}{-0.000}{10.9091}{mi}{99}\RG{4.721}{-1.637}{7.9701}{mi}{110}\RG{13.387}{-0.000}{10.9091}{rm}{61}\RG{24.885}{-0.000}{10.9091}{rm}{49}%
}\hskip.1em\allowbreak{} ，于是得到所求公式。\allowbreak{}\hskip.1em \RB{8.366}{7.995}{0.500}{%
\RG{0.000}{-0.000}{10.9091}{am}{3}%
}\hskip.1em\allowbreak{}

方法小结：先清除分母，再用交错性与次数确定全部因子；一个最低次数项即可确定常数。重复的\allowbreak{}\hskip.1em \RB{9.262}{5.325}{2.213}{%
\RG{0.000}{0.003}{10.9091}{mi}{120}\RG{6.235}{-1.633}{7.9701}{mi}{105}%
}\hskip.1em\allowbreak{} 或\allowbreak{}\hskip.1em \RB{9.141}{5.325}{3.759}{%
\RG{0.000}{0.003}{10.9091}{mi}{121}\RG{5.349}{-1.633}{7.9701}{mi}{106}%
}\hskip.1em\allowbreak{} 自动给出行列式为零。题目来源：谢启鸿《复旦大学高等代数习题集（2023版）》，第123页。

% Original page 5.
\section*{2022A02}
设\allowbreak{}\hskip.1em \RB{12.067}{8.191}{2.217}{%
\RG{0.000}{0.000}{10.9091}{mi}{83}\RG{6.689}{-1.637}{7.9701}{mi}{110}%
}\hskip.1em\allowbreak{} 为\allowbreak{}\hskip.1em \RB{54.443}{8.682}{3.227}{%
\RG{0.000}{0.000}{10.9091}{rm}{40}\RG{4.233}{0.000}{10.9091}{rm}{49}\RG{9.697}{0.000}{10.9091}{mi}{59}\RG{14.546}{0.000}{10.9091}{rm}{50}\RG{20.000}{0.000}{10.9091}{mi}{59}\RG{24.833}{0.000}{10.9091}{mi}{58}\RG{29.677}{0.000}{10.9091}{mi}{58}\RG{34.509}{0.000}{10.9091}{mi}{58}\RG{39.353}{0.000}{10.9091}{mi}{59}\RG{44.197}{0.000}{10.9091}{mi}{110}\RG{50.790}{0.000}{10.9091}{rm}{41}%
}\hskip.1em\allowbreak{} 的全排列集合，\allowbreak{}\hskip.1em \RB{63.820}{8.682}{3.227}{%
\RG{0.000}{0.000}{10.9091}{mi}{78}\RG{9.955}{0.000}{10.9091}{rm}{40}\RG{14.197}{0.000}{10.9091}{mi}{107}\RG{19.876}{-1.637}{7.9701}{rm}{49}\RG{24.609}{0.000}{10.9091}{mi}{59}\RG{29.442}{0.000}{10.9091}{mi}{58}\RG{34.286}{0.000}{10.9091}{mi}{58}\RG{39.118}{0.000}{10.9091}{mi}{58}\RG{43.962}{0.000}{10.9091}{mi}{59}\RG{48.806}{0.000}{10.9091}{mi}{107}\RG{54.530}{-1.637}{7.9701}{mi}{110}\RG{60.167}{0.000}{10.9091}{rm}{41}%
}\hskip.1em\allowbreak{} 表示排列的逆序数。

(1)若\allowbreak{}\hskip.1em \RB{42.814}{8.068}{2.968}{%
\RG{0.000}{-0.003}{10.9091}{mi}{107}\RG{5.679}{-1.639}{7.9701}{mi}{105}\RG{12.091}{-0.003}{10.9091}{mi}{62}\RG{23.589}{-0.003}{10.9091}{mi}{107}\RG{29.285}{-1.639}{7.9701}{mi}{105}\RG{32.168}{-1.639}{7.9701}{rm}{43}\RG{38.752}{-1.639}{7.9701}{rm}{49}%
}\hskip.1em\allowbreak{} ，证明交换这两个相邻元素后，逆序数减少\allowbreak{}\hskip.1em \RB{5.071}{7.763}{0.503}{%
\RG{0.000}{-0.003}{10.9091}{rm}{49}%
}\hskip.1em\allowbreak{} 。(2)利用(1)并对逆序数归纳，证明一个排列可经恰好等于其逆序数的相邻对换化为常序排列

\[
\RB{66.443}{9.231}{3.223}{%
\RG{0.000}{0.004}{10.9091}{rm}{40}\RG{4.233}{0.004}{10.9091}{rm}{49}\RG{9.697}{0.004}{10.9091}{mi}{59}\RG{14.546}{0.004}{10.9091}{rm}{50}\RG{20.000}{0.004}{10.9091}{mi}{59}\RG{24.833}{0.004}{10.9091}{mi}{58}\RG{29.677}{0.004}{10.9091}{mi}{58}\RG{34.520}{0.004}{10.9091}{mi}{58}\RG{39.353}{0.004}{10.9091}{mi}{59}\RG{44.197}{0.004}{10.9091}{mi}{110}\RG{50.790}{0.004}{10.9091}{rm}{41}\RT{55.033}{0.004}{10.9091}{。}%
}
\]
(3)利用(1)证明：对每个整数\allowbreak{}\hskip.1em \RB{91.781}{8.679}{3.230}{%
\RG{0.000}{-0.003}{10.9091}{rm}{49}\RG{8.485}{-0.003}{10.9091}{sy}{20}\RG{20.000}{-0.003}{10.9091}{mi}{107}\RG{29.053}{-0.003}{10.9091}{sy}{20}\RG{40.568}{-0.003}{10.9091}{mi}{110}\RG{47.116}{-0.003}{10.9091}{rm}{40}\RG{51.358}{-0.003}{10.9091}{mi}{110}\RG{60.332}{-0.003}{10.9091}{sy}{0}\RG{71.241}{-0.003}{10.9091}{rm}{49}\RG{76.695}{-0.003}{10.9091}{rm}{41}\RG{80.928}{-0.003}{10.9091}{rm}{47}\RG{86.383}{-0.003}{10.9091}{rm}{50}%
}\hskip.1em\allowbreak{} ，均存在逆序数恰为\allowbreak{}\hskip.1em \RB{6.042}{8.068}{0.623}{%
\RG{0.000}{-0.003}{10.9091}{mi}{107}%
}\hskip.1em\allowbreak{} 的排列。(4)设逆序数为\allowbreak{}\hskip.1em \RB{6.042}{8.070}{0.621}{%
\RG{0.000}{-0.001}{10.9091}{mi}{107}%
}\hskip.1em\allowbreak{} 的排列有\allowbreak{}\hskip.1em \RB{35.269}{8.681}{3.228}{%
\RG{0.000}{-0.001}{10.9091}{mi}{78}\RG{9.954}{-0.001}{10.9091}{rm}{40}\RG{14.197}{-0.001}{10.9091}{mi}{110}\RG{20.743}{-0.001}{10.9091}{mi}{59}\RG{25.575}{-0.001}{10.9091}{mi}{107}\RG{31.616}{-0.001}{10.9091}{rm}{41}%
}\hskip.1em\allowbreak{} 个，证明

\[
\RB{150.834}{33.697}{36.164}{%
\RG{0.000}{25.015}{10.9091}{mi}{78}\RG{9.955}{25.015}{10.9091}{rm}{40}\RG{14.197}{25.015}{10.9091}{mi}{110}\RG{20.743}{25.015}{10.9091}{mi}{59}\RG{25.594}{25.015}{10.9091}{rm}{49}\RG{31.049}{25.015}{10.9091}{rm}{41}\RG{38.303}{25.015}{10.9091}{rm}{61}\RG{49.836}{25.015}{10.9091}{mi}{110}\RG{58.808}{25.015}{10.9091}{sy}{0}\RG{69.718}{25.015}{10.9091}{rm}{49}\RG{75.173}{25.015}{10.9091}{mi}{59}\RG{57.580}{11.341}{7.9701}{rm}{50}\RG{51.031}{7.383}{10.9091}{mi}{110}\RG{64.737}{7.383}{10.9091}{sy}{0}\RG{75.646}{7.383}{10.9091}{mi}{110}\RG{84.618}{7.383}{10.9091}{sy}{0}\RG{95.527}{7.383}{10.9091}{rm}{50}\RG{0.000}{0.003}{10.9091}{mi}{78}\RG{9.955}{0.003}{10.9091}{rm}{40}\RG{14.197}{0.003}{10.9091}{mi}{110}\RG{20.743}{0.003}{10.9091}{mi}{59}\RG{25.594}{0.003}{10.9091}{rm}{50}\RG{31.049}{0.003}{10.9091}{rm}{41}\RG{38.303}{0.003}{10.9091}{rm}{61}\RG{113.086}{0.003}{10.9091}{rm}{40}\RG{117.329}{0.003}{10.9091}{mi}{110}\RG{126.907}{0.003}{10.9091}{sy}{21}\RG{138.422}{0.003}{10.9091}{rm}{50}\RG{143.877}{0.003}{10.9091}{rm}{41}\RG{148.119}{0.003}{10.9091}{mi}{59}\RG{73.280}{-7.481}{10.9091}{rm}{50}\RG{57.580}{-16.602}{7.9701}{rm}{51}\RG{51.031}{-20.561}{10.9091}{mi}{110}\RG{64.737}{-20.561}{10.9091}{sy}{0}\RG{75.646}{-20.561}{10.9091}{rm}{55}\RG{81.101}{-20.561}{10.9091}{mi}{110}\RG{0.000}{-27.941}{10.9091}{mi}{78}\RG{9.955}{-27.941}{10.9091}{rm}{40}\RG{14.197}{-27.941}{10.9091}{mi}{110}\RG{20.743}{-27.941}{10.9091}{mi}{59}\RG{25.594}{-27.941}{10.9091}{rm}{51}\RG{31.049}{-27.941}{10.9091}{rm}{41}\RG{38.303}{-27.941}{10.9091}{rm}{61}\RG{99.753}{-27.941}{10.9091}{rm}{40}\RG{103.996}{-27.941}{10.9091}{mi}{110}\RG{113.574}{-27.941}{10.9091}{sy}{21}\RG{125.089}{-27.941}{10.9091}{rm}{51}\RG{130.544}{-27.941}{10.9091}{rm}{41}\RG{134.786}{-27.941}{10.9091}{mi}{58}\RG{66.613}{-35.424}{10.9091}{rm}{54}\RR{51.031}{2.730}{49.950}{0.436}\RR{51.031}{-25.214}{36.617}{0.436}%
}
\]
证明：(1)令\allowbreak{}\hskip.1em \RB{86.759}{8.070}{2.966}{%
\RG{0.000}{-0.001}{10.9091}{mi}{97}\RG{9.530}{-0.001}{10.9091}{rm}{61}\RG{21.780}{-0.001}{10.9091}{mi}{107}\RG{27.459}{-1.637}{7.9701}{mi}{105}\RG{34.605}{-0.001}{10.9091}{mi}{62}\RG{46.845}{-0.001}{10.9091}{mi}{98}\RG{55.300}{-0.001}{10.9091}{rm}{61}\RG{67.550}{-0.001}{10.9091}{mi}{107}\RG{73.229}{-1.637}{7.9701}{mi}{105}\RG{76.113}{-1.637}{7.9701}{rm}{43}\RG{82.696}{-1.637}{7.9701}{rm}{49}%
}\hskip.1em\allowbreak{} 。对位于这两个位置之前的元素\allowbreak{}\hskip.1em \RB{5.191}{5.321}{0.621}{%
\RG{0.000}{-0.001}{10.9091}{mi}{99}%
}\hskip.1em\allowbreak{} ，它与\allowbreak{}\hskip.1em \RB{15.620}{8.070}{2.606}{%
\RG{0.000}{-0.001}{10.9091}{mi}{97}\RG{5.760}{-0.001}{10.9091}{mi}{59}\RG{10.593}{-0.001}{10.9091}{mi}{98}%
}\hskip.1em\allowbreak{} 所贡献的逆序数为1\allowbreak{}\hskip.1em \RB{26.040}{6.864}{2.642}{%
\RG{0.000}{-1.632}{7.9701}{mi}{99}\RG{3.666}{-1.632}{7.9701}{mi}{62}\RG{10.250}{-1.632}{7.9701}{mi}{97}\RG{17.674}{0.004}{10.9091}{rm}{43}%
}\hskip.1em\allowbreak{} 1\allowbreak{}\hskip.1em \RB{14.256}{4.258}{2.783}{%
\RG{0.000}{-1.773}{7.9701}{mi}{99}\RG{3.666}{-1.773}{7.9701}{mi}{62}\RG{10.250}{-1.773}{7.9701}{mi}{98}%
}\hskip.1em\allowbreak{} ；交换后这一和不变。对位于它们之后的元素\allowbreak{}\hskip.1em \RB{5.191}{5.326}{0.616}{%
\RG{0.000}{0.004}{10.9091}{mi}{99}%
}\hskip.1em\allowbreak{} ，相应贡献为1\allowbreak{}\hskip.1em \RB{26.040}{6.864}{2.642}{%
\RG{0.000}{-1.632}{7.9701}{mi}{97}\RG{4.495}{-1.632}{7.9701}{mi}{62}\RG{11.079}{-1.632}{7.9701}{mi}{99}\RG{17.674}{0.004}{10.9091}{rm}{43}%
}\hskip.1em\allowbreak{} 1\allowbreak{}\hskip.1em \RB{14.376}{4.258}{2.783}{%
\RG{0.000}{-1.773}{7.9701}{mi}{98}\RG{3.618}{-1.773}{7.9701}{mi}{62}\RG{10.202}{-1.773}{7.9701}{mi}{99}%
}\hskip.1em\allowbreak{} ，也不变。

其余不涉及\allowbreak{}\hskip.1em \RB{15.620}{8.071}{2.605}{%
\RG{0.000}{-0.000}{10.9091}{mi}{97}\RG{5.760}{-0.000}{10.9091}{mi}{59}\RG{10.593}{-0.000}{10.9091}{mi}{98}%
}\hskip.1em\allowbreak{} 的元素对完全不变，而\allowbreak{}\hskip.1em \RB{23.192}{8.682}{3.227}{%
\RG{0.000}{-0.000}{10.9091}{rm}{40}\RG{4.243}{-0.000}{10.9091}{mi}{97}\RG{10.003}{-0.000}{10.9091}{mi}{59}\RG{14.836}{-0.000}{10.9091}{mi}{98}\RG{19.539}{-0.000}{10.9091}{rm}{41}%
}\hskip.1em\allowbreak{} 本身由逆序变为顺序。故逆序数恰好减少\allowbreak{}\hskip.1em \RB{5.071}{7.766}{0.500}{%
\RG{0.000}{-0.000}{10.9091}{rm}{49}%
}\hskip.1em\allowbreak{} 。(2)对逆序数\allowbreak{}\hskip.1em \RB{9.751}{5.318}{0.624}{%
\RG{0.000}{-0.004}{10.9091}{mi}{109}%
}\hskip.1em\allowbreak{} 归纳。若\allowbreak{}\hskip.1em \RB{29.710}{7.762}{0.744}{%
\RG{0.000}{-0.004}{10.9091}{mi}{109}\RG{12.650}{-0.004}{10.9091}{rm}{61}\RG{24.192}{-0.004}{10.9091}{rm}{48}%
}\hskip.1em\allowbreak{} ，排列中没有逆序，只能是常序排列，无须交换。若\allowbreak{}\hskip.1em \RB{29.724}{7.762}{0.940}{%
\RG{0.000}{-0.004}{10.9091}{mi}{109}\RG{12.644}{-0.004}{10.9091}{mi}{62}\RG{24.206}{-0.004}{10.9091}{rm}{48}%
}\hskip.1em\allowbreak{} ，排列必存在相邻逆序：否则\allowbreak{}\hskip.1em \RB{89.050}{8.072}{2.215}{%
\RG{0.000}{0.001}{10.9091}{mi}{107}\RG{5.679}{-1.635}{7.9701}{rm}{49}\RG{13.576}{0.001}{10.9091}{mi}{60}\RG{25.216}{0.001}{10.9091}{mi}{107}\RG{30.905}{-1.635}{7.9701}{rm}{50}\RG{38.802}{0.001}{10.9091}{mi}{60}\RG{50.452}{0.001}{10.9091}{sy}{1}\RG{55.285}{0.001}{10.9091}{sy}{1}\RG{60.129}{0.001}{10.9091}{sy}{1}\RG{66.344}{0.001}{10.9091}{mi}{60}\RG{77.984}{0.001}{10.9091}{mi}{107}\RG{83.673}{-1.635}{7.9701}{mi}{110}%
}\hskip.1em\allowbreak{} ，逆序数便为零。交换一组相邻逆序，由(1)得到逆序数为\allowbreak{}\hskip.1em \RB{27.983}{7.762}{0.624}{%
\RG{0.000}{-0.004}{10.9091}{mi}{109}\RG{12.002}{-0.004}{10.9091}{sy}{0}\RG{22.912}{-0.004}{10.9091}{rm}{49}%
}\hskip.1em\allowbreak{} 的排列；由归纳假设，再经\allowbreak{}\hskip.1em \RB{27.982}{7.762}{0.624}{%
\RG{0.000}{-0.004}{10.9091}{mi}{109}\RG{12.002}{-0.004}{10.9091}{sy}{0}\RG{22.911}{-0.004}{10.9091}{rm}{49}%
}\hskip.1em\allowbreak{} 次相邻对换可化为常序排列。总计恰为\allowbreak{}\hskip.1em \RB{9.751}{5.318}{0.624}{%
\RG{0.000}{-0.004}{10.9091}{mi}{109}%
}\hskip.1em\allowbreak{} 次。

(3)从逆序排列\allowbreak{}\hskip.1em \RB{74.340}{8.685}{3.224}{%
\RG{0.000}{0.003}{10.9091}{rm}{40}\RG{4.243}{0.003}{10.9091}{mi}{110}\RG{10.789}{0.003}{10.9091}{mi}{59}\RG{15.621}{0.003}{10.9091}{mi}{110}\RG{24.619}{0.003}{10.9091}{sy}{0}\RG{35.536}{0.003}{10.9091}{rm}{49}\RG{40.990}{0.003}{10.9091}{mi}{59}\RG{45.823}{0.003}{10.9091}{mi}{58}\RG{50.667}{0.003}{10.9091}{mi}{58}\RG{55.500}{0.003}{10.9091}{mi}{58}\RG{60.343}{0.003}{10.9091}{mi}{59}\RG{65.233}{0.003}{10.9091}{rm}{49}\RG{70.687}{0.003}{10.9091}{rm}{41}%
}\hskip.1em\allowbreak{} 出发，初始逆序数为\allowbreak{}\hskip.1em \RB{77.583}{8.685}{3.224}{%
\RG{0.000}{0.003}{10.9091}{mi}{77}\RG{14.821}{0.003}{10.9091}{rm}{61}\RG{26.355}{0.003}{10.9091}{mi}{110}\RG{32.903}{0.003}{10.9091}{rm}{40}\RG{37.146}{0.003}{10.9091}{mi}{110}\RG{46.125}{0.003}{10.9091}{sy}{0}\RG{57.042}{0.003}{10.9091}{rm}{49}\RG{62.497}{0.003}{10.9091}{rm}{41}\RG{66.730}{0.003}{10.9091}{rm}{47}\RG{72.184}{0.003}{10.9091}{rm}{50}%
}\hskip.1em\allowbreak{} 。反复交换一组相邻逆序，直至常序排列；由(1)、(2)，沿途逆序数依次为

\[
\RB{89.170}{7.955}{2.601}{%
\RG{0.000}{0.004}{10.9091}{mi}{77}\RG{11.160}{0.004}{10.9091}{mi}{59}\RG{16.004}{0.004}{10.9091}{mi}{77}\RG{30.212}{0.004}{10.9091}{sy}{0}\RG{41.121}{0.004}{10.9091}{rm}{49}\RG{46.576}{0.004}{10.9091}{mi}{59}\RG{51.409}{0.004}{10.9091}{mi}{58}\RG{56.253}{0.004}{10.9091}{mi}{58}\RG{61.086}{0.004}{10.9091}{mi}{58}\RG{65.929}{0.004}{10.9091}{mi}{59}\RG{70.819}{0.004}{10.9091}{rm}{49}\RG{76.273}{0.004}{10.9091}{mi}{59}\RG{81.122}{0.004}{10.9091}{rm}{48}\RG{86.576}{0.004}{10.9091}{mi}{58}%
}
\]
所以区间中的每个整数都能实现。

(4)为避免重复计数，对每个数值\allowbreak{}\hskip.1em \RB{4.831}{7.707}{2.740}{%
\RG{0.000}{-0.004}{10.9091}{mi}{106}%
}\hskip.1em\allowbreak{} 定义

\allowbreak{}\hskip.1em \RB{20.858}{5.324}{3.760}{%
\RG{0.000}{0.002}{10.9091}{mi}{101}\RG{5.080}{-1.634}{7.9701}{mi}{106}\RG{12.492}{0.002}{10.9091}{rm}{61}%
}\hskip.1em\allowbreak{} \#\allowbreak{}\hskip.1em \RB{5.158}{8.684}{3.225}{%
\RG{0.000}{0.002}{10.9091}{sy}{102}%
}\hskip.1em\allowbreak{} 位于\allowbreak{}\hskip.1em \RB{4.831}{7.713}{2.734}{%
\RG{0.000}{0.002}{10.9091}{mi}{106}%
}\hskip.1em\allowbreak{} 右边且小于\allowbreak{}\hskip.1em \RB{4.831}{7.713}{2.734}{%
\RG{0.000}{0.002}{10.9091}{mi}{106}%
}\hskip.1em\allowbreak{} 的元素\allowbreak{}\hskip.1em \RB{8.049}{8.684}{3.225}{%
\RG{0.000}{0.002}{10.9091}{sy}{103}\RG{5.454}{0.002}{10.9091}{mi}{58}%
}\hskip.1em\allowbreak{}

则\allowbreak{}\hskip.1em \RB{67.528}{7.764}{3.764}{%
\RG{0.000}{-0.002}{10.9091}{rm}{48}\RG{8.485}{-0.002}{10.9091}{sy}{20}\RG{20.000}{-0.002}{10.9091}{mi}{101}\RG{25.079}{-1.638}{7.9701}{mi}{106}\RG{32.492}{-0.002}{10.9091}{sy}{20}\RG{44.007}{-0.002}{10.9091}{mi}{106}\RG{51.548}{-0.002}{10.9091}{sy}{0}\RG{62.458}{-0.002}{10.9091}{rm}{49}%
}\hskip.1em\allowbreak{} ，并且每个逆序都由其较大元素唯一计入，故

\[
\RB{164.617}{8.686}{3.223}{%
\RG{0.000}{0.004}{10.9091}{mi}{78}\RG{9.955}{0.004}{10.9091}{rm}{40}\RG{14.197}{0.004}{10.9091}{mi}{107}\RG{24.609}{0.004}{10.9091}{mi}{59}\RG{29.442}{0.004}{10.9091}{mi}{58}\RG{34.286}{0.004}{10.9091}{mi}{58}\RG{39.118}{0.004}{10.9091}{mi}{58}\RG{43.962}{0.004}{10.9091}{mi}{59}\RG{48.806}{0.004}{10.9091}{mi}{107}\RG{60.168}{0.004}{10.9091}{rm}{41}\RG{67.422}{0.004}{10.9091}{rm}{61}\RG{78.955}{0.004}{10.9091}{mi}{101}\RG{91.192}{0.004}{10.9091}{rm}{43}\RG{102.101}{0.004}{10.9091}{mi}{101}\RG{114.337}{0.004}{10.9091}{rm}{43}\RG{125.246}{0.004}{10.9091}{sy}{1}\RG{130.079}{0.004}{10.9091}{sy}{1}\RG{134.922}{0.004}{10.9091}{sy}{1}\RG{140.397}{0.004}{10.9091}{rm}{43}\RG{151.306}{0.004}{10.9091}{mi}{101}\RG{162.022}{0.004}{10.9091}{mi}{58}\RG{19.876}{-1.632}{7.9701}{rm}{49}\RG{54.530}{-1.632}{7.9701}{mi}{110}\RG{84.035}{-1.632}{7.9701}{rm}{49}\RG{107.180}{-1.632}{7.9701}{rm}{50}\RG{156.386}{-1.632}{7.9701}{mi}{110}%
}
\]
这些整数独立决定排列：从单元素排列\allowbreak{}\hskip.1em \RB{13.340}{8.682}{3.227}{%
\RG{0.000}{0.000}{10.9091}{rm}{40}\RG{4.233}{0.000}{10.9091}{rm}{49}\RG{9.687}{0.000}{10.9091}{rm}{41}%
}\hskip.1em\allowbreak{} 开始，依次插入\allowbreak{}\hskip.1em \RB{56.776}{7.766}{2.736}{%
\RG{0.000}{0.000}{10.9091}{mi}{106}\RG{8.514}{0.000}{10.9091}{rm}{61}\RG{20.383}{0.000}{10.9091}{rm}{50}\RG{25.850}{0.000}{10.9091}{mi}{59}\RG{30.683}{0.000}{10.9091}{mi}{58}\RG{35.527}{0.000}{10.9091}{mi}{58}\RG{40.359}{0.000}{10.9091}{mi}{58}\RG{45.203}{0.000}{10.9091}{mi}{59}\RG{50.047}{0.000}{10.9091}{mi}{110}%
}\hskip.1em\allowbreak{} ，把\allowbreak{}\hskip.1em \RB{4.831}{7.711}{2.736}{%
\RG{0.000}{0.000}{10.9091}{mi}{106}%
}\hskip.1em\allowbreak{} 插在恰有\allowbreak{}\hskip.1em \RB{8.871}{5.322}{3.762}{%
\RG{0.000}{0.000}{10.9091}{mi}{101}\RG{5.079}{-1.636}{7.9701}{mi}{106}%
}\hskip.1em\allowbreak{} 个已有元素位于其右侧的位置。该位置唯一，且插入更大的元素不会改变此前的\allowbreak{}\hskip.1em \RB{8.107}{5.317}{2.221}{%
\RG{0.000}{-0.005}{10.9091}{mi}{101}\RG{5.080}{-1.641}{7.9701}{mi}{105}%
}\hskip.1em\allowbreak{} 。这就建立了排列与满足\allowbreak{}\hskip.1em \RB{67.528}{7.766}{3.762}{%
\RG{0.000}{0.000}{10.9091}{rm}{48}\RG{8.485}{0.000}{10.9091}{sy}{20}\RG{20.000}{0.000}{10.9091}{mi}{101}\RG{25.079}{-1.636}{7.9701}{mi}{106}\RG{32.492}{0.000}{10.9091}{sy}{20}\RG{44.007}{0.000}{10.9091}{mi}{106}\RG{51.548}{0.000}{10.9091}{sy}{0}\RG{62.458}{0.000}{10.9091}{rm}{49}%
}\hskip.1em\allowbreak{} 的整数列之间的一一对应。

% Original page 6.
\section*{2022A02（续）}
现在按各个非零\allowbreak{}\hskip.1em \RB{8.872}{5.322}{3.762}{%
\RG{0.000}{0.000}{10.9091}{mi}{101}\RG{5.080}{-1.636}{7.9701}{mi}{106}%
}\hskip.1em\allowbreak{} 的取值分类，约定当\allowbreak{}\hskip.1em \RB{25.318}{8.071}{0.936}{%
\RG{0.000}{0.000}{10.9091}{mi}{97}\RG{8.782}{0.000}{10.9091}{mi}{60}\RG{20.291}{0.000}{10.9091}{mi}{98}%
}\hskip.1em\allowbreak{} 时\allowbreak{}\hskip.1em \RB{34.546}{9.773}{4.342}{%
\RG{0.000}{8.837}{10.9092}{ex}{0}\RG{5.000}{4.841}{7.9701}{mi}{97}\RG{5.438}{-3.762}{7.9701}{mi}{98}\RG{9.498}{8.837}{10.9092}{ex}{1}\RG{17.529}{0.000}{10.9091}{rm}{61}\RG{29.027}{0.000}{10.9091}{rm}{48}%
}\hskip.1em\allowbreak{} 。

逆序数为\allowbreak{}\hskip.1em \RB{5.071}{7.763}{0.503}{%
\RG{0.000}{-0.003}{10.9091}{rm}{49}%
}\hskip.1em\allowbreak{} 时，唯一一个\allowbreak{}\hskip.1em \RB{29.061}{7.763}{3.766}{%
\RG{0.000}{-0.003}{10.9091}{mi}{101}\RG{5.079}{-1.640}{7.9701}{mi}{106}\RG{12.492}{-0.003}{10.9091}{rm}{61}\RG{23.990}{-0.003}{10.9091}{rm}{49}%
}\hskip.1em\allowbreak{} ，其下标可为\allowbreak{}\hskip.1em \RB{36.380}{7.763}{2.608}{%
\RG{0.000}{-0.003}{10.9091}{rm}{50}\RG{5.454}{-0.003}{10.9091}{mi}{59}\RG{10.287}{-0.003}{10.9091}{mi}{58}\RG{15.130}{-0.003}{10.9091}{mi}{58}\RG{19.974}{-0.003}{10.9091}{mi}{58}\RG{24.807}{-0.003}{10.9091}{mi}{59}\RG{29.651}{-0.003}{10.9091}{mi}{110}%
}\hskip.1em\allowbreak{} ，所以

\[
\RB{77.767}{8.685}{3.224}{%
\RG{0.000}{0.003}{10.9091}{mi}{78}\RG{9.954}{0.003}{10.9091}{rm}{40}\RG{14.197}{0.003}{10.9091}{mi}{110}\RG{20.743}{0.003}{10.9091}{mi}{59}\RG{25.593}{0.003}{10.9091}{rm}{49}\RG{31.048}{0.003}{10.9091}{rm}{41}\RG{38.302}{0.003}{10.9091}{rm}{61}\RG{49.836}{0.003}{10.9091}{mi}{110}\RG{58.808}{0.003}{10.9091}{sy}{0}\RG{69.718}{0.003}{10.9091}{rm}{49}\RG{75.173}{0.003}{10.9091}{mi}{58}%
}
\]
逆序数为\allowbreak{}\hskip.1em \RB{5.398}{7.765}{0.501}{%
\RG{0.000}{-0.001}{10.9091}{rm}{50}%
}\hskip.1em\allowbreak{} 时，只有两类：

•一个\allowbreak{}\hskip.1em \RB{29.389}{7.762}{3.766}{%
\RG{0.000}{-0.004}{10.9091}{mi}{101}\RG{5.080}{-1.640}{7.9701}{mi}{106}\RG{12.492}{-0.004}{10.9091}{rm}{61}\RG{23.990}{-0.004}{10.9091}{rm}{50}%
}\hskip.1em\allowbreak{} ：可取\allowbreak{}\hskip.1em \RB{56.043}{7.762}{2.740}{%
\RG{0.000}{-0.004}{10.9091}{mi}{106}\RG{8.147}{-0.004}{10.9091}{rm}{61}\RG{19.645}{-0.004}{10.9091}{rm}{51}\RG{25.117}{-0.004}{10.9091}{mi}{59}\RG{29.950}{-0.004}{10.9091}{mi}{58}\RG{34.794}{-0.004}{10.9091}{mi}{58}\RG{39.626}{-0.004}{10.9091}{mi}{58}\RG{44.470}{-0.004}{10.9091}{mi}{59}\RG{49.314}{-0.004}{10.9091}{mi}{110}%
}\hskip.1em\allowbreak{} ，共有\allowbreak{}\hskip.1em \RB{25.279}{7.762}{0.624}{%
\RG{0.000}{-0.004}{10.9091}{mi}{110}\RG{8.972}{-0.004}{10.9091}{sy}{0}\RG{19.881}{-0.004}{10.9091}{rm}{50}%
}\hskip.1em\allowbreak{} 种；

•两个\allowbreak{}\hskip.1em \RB{29.061}{7.761}{3.767}{%
\RG{0.000}{-0.005}{10.9091}{mi}{101}\RG{5.080}{-1.641}{7.9701}{mi}{106}\RG{12.492}{-0.005}{10.9091}{rm}{61}\RG{23.990}{-0.005}{10.9091}{rm}{49}%
}\hskip.1em\allowbreak{} ：从\allowbreak{}\hskip.1em \RB{36.380}{7.761}{2.610}{%
\RG{0.000}{-0.005}{10.9091}{rm}{50}\RG{5.454}{-0.005}{10.9091}{mi}{59}\RG{10.287}{-0.005}{10.9091}{mi}{58}\RG{15.131}{-0.005}{10.9091}{mi}{58}\RG{19.974}{-0.005}{10.9091}{mi}{58}\RG{24.807}{-0.005}{10.9091}{mi}{59}\RG{29.651}{-0.005}{10.9091}{mi}{110}%
}\hskip.1em\allowbreak{} 中选两个下标，共有\allowbreak{}\hskip.1em \RB{24.786}{10.636}{4.312}{%
\RG{0.000}{8.832}{10.9092}{ex}{0}\RG{5.000}{4.836}{7.9701}{mi}{110}\RG{10.138}{4.836}{7.9701}{sy}{0}\RG{10.862}{-3.767}{7.9701}{rm}{50}\RG{16.724}{4.836}{7.9701}{rm}{49}\RG{20.959}{8.832}{10.9092}{ex}{1}%
}\hskip.1em\allowbreak{} 种。

故当\allowbreak{}\hskip.1em \RB{26.491}{7.769}{1.992}{%
\RG{0.000}{0.003}{10.9091}{mi}{110}\RG{9.578}{0.003}{10.9091}{sy}{21}\RG{21.093}{0.003}{10.9091}{rm}{50}%
}\hskip.1em\allowbreak{} 时，

\[
\RB{207.868}{17.135}{10.857}{%
\RG{96.990}{15.378}{10.9092}{ex}{18}\RG{130.357}{15.378}{10.9092}{ex}{19}\RG{160.676}{11.335}{7.9701}{rm}{50}\RG{105.021}{7.376}{10.9091}{mi}{110}\RG{113.993}{7.376}{10.9091}{sy}{0}\RG{124.902}{7.376}{10.9091}{rm}{49}\RG{154.128}{7.376}{10.9091}{mi}{110}\RG{167.832}{7.376}{10.9091}{sy}{0}\RG{178.741}{7.376}{10.9091}{mi}{110}\RG{187.715}{7.376}{10.9091}{sy}{0}\RG{198.624}{7.376}{10.9091}{rm}{50}\RG{0.000}{-0.004}{10.9091}{mi}{78}\RG{9.954}{-0.004}{10.9091}{rm}{40}\RG{14.197}{-0.004}{10.9091}{mi}{110}\RG{20.743}{-0.004}{10.9091}{mi}{59}\RG{25.593}{-0.004}{10.9091}{rm}{50}\RG{31.048}{-0.004}{10.9091}{rm}{41}\RG{38.302}{-0.004}{10.9091}{rm}{61}\RG{49.812}{-0.004}{10.9091}{rm}{40}\RG{54.078}{-0.004}{10.9091}{mi}{110}\RG{63.050}{-0.004}{10.9091}{sy}{0}\RG{73.960}{-0.004}{10.9091}{rm}{50}\RG{79.414}{-0.004}{10.9091}{rm}{41}\RG{86.069}{-0.004}{10.9091}{rm}{43}\RG{141.417}{-0.004}{10.9091}{rm}{61}\RG{205.274}{-0.004}{10.9091}{mi}{58}\RG{114.961}{-7.487}{10.9091}{rm}{50}\RG{176.375}{-7.487}{10.9091}{rm}{50}\RR{154.128}{2.723}{49.949}{0.436}%
}
\]
逆序数为\allowbreak{}\hskip.1em \RB{5.485}{7.764}{0.742}{%
\RG{0.000}{-0.002}{10.9091}{rm}{51}%
}\hskip.1em\allowbreak{} 时，只有三类：

•一个\allowbreak{}\hskip.1em \RB{29.476}{7.762}{3.767}{%
\RG{0.000}{-0.004}{10.9091}{mi}{101}\RG{5.080}{-1.641}{7.9701}{mi}{106}\RG{12.492}{-0.004}{10.9091}{rm}{61}\RG{23.990}{-0.004}{10.9091}{rm}{51}%
}\hskip.1em\allowbreak{} ：\allowbreak{}\hskip.1em \RB{56.043}{7.882}{2.740}{%
\RG{0.000}{-0.004}{10.9091}{mi}{106}\RG{8.147}{-0.004}{10.9091}{rm}{61}\RG{19.645}{-0.004}{10.9091}{rm}{52}\RG{25.117}{-0.004}{10.9091}{mi}{59}\RG{29.950}{-0.004}{10.9091}{mi}{58}\RG{34.794}{-0.004}{10.9091}{mi}{58}\RG{39.626}{-0.004}{10.9091}{mi}{58}\RG{44.470}{-0.004}{10.9091}{mi}{59}\RG{49.314}{-0.004}{10.9091}{mi}{110}%
}\hskip.1em\allowbreak{} ，共有\allowbreak{}\hskip.1em \RB{25.367}{7.762}{0.744}{%
\RG{0.000}{-0.004}{10.9091}{mi}{110}\RG{8.972}{-0.004}{10.9091}{sy}{0}\RG{19.881}{-0.004}{10.9091}{rm}{51}%
}\hskip.1em\allowbreak{} 种；

•一个\allowbreak{}\hskip.1em \RB{29.389}{7.765}{3.764}{%
\RG{0.000}{-0.001}{10.9091}{mi}{101}\RG{5.080}{-1.638}{7.9701}{mi}{106}\RG{12.492}{-0.001}{10.9091}{rm}{61}\RG{23.990}{-0.001}{10.9091}{rm}{50}%
}\hskip.1em\allowbreak{} 、另一个\allowbreak{}\hskip.1em \RB{28.636}{7.765}{2.375}{%
\RG{0.000}{-0.001}{10.9091}{mi}{101}\RG{5.080}{-1.779}{7.9701}{mi}{96}\RG{12.067}{-0.001}{10.9091}{rm}{61}\RG{23.565}{-0.001}{10.9091}{rm}{49}%
}\hskip.1em\allowbreak{} ：\allowbreak{}\hskip.1em \RB{4.831}{7.710}{2.737}{%
\RG{0.000}{-0.001}{10.9091}{mi}{106}%
}\hskip.1em\allowbreak{} 有\allowbreak{}\hskip.1em \RB{25.203}{7.765}{0.621}{%
\RG{0.000}{-0.001}{10.9091}{mi}{110}\RG{8.934}{-0.001}{10.9091}{sy}{0}\RG{19.805}{-0.001}{10.9091}{rm}{50}%
}\hskip.1em\allowbreak{} 种选择，再从\allowbreak{}\hskip.1em \RB{36.380}{7.765}{2.606}{%
\RG{0.000}{-0.001}{10.9091}{rm}{50}\RG{5.454}{-0.001}{10.9091}{mi}{59}\RG{10.287}{-0.001}{10.9091}{mi}{58}\RG{15.130}{-0.001}{10.9091}{mi}{58}\RG{19.974}{-0.001}{10.9091}{mi}{58}\RG{24.807}{-0.001}{10.9091}{mi}{59}\RG{29.651}{-0.001}{10.9091}{mi}{110}%
}\hskip.1em\allowbreak{} 中除去\allowbreak{}\hskip.1em \RB{4.831}{7.710}{2.737}{%
\RG{0.000}{-0.001}{10.9091}{mi}{106}%
}\hskip.1em\allowbreak{} 后选\allowbreak{}\hskip.1em \RB{4.842}{8.190}{0.632}{%
\RG{0.000}{-0.001}{10.9091}{mi}{96}%
}\hskip.1em\allowbreak{} ，有\allowbreak{}\hskip.1em \RB{25.203}{7.765}{0.621}{%
\RG{0.000}{-0.001}{10.9091}{mi}{110}\RG{8.934}{-0.001}{10.9091}{sy}{0}\RG{19.805}{-0.001}{10.9091}{rm}{50}%
}\hskip.1em\allowbreak{} 种选择，共\allowbreak{}\hskip.1em \RB{38.123}{9.762}{3.223}{%
\RG{0.000}{0.004}{10.9091}{rm}{40}\RG{4.243}{0.004}{10.9091}{mi}{110}\RG{13.215}{0.004}{10.9091}{sy}{0}\RG{24.124}{0.004}{10.9091}{rm}{50}\RG{29.579}{0.004}{10.9091}{rm}{41}\RG{33.821}{3.962}{7.9701}{rm}{50}%
}\hskip.1em\allowbreak{} 种；

•三个\allowbreak{}\hskip.1em \RB{29.061}{7.763}{3.766}{%
\RG{0.000}{-0.003}{10.9091}{mi}{101}\RG{5.080}{-1.640}{7.9701}{mi}{106}\RG{12.492}{-0.003}{10.9091}{rm}{61}\RG{23.990}{-0.003}{10.9091}{rm}{49}%
}\hskip.1em\allowbreak{} ：共有\allowbreak{}\hskip.1em \RB{24.786}{10.637}{4.432}{%
\RG{0.000}{8.833}{10.9092}{ex}{0}\RG{5.000}{4.837}{7.9701}{mi}{110}\RG{10.138}{4.837}{7.9701}{sy}{0}\RG{10.862}{-3.765}{7.9701}{rm}{51}\RG{16.724}{4.837}{7.9701}{rm}{49}\RG{20.959}{8.833}{10.9092}{ex}{1}%
}\hskip.1em\allowbreak{} 种。

因此，当\allowbreak{}\hskip.1em \RB{26.579}{7.769}{1.991}{%
\RG{0.000}{0.004}{10.9091}{mi}{110}\RG{9.578}{0.004}{10.9091}{sy}{21}\RG{21.093}{0.004}{10.9091}{rm}{51}%
}\hskip.1em\allowbreak{} 时，

\[
\RB{188.493}{16.323}{39.042}{%
\RG{148.877}{15.387}{10.9092}{ex}{18}\RG{182.243}{15.387}{10.9092}{ex}{19}\RG{156.907}{7.385}{10.9091}{mi}{110}\RG{165.879}{7.385}{10.9091}{sy}{0}\RG{176.788}{7.385}{10.9091}{rm}{49}\RG{130.811}{4.509}{7.9701}{rm}{50}\RG{0.000}{0.005}{10.9091}{mi}{78}\RG{9.954}{0.005}{10.9091}{rm}{40}\RG{14.197}{0.005}{10.9091}{mi}{110}\RG{20.743}{0.005}{10.9091}{mi}{59}\RG{25.593}{0.005}{10.9091}{rm}{51}\RG{31.048}{0.005}{10.9091}{rm}{41}\RG{38.302}{0.005}{10.9091}{rm}{61}\RG{49.812}{0.005}{10.9091}{rm}{40}\RG{54.078}{0.005}{10.9091}{mi}{110}\RG{63.050}{0.005}{10.9091}{sy}{0}\RG{73.960}{0.005}{10.9091}{rm}{51}\RG{79.414}{0.005}{10.9091}{rm}{41}\RG{86.069}{0.005}{10.9091}{rm}{43}\RG{96.967}{0.005}{10.9091}{rm}{40}\RG{101.233}{0.005}{10.9091}{mi}{110}\RG{110.205}{0.005}{10.9091}{sy}{0}\RG{121.114}{0.005}{10.9091}{rm}{50}\RG{126.569}{0.005}{10.9091}{rm}{41}\RG{137.968}{0.005}{10.9091}{rm}{43}\RG{166.848}{-7.478}{10.9091}{rm}{51}\RG{57.579}{-19.480}{7.9701}{rm}{51}\RG{51.031}{-23.439}{10.9091}{mi}{110}\RG{64.736}{-23.439}{10.9091}{sy}{0}\RG{75.645}{-23.439}{10.9091}{rm}{55}\RG{81.100}{-23.439}{10.9091}{mi}{110}\RG{38.320}{-30.819}{10.9091}{rm}{61}\RG{88.844}{-30.819}{10.9091}{mi}{58}\RG{66.613}{-38.302}{10.9091}{rm}{54}\RR{51.031}{-28.092}{36.616}{0.436}%
}
\]
这就证明了全部结论。边界情形\allowbreak{}\hskip.1em \RB{26.475}{7.762}{0.624}{%
\RG{0.000}{-0.004}{10.9091}{mi}{110}\RG{9.578}{-0.004}{10.9091}{rm}{61}\RG{21.076}{-0.004}{10.9091}{rm}{50}%
}\hskip.1em\allowbreak{} 时\allowbreak{}\hskip.1em \RB{54.237}{8.678}{3.231}{%
\RG{0.000}{-0.004}{10.9091}{mi}{78}\RG{9.955}{-0.004}{10.9091}{rm}{40}\RG{14.188}{-0.004}{10.9091}{rm}{50}\RG{19.652}{-0.004}{10.9091}{mi}{59}\RG{24.500}{-0.004}{10.9091}{rm}{50}\RG{29.955}{-0.004}{10.9091}{rm}{41}\RG{37.209}{-0.004}{10.9091}{rm}{61}\RG{48.719}{-0.004}{10.9091}{rm}{48}%
}\hskip.1em\allowbreak{} ，\allowbreak{}\hskip.1em \RB{26.563}{7.762}{0.744}{%
\RG{0.000}{-0.004}{10.9091}{mi}{110}\RG{9.579}{-0.004}{10.9091}{rm}{61}\RG{21.077}{-0.004}{10.9091}{rm}{51}%
}\hskip.1em\allowbreak{} 时\allowbreak{}\hskip.1em \RB{53.790}{8.678}{3.231}{%
\RG{0.000}{-0.004}{10.9091}{mi}{78}\RG{9.954}{-0.004}{10.9091}{rm}{40}\RG{14.187}{-0.004}{10.9091}{rm}{51}\RG{19.651}{-0.004}{10.9091}{mi}{59}\RG{24.500}{-0.004}{10.9091}{rm}{51}\RG{29.955}{-0.004}{10.9091}{rm}{41}\RG{37.209}{-0.004}{10.9091}{rm}{61}\RG{48.719}{-0.004}{10.9091}{rm}{49}%
}\hskip.1em\allowbreak{} ，也与公式一致。\allowbreak{}\hskip.1em \RB{8.366}{7.991}{0.504}{%
\RG{0.000}{-0.004}{10.9091}{am}{3}%
}\hskip.1em\allowbreak{}

方法小结：相邻逆序交换使逆序数恰减一，适合归纳与构造；精确计数则改用独立的逆序编码\allowbreak{}\hskip.1em \RB{67.528}{7.766}{3.762}{%
\RG{0.000}{0.000}{10.9091}{rm}{48}\RG{8.485}{0.000}{10.9091}{sy}{20}\RG{20.000}{0.000}{10.9091}{mi}{101}\RG{25.079}{-1.636}{7.9701}{mi}{106}\RG{32.492}{0.000}{10.9091}{sy}{20}\RG{44.007}{0.000}{10.9091}{mi}{106}\RG{51.548}{0.000}{10.9091}{sy}{0}\RG{62.458}{0.000}{10.9091}{rm}{49}%
}\hskip.1em\allowbreak{} ，把排列问题化为小整数的分拆计数。题目来源：谢启鸿《复旦大学高等代数习题集（2023版）》，第123页。

% Original page 7.
\section*{2022A03}
\allowbreak{}\hskip.1em \RB{6.729}{5.322}{0.620}{%
\RG{0.000}{0.000}{10.9091}{mi}{110}%
}\hskip.1em\allowbreak{} 阶行列式的组合定义为

\[
\RB{236.909}{10.978}{16.293}{%
\RG{57.976}{10.359}{10.9092}{ex}{88}\RG{114.435}{4.500}{7.9701}{mi}{78}\RG{122.006}{4.500}{7.9701}{rm}{40}\RG{125.299}{4.500}{7.9701}{mi}{107}\RG{133.860}{4.500}{7.9701}{mi}{59}\RG{136.211}{4.500}{7.9701}{mi}{58}\RG{138.562}{4.500}{7.9701}{mi}{58}\RG{140.914}{4.500}{7.9701}{mi}{58}\RG{143.265}{4.500}{7.9701}{mi}{59}\RG{145.616}{4.500}{7.9701}{mi}{107}\RG{155.129}{4.500}{7.9701}{rm}{41}\RG{150.031}{3.504}{5.9776}{mi}{110}\RG{129.708}{3.393}{5.9776}{rm}{49}\RT{0.000}{-0.004}{10.9091}{d}\RT{6.066}{-0.004}{10.9091}{e}\RT{10.909}{-0.004}{10.9091}{t}\RG{16.971}{-0.004}{10.9091}{mi}{65}\RG{28.183}{-0.004}{10.9091}{rm}{61}\RG{92.011}{-0.004}{10.9091}{rm}{40}\RG{96.254}{-0.004}{10.9091}{sy}{0}\RG{104.738}{-0.004}{10.9091}{rm}{49}\RG{110.192}{-0.004}{10.9091}{rm}{41}\RG{158.921}{-0.004}{10.9091}{mi}{97}\RG{177.980}{-0.004}{10.9091}{mi}{97}\RG{198.858}{-0.004}{10.9091}{sy}{1}\RG{203.690}{-0.004}{10.9091}{sy}{1}\RG{208.534}{-0.004}{10.9091}{sy}{1}\RG{213.404}{-0.004}{10.9091}{mi}{97}\RG{234.315}{-0.004}{10.9091}{mi}{58}\RG{164.687}{-1.782}{7.9701}{mi}{107}\RG{173.248}{-1.782}{7.9701}{rm}{49}\RG{183.747}{-1.782}{7.9701}{mi}{107}\RG{192.308}{-1.782}{7.9701}{rm}{50}\RG{219.170}{-1.782}{7.9701}{mi}{107}\RG{228.679}{-1.782}{7.9701}{mi}{110}\RG{223.580}{-2.778}{5.9776}{mi}{110}\RG{169.097}{-2.888}{5.9776}{rm}{49}\RG{188.156}{-2.888}{5.9776}{rm}{50}\RG{39.699}{-13.800}{7.9701}{rm}{40}\RG{42.992}{-13.800}{7.9701}{mi}{107}\RG{51.552}{-13.800}{7.9701}{mi}{59}\RG{53.904}{-13.800}{7.9701}{mi}{58}\RG{56.255}{-13.800}{7.9701}{mi}{58}\RG{58.606}{-13.800}{7.9701}{mi}{58}\RG{60.957}{-13.800}{7.9701}{mi}{59}\RG{63.308}{-13.800}{7.9701}{mi}{107}\RG{72.823}{-13.800}{7.9701}{rm}{41}\RG{76.116}{-13.800}{7.9701}{sy}{50}\RG{81.762}{-13.800}{7.9701}{mi}{83}\RG{67.724}{-14.796}{5.9776}{mi}{110}\RG{86.912}{-14.796}{5.9776}{mi}{110}\RG{47.401}{-14.907}{5.9776}{rm}{49}%
}
\]
证明：当\allowbreak{}\hskip.1em \RB{26.579}{7.770}{1.991}{%
\RG{0.000}{0.004}{10.9091}{mi}{110}\RG{9.578}{0.004}{10.9091}{sy}{21}\RG{21.093}{0.004}{10.9091}{rm}{51}%
}\hskip.1em\allowbreak{} 时，不存在实矩阵\allowbreak{}\hskip.1em \RB{43.655}{8.686}{3.758}{%
\RG{0.000}{0.004}{10.9091}{mi}{65}\RG{11.212}{0.004}{10.9091}{rm}{61}\RG{22.710}{0.004}{10.9091}{rm}{40}\RG{26.970}{0.004}{10.9091}{mi}{97}\RG{32.736}{-1.632}{7.9701}{mi}{105}\RG{35.613}{-1.632}{7.9701}{mi}{106}\RG{40.002}{0.004}{10.9091}{rm}{41}%
}\hskip.1em\allowbreak{} ，使每个排列所对应的带符号项都严格为正。

证明：反设存在这样的矩阵。固定第\allowbreak{}\hskip.1em \RB{36.380}{7.886}{2.605}{%
\RG{0.000}{-0.000}{10.9091}{rm}{52}\RG{5.454}{-0.000}{10.9091}{mi}{59}\RG{10.287}{-0.000}{10.9091}{mi}{58}\RG{15.130}{-0.000}{10.9091}{mi}{58}\RG{19.974}{-0.000}{10.9091}{mi}{58}\RG{24.807}{-0.000}{10.9091}{mi}{59}\RG{29.651}{-0.000}{10.9091}{mi}{110}%
}\hskip.1em\allowbreak{} 列所选择的行分别为\allowbreak{}\hskip.1em \RB{36.381}{7.886}{2.605}{%
\RG{0.000}{-0.000}{10.9091}{rm}{52}\RG{5.455}{-0.000}{10.9091}{mi}{59}\RG{10.288}{-0.000}{10.9091}{mi}{58}\RG{15.131}{-0.000}{10.9091}{mi}{58}\RG{19.964}{-0.000}{10.9091}{mi}{58}\RG{24.808}{-0.000}{10.9091}{mi}{59}\RG{29.652}{-0.000}{10.9091}{mi}{110}%
}\hskip.1em\allowbreak{} ，只让前三列所选的行在\allowbreak{}\hskip.1em \RB{26.092}{7.771}{2.600}{%
\RG{0.000}{0.005}{10.9091}{rm}{49}\RG{5.454}{0.005}{10.9091}{mi}{59}\RG{10.303}{0.005}{10.9091}{rm}{50}\RG{15.757}{0.005}{10.9091}{mi}{59}\RG{20.606}{0.005}{10.9091}{rm}{51}%
}\hskip.1em\allowbreak{} 中排列。记

\[
\RB{52.450}{17.650}{15.214}{%
\RG{24.049}{13.635}{7.9701}{mi}{110}\RG{19.649}{10.362}{10.9092}{ex}{89}\RG{0.000}{-0.001}{10.9091}{mi}{99}\RG{7.751}{-0.001}{10.9091}{rm}{61}\RG{35.789}{-0.001}{10.9091}{mi}{97}\RG{49.822}{-0.001}{10.9091}{rm}{59}\RG{41.555}{-1.638}{7.9701}{mi}{106}\RG{45.429}{-1.638}{7.9701}{mi}{106}\RG{19.266}{-13.088}{7.9701}{mi}{106}\RG{23.150}{-13.088}{7.9701}{rm}{61}\RG{29.734}{-13.088}{7.9701}{rm}{52}%
}
\]
当\allowbreak{}\hskip.1em \RB{26.562}{7.764}{0.742}{%
\RG{0.000}{-0.002}{10.9091}{mi}{110}\RG{9.578}{-0.002}{10.9091}{rm}{61}\RG{21.076}{-0.002}{10.9091}{rm}{51}%
}\hskip.1em\allowbreak{} 时，空乘积规定为\allowbreak{}\hskip.1em \RB{24.320}{7.764}{0.622}{%
\RG{0.000}{-0.002}{10.9091}{mi}{99}\RG{7.751}{-0.002}{10.9091}{rm}{61}\RG{19.249}{-0.002}{10.9091}{rm}{49}%
}\hskip.1em\allowbreak{} 。恒等排列对应的项严格为正，故\allowbreak{}\hskip.1em \RB{24.768}{8.309}{2.858}{%
\RG{0.000}{-0.002}{10.9091}{mi}{99}\RG{7.751}{-0.002}{10.9091}{sy}{54}\RG{7.751}{-0.002}{10.9091}{rm}{61}\RG{19.249}{-0.002}{10.9091}{rm}{48}%
}\hskip.1em\allowbreak{} 。

对\allowbreak{}\hskip.1em \RB{11.063}{8.195}{2.299}{%
\RG{0.000}{0.004}{10.9091}{mi}{83}\RG{6.689}{-1.632}{7.9701}{rm}{51}%
}\hskip.1em\allowbreak{} 中的每个排列\allowbreak{}\hskip.1em \RB{6.686}{5.206}{0.616}{%
\RG{0.000}{0.004}{10.9091}{mi}{27}%
}\hskip.1em\allowbreak{} ，对应的\allowbreak{}\hskip.1em \RB{6.729}{5.326}{0.616}{%
\RG{0.000}{0.004}{10.9091}{mi}{110}%
}\hskip.1em\allowbreak{} 元排列为\allowbreak{}\hskip.1em \RB{120.380}{8.686}{3.223}{%
\RG{0.000}{0.004}{10.9091}{rm}{40}\RG{4.243}{0.004}{10.9091}{mi}{27}\RG{10.868}{0.004}{10.9091}{rm}{40}\RG{15.101}{0.004}{10.9091}{rm}{49}\RG{20.555}{0.004}{10.9091}{rm}{41}\RG{24.807}{0.004}{10.9091}{mi}{59}\RG{29.640}{0.004}{10.9091}{mi}{27}\RG{36.281}{0.004}{10.9091}{rm}{40}\RG{40.514}{0.004}{10.9091}{rm}{50}\RG{45.969}{0.004}{10.9091}{rm}{41}\RG{50.220}{0.004}{10.9091}{mi}{59}\RG{55.053}{0.004}{10.9091}{mi}{27}\RG{61.695}{0.004}{10.9091}{rm}{40}\RG{65.927}{0.004}{10.9091}{rm}{51}\RG{71.382}{0.004}{10.9091}{rm}{41}\RG{75.634}{0.004}{10.9091}{mi}{59}\RG{80.482}{0.004}{10.9091}{rm}{52}\RG{85.937}{0.004}{10.9091}{mi}{59}\RG{90.769}{0.004}{10.9091}{mi}{58}\RG{95.613}{0.004}{10.9091}{mi}{58}\RG{100.446}{0.004}{10.9091}{mi}{58}\RG{105.290}{0.004}{10.9091}{mi}{59}\RG{110.133}{0.004}{10.9091}{mi}{110}\RG{116.727}{0.004}{10.9091}{rm}{41}%
}\hskip.1em\allowbreak{} ，其逆序数与\allowbreak{}\hskip.1em \RB{6.686}{5.206}{0.616}{%
\RG{0.000}{0.004}{10.9091}{mi}{27}%
}\hskip.1em\allowbreak{} 相同，故假设给出六个正数

\[
\RB{134.863}{19.431}{15.218}{%
\RG{72.884}{13.631}{7.9701}{rm}{51}\RG{68.031}{10.358}{10.9092}{ex}{89}\RG{0.000}{-0.005}{10.9091}{mi}{84}\RG{15.029}{-0.005}{10.9091}{rm}{61}\RG{26.544}{-0.005}{10.9091}{mi}{99}\RT{34.901}{-0.005}{10.9091}{s}\RT{39.200}{-0.005}{10.9091}{g}\RT{44.654}{-0.005}{10.9091}{n}\RG{50.720}{-0.005}{10.9091}{rm}{40}\RG{54.962}{-0.005}{10.9091}{mi}{27}\RG{61.588}{-0.005}{10.9091}{rm}{41}\RG{84.171}{-0.005}{10.9091}{mi}{97}\RG{115.298}{-0.005}{10.9091}{mi}{62}\RG{126.814}{-0.005}{10.9091}{rm}{48}\RG{132.268}{-0.005}{10.9091}{mi}{58}\RG{6.375}{-1.642}{7.9701}{mi}{27}\RG{89.938}{-2.225}{7.9701}{mi}{27}\RG{95.063}{-2.225}{7.9701}{rm}{40}\RG{98.356}{-2.225}{7.9701}{mi}{106}\RG{102.240}{-2.225}{7.9701}{rm}{41}\RG{105.534}{-2.225}{7.9701}{mi}{59}\RG{107.885}{-2.225}{7.9701}{mi}{106}\RG{67.649}{-13.092}{7.9701}{mi}{106}\RG{71.533}{-13.092}{7.9701}{rm}{61}\RG{78.116}{-13.092}{7.9701}{rm}{49}%
}
\]
把这六个不等式相乘。一方面，左端乘积应当为正；另一方面，\allowbreak{}\hskip.1em \RB{11.063}{8.189}{2.306}{%
\RG{0.000}{-0.002}{10.9091}{mi}{83}\RG{6.689}{-1.639}{7.9701}{rm}{51}%
}\hskip.1em\allowbreak{} 中有三个奇排列、三个偶排列，故六个符号的乘积为\allowbreak{}\hskip.1em \RB{13.555}{7.769}{0.500}{%
\RG{0.000}{0.003}{10.9091}{sy}{0}\RG{8.484}{0.003}{10.9091}{rm}{49}%
}\hskip.1em\allowbreak{} 。

此外，对任意固定的\allowbreak{}\hskip.1em \RB{54.065}{7.765}{2.737}{%
\RG{0.000}{-0.001}{10.9091}{rm}{49}\RG{8.562}{-0.001}{10.9091}{sy}{20}\RG{20.155}{-0.001}{10.9091}{mi}{105}\RG{23.908}{-0.001}{10.9091}{mi}{59}\RG{28.741}{-0.001}{10.9091}{mi}{106}\RG{36.986}{-0.001}{10.9091}{sy}{20}\RG{48.579}{-0.001}{10.9091}{rm}{51}%
}\hskip.1em\allowbreak{} ，满足\allowbreak{}\hskip.1em \RB{38.624}{8.681}{3.228}{%
\RG{0.000}{-0.001}{10.9091}{mi}{27}\RG{6.625}{-0.001}{10.9091}{rm}{40}\RG{10.868}{-0.001}{10.9091}{mi}{106}\RG{15.985}{-0.001}{10.9091}{rm}{41}\RG{23.316}{-0.001}{10.9091}{rm}{61}\RG{34.927}{-0.001}{10.9091}{mi}{105}%
}\hskip.1em\allowbreak{} 的排列恰有\allowbreak{}\hskip.1em \RB{5.398}{7.765}{0.501}{%
\RG{0.000}{-0.001}{10.9091}{rm}{50}%
}\hskip.1em\allowbreak{} 个，所以每个\allowbreak{}\hskip.1em \RB{12.435}{5.321}{3.764}{%
\RG{0.000}{-0.001}{10.9091}{mi}{97}\RG{5.766}{-1.638}{7.9701}{mi}{105}\RG{8.643}{-1.638}{7.9701}{mi}{106}%
}\hskip.1em\allowbreak{} 在六项的乘积中恰出现两次。因此

\[
\RB{136.099}{19.437}{15.213}{%
\RG{73.044}{13.637}{7.9701}{rm}{51}\RG{89.185}{13.637}{7.9701}{rm}{51}\RG{3.067}{10.364}{10.9092}{ex}{89}\RG{68.192}{10.364}{10.9092}{ex}{89}\RG{84.332}{10.364}{10.9092}{ex}{89}\RG{106.239}{5.308}{7.9701}{mi}{59}\RG{108.591}{5.308}{7.9701}{rm}{50}\RG{61.641}{4.505}{7.9701}{rm}{54}\RG{21.891}{0.000}{10.9091}{mi}{84}\RG{36.919}{0.000}{10.9091}{rm}{61}\RG{48.435}{0.000}{10.9091}{sy}{0}\RG{56.919}{0.000}{10.9091}{mi}{99}\RG{100.472}{0.000}{10.9091}{mi}{97}\RG{116.534}{0.000}{10.9091}{mi}{60}\RG{128.049}{0.000}{10.9091}{rm}{48}\RG{133.504}{0.000}{10.9091}{mi}{58}\RG{28.266}{-1.636}{7.9701}{mi}{27}\RG{106.239}{-3.256}{7.9701}{mi}{105}\RG{109.116}{-3.256}{7.9701}{mi}{106}\RG{68.310}{-13.087}{7.9701}{mi}{105}\RG{71.193}{-13.087}{7.9701}{rm}{61}\RG{77.776}{-13.087}{7.9701}{rm}{49}\RG{83.949}{-13.087}{7.9701}{mi}{106}\RG{87.833}{-13.087}{7.9701}{rm}{61}\RG{94.416}{-13.087}{7.9701}{rm}{49}\RG{0.000}{-13.264}{7.9701}{mi}{27}\RG{5.126}{-13.264}{7.9701}{sy}{50}\RG{10.771}{-13.264}{7.9701}{mi}{83}\RG{15.922}{-14.371}{5.9776}{rm}{51}%
}
\]
严格小于零是因为每个\allowbreak{}\hskip.1em \RB{12.436}{5.317}{3.767}{%
\RG{0.000}{-0.005}{10.9091}{mi}{97}\RG{5.767}{-1.641}{7.9701}{mi}{105}\RG{8.644}{-1.641}{7.9701}{mi}{106}%
}\hskip.1em\allowbreak{} 都出现在某个严格正的\allowbreak{}\hskip.1em \RB{11.665}{7.881}{2.221}{%
\RG{0.000}{-0.005}{10.9091}{mi}{84}\RG{6.375}{-1.641}{7.9701}{mi}{27}%
}\hskip.1em\allowbreak{} 中，故均不为零。这与六个正数之积为正矛盾，命题得证。\allowbreak{}\hskip.1em \RB{8.366}{7.995}{0.500}{%
\RG{0.000}{0.000}{10.9091}{am}{3}%
}\hskip.1em\allowbreak{}

方法小结：只需固定其余指标，将问题压缩到三个指标。六个行列式项相乘后，矩阵元素全部变成平方，而排列符号仍留下一个负号。

题目来源：谢启鸿《复旦大学高等代数习题集（2023版）》，第123页。

% Original page 8.
\section*{2022A04}
设\allowbreak{}\hskip.1em \RB{44.398}{8.682}{3.763}{%
\RG{0.000}{0.000}{10.9091}{mi}{65}\RG{11.584}{0.000}{10.9091}{rm}{61}\RG{23.453}{0.000}{10.9091}{rm}{40}\RG{27.714}{0.000}{10.9091}{mi}{97}\RG{33.480}{-1.637}{7.9701}{mi}{105}\RG{36.358}{-1.637}{7.9701}{mi}{106}\RG{40.745}{0.000}{10.9091}{rm}{41}%
}\hskip.1em\allowbreak{} 为\allowbreak{}\hskip.1em \RB{6.729}{5.322}{0.620}{%
\RG{0.000}{0.000}{10.9091}{mi}{110}%
}\hskip.1em\allowbreak{} 阶非负矩阵，且每行、每列的元素之和都为\allowbreak{}\hskip.1em \RB{5.071}{7.766}{0.500}{%
\RG{0.000}{0.000}{10.9091}{rm}{49}%
}\hskip.1em\allowbreak{} 。设\allowbreak{}\hskip.1em \RB{76.501}{8.682}{3.227}{%
\RG{0.000}{0.000}{10.9091}{mi}{120}\RG{9.637}{0.000}{10.9091}{rm}{61}\RG{21.506}{0.000}{10.9091}{rm}{40}\RG{25.767}{0.000}{10.9091}{mi}{120}\RG{32.002}{-1.637}{7.9701}{rm}{49}\RG{36.734}{0.000}{10.9091}{mi}{59}\RG{41.567}{0.000}{10.9091}{mi}{58}\RG{46.411}{0.000}{10.9091}{mi}{58}\RG{51.254}{0.000}{10.9091}{mi}{58}\RG{56.087}{0.000}{10.9091}{mi}{59}\RG{60.931}{0.000}{10.9091}{mi}{120}\RG{67.212}{-1.637}{7.9701}{mi}{110}\RG{72.849}{0.000}{10.9091}{rm}{41}%
}\hskip.1em\allowbreak{} T为非负向量，\allowbreak{}\hskip.1em \RB{102.452}{8.677}{3.232}{%
\RG{0.000}{-0.005}{10.9091}{mi}{121}\RG{8.770}{-0.005}{10.9091}{rm}{61}\RG{20.285}{-0.005}{10.9091}{mi}{65}\RG{28.467}{-0.005}{10.9091}{mi}{120}\RG{37.732}{-0.005}{10.9091}{rm}{61}\RG{49.231}{-0.005}{10.9091}{rm}{40}\RG{53.490}{-0.005}{10.9091}{mi}{121}\RG{58.838}{-1.642}{7.9701}{rm}{49}\RG{63.572}{-0.005}{10.9091}{mi}{59}\RG{68.404}{-0.005}{10.9091}{mi}{58}\RG{73.248}{-0.005}{10.9091}{mi}{58}\RG{78.081}{-0.005}{10.9091}{mi}{58}\RG{82.924}{-0.005}{10.9091}{mi}{59}\RG{87.768}{-0.005}{10.9091}{mi}{121}\RG{93.163}{-1.642}{7.9701}{mi}{110}\RG{98.799}{-0.005}{10.9091}{rm}{41}%
}\hskip.1em\allowbreak{} T。证明

\[
\RB{68.768}{30.744}{2.121}{%
\RG{4.401}{26.729}{7.9701}{mi}{110}\RG{43.816}{26.729}{7.9701}{mi}{110}\RG{0.000}{23.456}{10.9092}{ex}{89}\RG{39.416}{23.456}{10.9092}{ex}{89}\RG{15.758}{13.092}{10.9091}{mi}{121}\RG{27.518}{13.092}{10.9091}{sy}{21}\RG{55.556}{13.092}{10.9091}{mi}{120}\RG{66.174}{13.092}{10.9091}{mi}{58}\RG{21.106}{11.456}{7.9701}{mi}{105}\RG{61.792}{11.456}{7.9701}{mi}{106}\RG{0.118}{0.005}{7.9701}{mi}{105}\RG{3.001}{0.005}{7.9701}{rm}{61}\RG{9.584}{0.005}{7.9701}{rm}{49}\RG{39.033}{0.005}{7.9701}{mi}{106}\RG{42.917}{0.005}{7.9701}{rm}{61}\RG{49.501}{0.005}{7.9701}{rm}{49}%
}
\]
证明：若有某个\allowbreak{}\hskip.1em \RB{30.679}{7.769}{3.760}{%
\RG{0.000}{0.003}{10.9091}{mi}{120}\RG{6.234}{-1.634}{7.9701}{mi}{106}\RG{13.651}{0.003}{10.9091}{rm}{61}\RG{25.160}{0.003}{10.9091}{rm}{48}%
}\hskip.1em\allowbreak{} ，则右端为零；由\allowbreak{}\hskip.1em \RB{33.103}{7.769}{3.760}{%
\RG{0.000}{0.003}{10.9091}{mi}{97}\RG{5.766}{-1.634}{7.9701}{mi}{105}\RG{8.643}{-1.634}{7.9701}{mi}{106}\RG{16.066}{0.003}{10.9091}{sy}{21}\RG{27.585}{0.003}{10.9091}{rm}{48}%
}\hskip.1em\allowbreak{} 、\allowbreak{}\hskip.1em \RB{30.689}{7.769}{3.760}{%
\RG{0.000}{0.003}{10.9091}{mi}{120}\RG{6.235}{-1.634}{7.9701}{mi}{106}\RG{13.652}{0.003}{10.9091}{sy}{21}\RG{25.171}{0.003}{10.9091}{rm}{48}%
}\hskip.1em\allowbreak{} 可知各\allowbreak{}\hskip.1em \RB{28.802}{7.769}{2.733}{%
\RG{0.000}{0.003}{10.9091}{mi}{121}\RG{5.348}{-1.634}{7.9701}{mi}{105}\RG{11.764}{0.003}{10.9091}{sy}{21}\RG{23.284}{0.003}{10.9091}{rm}{48}%
}\hskip.1em\allowbreak{} ，故不等式成立。以下设所有\allowbreak{}\hskip.1em \RB{30.680}{7.764}{3.765}{%
\RG{0.000}{-0.002}{10.9091}{mi}{120}\RG{6.235}{-1.639}{7.9701}{mi}{106}\RG{13.647}{-0.002}{10.9091}{mi}{62}\RG{25.162}{-0.002}{10.9091}{rm}{48}%
}\hskip.1em\allowbreak{} 。

对每个固定的\allowbreak{}\hskip.1em \RB{3.696}{7.715}{0.616}{%
\RG{0.000}{0.004}{10.9091}{mi}{105}%
}\hskip.1em\allowbreak{} ，有\allowbreak{}\hskip.1em \RB{33.095}{7.770}{3.759}{%
\RG{0.000}{0.004}{10.9091}{mi}{97}\RG{5.766}{-1.633}{7.9701}{mi}{105}\RG{8.643}{-1.633}{7.9701}{mi}{106}\RG{16.062}{0.004}{10.9091}{sy}{21}\RG{27.577}{0.004}{10.9091}{rm}{48}%
}\hskip.1em\allowbreak{} 、\allowbreak{}\hskip.1em \RB{50.348}{8.686}{5.348}{%
\RG{0.000}{8.186}{10.9092}{ex}{80}\RG{11.516}{-3.222}{7.9701}{mi}{106}\RG{17.716}{0.004}{10.9091}{mi}{97}\RG{23.482}{-1.633}{7.9701}{mi}{105}\RG{26.359}{-1.633}{7.9701}{mi}{106}\RG{33.778}{0.004}{10.9091}{rm}{61}\RG{45.277}{0.004}{10.9091}{rm}{49}%
}\hskip.1em\allowbreak{} ，所以加权算术–几何平均不等式给出

\[
\RB{115.961}{17.656}{15.209}{%
\RG{28.585}{13.641}{7.9701}{mi}{110}\RG{83.829}{13.641}{7.9701}{mi}{110}\RG{23.275}{10.368}{10.9092}{ex}{88}\RG{79.429}{10.368}{10.9092}{ex}{89}\RG{101.804}{6.139}{7.9701}{mi}{97}\RG{106.302}{4.924}{5.9776}{mi}{105}\RG{108.962}{4.924}{5.9776}{mi}{106}\RG{0.000}{0.005}{10.9091}{mi}{121}\RG{11.760}{0.005}{10.9091}{rm}{61}\RG{40.851}{0.005}{10.9091}{mi}{97}\RG{53.883}{0.005}{10.9091}{mi}{120}\RG{67.531}{0.005}{10.9091}{sy}{21}\RG{95.569}{0.005}{10.9091}{mi}{120}\RG{113.366}{0.005}{10.9091}{mi}{58}\RG{5.349}{-1.632}{7.9701}{mi}{105}\RG{46.617}{-1.632}{7.9701}{mi}{105}\RG{49.495}{-1.632}{7.9701}{mi}{106}\RG{60.119}{-1.632}{7.9701}{mi}{106}\RG{101.804}{-3.252}{7.9701}{mi}{106}\RG{23.802}{-13.083}{7.9701}{mi}{106}\RG{27.686}{-13.083}{7.9701}{rm}{61}\RG{34.270}{-13.083}{7.9701}{rm}{49}\RG{79.046}{-13.083}{7.9701}{mi}{106}\RG{82.930}{-13.083}{7.9701}{rm}{61}\RG{89.513}{-13.083}{7.9701}{rm}{49}%
}
\]
为使论证自足，可用log\allowbreak{}\hskip.1em \RB{40.829}{7.770}{1.991}{%
\RG{0.000}{0.004}{10.9091}{mi}{116}\RG{6.970}{0.004}{10.9091}{sy}{20}\RG{18.485}{0.004}{10.9091}{mi}{116}\RG{24.849}{0.004}{10.9091}{sy}{0}\RG{35.758}{0.004}{10.9091}{rm}{49}%
}\hskip.1em\allowbreak{} （\allowbreak{}\hskip.1em \RB{24.003}{7.770}{0.932}{%
\RG{0.000}{0.004}{10.9091}{mi}{116}\RG{6.960}{0.004}{10.9091}{mi}{62}\RG{18.485}{0.004}{10.9091}{rm}{48}%
}\hskip.1em\allowbreak{} ）直接验证这一行不等式。因为\allowbreak{}\hskip.1em \RB{28.793}{7.770}{2.732}{%
\RG{0.000}{0.004}{10.9091}{mi}{121}\RG{5.348}{-1.632}{7.9701}{mi}{105}\RG{11.760}{0.004}{10.9091}{mi}{62}\RG{23.275}{0.004}{10.9091}{rm}{48}%
}\hskip.1em\allowbreak{} ，

\[
\RB{186.031}{17.654}{52.135}{%
\RG{108.166}{15.385}{10.9092}{ex}{18}\RG{147.992}{15.385}{10.9092}{ex}{19}\RG{5.310}{13.639}{7.9701}{mi}{110}\RG{81.050}{13.639}{7.9701}{mi}{110}\RG{0.000}{10.367}{10.9092}{ex}{88}\RG{75.740}{10.367}{10.9092}{ex}{88}\RG{49.381}{7.443}{10.9091}{mi}{120}\RG{117.392}{7.443}{10.9091}{mi}{120}\RG{55.616}{5.807}{7.9701}{mi}{106}\RG{123.626}{5.807}{7.9701}{mi}{106}\RG{17.576}{0.003}{10.9091}{mi}{97}\RT{32.426}{0.003}{10.9091}{l}\RT{35.459}{0.003}{10.9091}{o}\RT{40.914}{0.003}{10.9091}{g}\RG{64.224}{0.003}{10.9091}{sy}{20}\RG{93.316}{0.003}{10.9091}{mi}{97}\RG{131.628}{0.003}{10.9091}{sy}{0}\RG{142.537}{0.003}{10.9091}{rm}{49}\RG{23.342}{-1.633}{7.9701}{mi}{105}\RG{26.219}{-1.633}{7.9701}{mi}{106}\RG{99.082}{-1.633}{7.9701}{mi}{105}\RG{101.959}{-1.633}{7.9701}{mi}{106}\RG{50.325}{-7.480}{10.9091}{mi}{121}\RG{118.335}{-7.480}{10.9091}{mi}{121}\RG{55.673}{-9.117}{7.9701}{mi}{105}\RG{123.684}{-9.117}{7.9701}{mi}{105}\RG{0.527}{-13.084}{7.9701}{mi}{106}\RG{4.411}{-13.084}{7.9701}{rm}{61}\RG{10.994}{-13.084}{7.9701}{rm}{49}\RG{76.267}{-13.084}{7.9701}{mi}{106}\RG{80.151}{-13.084}{7.9701}{rm}{61}\RG{86.734}{-13.084}{7.9701}{rm}{49}\RG{76.936}{-19.710}{10.9092}{ex}{80}\RG{132.829}{-26.558}{10.9092}{ex}{88}\RG{94.651}{-27.892}{10.9091}{mi}{97}\RG{107.683}{-27.892}{10.9091}{mi}{120}\RG{100.418}{-29.528}{7.9701}{mi}{105}\RG{103.295}{-29.528}{7.9701}{mi}{106}\RG{113.918}{-29.528}{7.9701}{mi}{106}\RG{88.451}{-31.117}{7.9701}{mi}{106}\RG{64.224}{-36.922}{10.9091}{rm}{61}\RG{121.920}{-36.922}{10.9091}{sy}{0}\RG{150.404}{-36.922}{10.9091}{mi}{97}\RG{166.466}{-36.922}{10.9091}{rm}{61}\RG{177.965}{-36.922}{10.9091}{rm}{48}\RG{183.437}{-36.922}{10.9091}{mi}{58}\RG{156.171}{-38.558}{7.9701}{mi}{105}\RG{159.048}{-38.558}{7.9701}{mi}{106}\RG{93.253}{-44.405}{10.9091}{mi}{121}\RG{98.601}{-46.041}{7.9701}{mi}{105}\RG{138.765}{-50.009}{7.9701}{mi}{106}\RR{49.381}{2.730}{10.617}{0.436}\RR{117.392}{2.730}{10.617}{0.436}\RR{76.936}{-34.194}{41.364}{0.436}%
}
\]
这里log\allowbreak{}\hskip.1em \RB{40.828}{7.769}{1.992}{%
\RG{0.000}{0.003}{10.9091}{mi}{116}\RG{6.969}{0.003}{10.9091}{sy}{20}\RG{18.485}{0.003}{10.9091}{mi}{116}\RG{24.848}{0.003}{10.9091}{sy}{0}\RG{35.757}{0.003}{10.9091}{rm}{49}%
}\hskip.1em\allowbreak{} 可由函数\allowbreak{}\hskip.1em \RB{33.222}{7.769}{0.617}{%
\RG{0.000}{0.003}{10.9091}{mi}{116}\RG{6.363}{0.003}{10.9091}{sy}{0}\RG{17.272}{0.003}{10.9091}{rm}{49}\RG{25.151}{0.003}{10.9091}{sy}{0}%
}\hskip.1em\allowbreak{} log\allowbreak{}\hskip.1em \RB{4.100}{7.332}{0.617}{%
\RG{0.000}{0.003}{10.9091}{mi}{116}%
}\hskip.1em\allowbreak{} 在\allowbreak{}\hskip.1em \RB{23.539}{7.769}{0.617}{%
\RG{0.000}{0.003}{10.9091}{mi}{116}\RG{6.970}{0.003}{10.9091}{rm}{61}\RG{18.468}{0.003}{10.9091}{rm}{49}%
}\hskip.1em\allowbreak{} 处取最小值\allowbreak{}\hskip.1em \RB{5.518}{7.769}{0.737}{%
\RG{0.000}{0.003}{10.9091}{rm}{48}%
}\hskip.1em\allowbreak{} 得到。

将各行的不等式相乘，并使用每列的元素和为\allowbreak{}\hskip.1em \RB{5.071}{7.765}{0.501}{%
\RG{0.000}{-0.001}{10.9091}{rm}{49}%
}\hskip.1em\allowbreak{} ，得

\[
\RB{131.146}{30.741}{39.996}{%
\RG{4.400}{26.726}{7.9701}{mi}{110}\RG{43.433}{26.726}{7.9701}{mi}{110}\RG{59.573}{26.726}{7.9701}{mi}{110}\RG{0.000}{23.453}{10.9092}{ex}{89}\RG{39.032}{23.453}{10.9092}{ex}{89}\RG{55.173}{23.453}{10.9092}{ex}{89}\RG{77.549}{19.224}{7.9701}{mi}{97}\RG{82.047}{18.009}{5.9776}{mi}{105}\RG{84.707}{18.009}{5.9776}{mi}{106}\RG{15.757}{13.089}{10.9091}{mi}{121}\RG{27.517}{13.089}{10.9091}{sy}{21}\RG{71.314}{13.089}{10.9091}{mi}{120}\RG{21.106}{11.453}{7.9701}{mi}{105}\RG{77.549}{9.833}{7.9701}{mi}{106}\RG{0.117}{0.002}{7.9701}{mi}{105}\RG{3.000}{0.002}{7.9701}{rm}{61}\RG{9.583}{0.002}{7.9701}{rm}{49}\RG{39.150}{0.002}{7.9701}{mi}{105}\RG{42.033}{0.002}{7.9701}{rm}{61}\RG{48.617}{0.002}{7.9701}{rm}{49}\RG{54.790}{0.002}{7.9701}{mi}{106}\RG{58.674}{0.002}{7.9701}{rm}{61}\RG{65.258}{0.002}{7.9701}{rm}{49}\RG{43.816}{-11.146}{7.9701}{mi}{110}\RG{106.194}{-11.146}{7.9701}{mi}{110}\RG{61.791}{-12.557}{7.9702}{ex}{80}\RG{39.415}{-14.419}{10.9092}{ex}{89}\RG{101.794}{-14.419}{10.9092}{ex}{89}\RG{75.303}{-18.534}{7.9701}{mi}{97}\RG{79.801}{-19.749}{5.9776}{mi}{105}\RG{82.461}{-19.749}{5.9776}{mi}{106}\RG{70.730}{-21.025}{5.9776}{mi}{105}\RG{27.517}{-24.782}{10.9091}{rm}{61}\RG{55.555}{-24.782}{10.9091}{mi}{120}\RG{89.896}{-24.782}{10.9091}{rm}{61}\RG{117.934}{-24.782}{10.9091}{mi}{120}\RG{128.551}{-24.782}{10.9091}{mi}{58}\RG{124.169}{-26.419}{7.9701}{mi}{106}\RG{61.791}{-28.039}{7.9701}{mi}{106}\RG{39.032}{-37.870}{7.9701}{mi}{106}\RG{42.916}{-37.870}{7.9701}{rm}{61}\RG{49.500}{-37.870}{7.9701}{rm}{49}\RG{101.411}{-37.870}{7.9701}{mi}{106}\RG{105.295}{-37.870}{7.9701}{rm}{61}\RG{111.878}{-37.870}{7.9701}{rm}{49}%
}
\]
命题得证。\allowbreak{}\hskip.1em \RB{8.366}{7.992}{0.503}{%
\RG{0.000}{-0.003}{10.9091}{am}{3}%
}\hskip.1em\allowbreak{}

顺带说明，当所有\allowbreak{}\hskip.1em \RB{30.853}{7.769}{3.759}{%
\RG{0.000}{0.003}{10.9091}{mi}{120}\RG{6.235}{-1.633}{7.9701}{mi}{106}\RG{13.734}{0.003}{10.9091}{mi}{62}\RG{25.335}{0.003}{10.9091}{rm}{48}%
}\hskip.1em\allowbreak{} 时，等号成立当且仅当每一行中所有满足\allowbreak{}\hskip.1em \RB{33.269}{7.769}{3.759}{%
\RG{0.000}{0.003}{10.9091}{mi}{97}\RG{5.767}{-1.633}{7.9701}{mi}{105}\RG{8.644}{-1.633}{7.9701}{mi}{106}\RG{16.149}{0.003}{10.9091}{mi}{62}\RG{27.751}{0.003}{10.9091}{rm}{48}%
}\hskip.1em\allowbreak{} 的位置对应的\allowbreak{}\hskip.1em \RB{10.027}{5.325}{3.759}{%
\RG{0.000}{0.003}{10.9091}{mi}{120}\RG{6.235}{-1.633}{7.9701}{mi}{106}%
}\hskip.1em\allowbreak{}都相等。这正是每一行加权算术–几何平均不等式的取等条件；不必要求全部\allowbreak{}\hskip.1em \RB{10.027}{5.320}{3.764}{%
\RG{0.000}{-0.002}{10.9091}{mi}{120}\RG{6.235}{-1.638}{7.9701}{mi}{106}%
}\hskip.1em\allowbreak{} 相等，例如\allowbreak{}\hskip.1em \RB{28.649}{8.309}{0.502}{%
\RG{0.000}{-0.002}{10.9091}{mi}{65}\RG{11.212}{-0.002}{10.9091}{rm}{61}\RG{22.727}{-0.002}{10.9091}{mi}{73}%
}\hskip.1em\allowbreak{}时任意\allowbreak{}\hskip.1em \RB{6.249}{5.325}{0.617}{%
\RG{0.000}{0.003}{10.9091}{mi}{120}%
}\hskip.1em\allowbreak{} 都取等。

方法小结：行和为\allowbreak{}\hskip.1em \RB{5.071}{7.768}{0.500}{%
\RG{0.000}{0.002}{10.9091}{rm}{49}%
}\hskip.1em\allowbreak{} 提供加权平均，列和为\allowbreak{}\hskip.1em \RB{5.071}{7.768}{0.500}{%
\RG{0.000}{0.002}{10.9091}{rm}{49}%
}\hskip.1em\allowbreak{} 保证相乘后每个\allowbreak{}\hskip.1em \RB{10.027}{5.324}{3.760}{%
\RG{0.000}{0.002}{10.9091}{mi}{120}\RG{6.235}{-1.634}{7.9701}{mi}{106}%
}\hskip.1em\allowbreak{} 的总指数恰为\allowbreak{}\hskip.1em \RB{5.071}{7.768}{0.500}{%
\RG{0.000}{0.002}{10.9091}{rm}{49}%
}\hskip.1em\allowbreak{} ；先单独处理零坐标即可避开对数的定义域问题。

题目来源：谢启鸿《复旦大学高等代数习题集（2023版）》，第123–124页。


\end{document}
